% Les ingérences en politique — Allemand Tle A — Cours signé OnBuch+
% Programme officiel MINESEC. Compiler avec : tectonic AL26-les-ingerences-en-politique.tex
\newcommand{\DOCMATIERE}{Allemand}
\newcommand{\DOCNIVEAU}{Tle A}
\newcommand{\DOCMODULE}{Chapitre 5 : Citoyenneté}
\newcommand{\DOCLECON}{AL26}
\newcommand{\DOCTITRE}{Les ingérences en politique}
% OnBuch+ — préambule « cours signés » ENRICHI (compilé avec tectonic / XeLaTeX)
% Système visuel « Pop » d'OnBuch+ : fond crème (jamais blanc), bordures encre
% épaisses, ombres « dures » décalées, titres Archivo Black en capitales.
% Version MATHÉMATIQUES : graphes (pgfplots/tikz), boîtes pédagogiques
% (définition, propriété, méthode, attention, le savais-tu, exercice résolu,
% exercice, TP, bilan), page de garde et signature OnBuch+ ancrée en bas.
%
% Macros attendues dans main.tex AVANT \input{preamble} :
%   \DOCMATIERE  \DOCNIVEAU  \DOCMODULE  \DOCLECON (numéro)  \DOCTITRE
\documentclass[11pt]{article}

\usepackage{fontspec}
\usepackage[ngerman,french]{babel} % français = langue principale ; l'allemand via \all{..} / textebox
\usepackage[a4paper,margin=2cm,top=3.4cm,bottom=2.8cm,headheight=1.4cm,footskip=1.3cm]{geometry}
\usepackage{amsmath,amssymb,amsfonts}
\usepackage{mathtools} % \xmapsto, \coloneqq... souvent utilisés par les modèles
\usepackage{cancel} % simplifications barrées (bilans de forces)
% \abs{z} (module, valeur absolue) et \norm{u} (norme), utilisés par les modèles
\providecommand{\abs}{}\let\abs\relax\DeclarePairedDelimiter{\abs}{\lvert}{\rvert}
\providecommand{\norm}{}\let\norm\relax\DeclarePairedDelimiter{\norm}{\lVert}{\rVert}
\usepackage[table]{xcolor} % [table] : fournit aussi \rowcolors (alternance de lignes)
\usepackage{pagecolor}
\usepackage{tikz}
\usetikzlibrary{arrows.meta,positioning,calc,shapes.geometric,shapes.misc,shapes.arrows,decorations.pathmorphing,decorations.pathreplacing,decorations.markings,patterns,shadows,mindmap,trees,fit,backgrounds,matrix,intersections,angles,quotes}
\usepackage{pgfplots}
\usepackage[version=4]{mhchem} % équations nucléaires \ce{^{235}_{92}U}
\usepackage[european,siunitx]{circuitikz} % schémas électriques
\pgfplotsset{compat=1.18}
\usepgfplotslibrary{fillbetween,groupplots} % aires entre courbes (calcul intégral)
\usepackage{enumitem}
\usepackage{fancyhdr}
% [most] charge toutes les bibliothèques tcolorbox (listings, minted, raster,
% documentation...) et gonfle inutilement la mémoire TeX sur un document long
% (~300 boîtes) : on ne charge que ce qui est réellement utilisé.
\usepackage[skins,breakable]{tcolorbox}
\usepackage{graphicx}
\usepackage{colortbl}
\usepackage{booktabs}
\usepackage{tabularx}
\usepackage{diagbox} % case d'en-tête diagonale des tableaux à double entrée
% Colonnes extensibles centrée / gauche / droite (C, L, R), souvent utilisées par les modèles
\newcolumntype{C}{>{\centering\arraybackslash}X}
\newcolumntype{L}{>{\raggedright\arraybackslash}X}
\newcolumntype{R}{>{\raggedleft\arraybackslash}X}
\usepackage{array}
\usepackage{multicol}
\usepackage{multirow}
\usepackage{siunitx}
% parse-numbers=false : accepte aussi \SI{2,5 \times 10^{-3}}{...}
\sisetup{locale=FR,output-decimal-marker={,},group-minimum-digits=4,parse-numbers=false}
\DeclareSIUnit\ohm{\ensuremath{\symup{\Omega}}} % U+2126 absent de la police
\DeclareSIUnit\division{div} % réglages d'oscilloscope (ms/div, V/div)
\DeclareSIUnit\tr{tr} % tours (vitesse de rotation, ex. \SI{1500}{\tr\per\minute})
\DeclareSIUnit\molar{M} % concentration molaire (ex. \SI{3}{\molar}, \SI{1}{\micro\molar})
\DeclareSIUnit\an{an} % année (le modèle confond parfois avec \year, un compteur TeX) — ex. \SI{5}{\kilo\meter\per\an}
\DeclareSIUnit\FCFA{FCFA} % franc CFA (ex. \SI{500}{\FCFA\per\kilo\gram})
\DeclareSIUnit\degreeBrix{\textdegree Brix} % teneur en sucre d'un sirop/jus (réfractométrie)
\DeclareSIUnit\month{mois} % le modèle utilise parfois \month, un compteur TeX, comme unité de durée
\usepackage{qrcode}
% \degree : commande fréquemment utilisée par les modèles pour un angle ou une
% température en dehors de \SI{}{} (non fournie par les paquets chargés).
\providecommand{\degree}{^{\circ}}
% \qed (fin de démonstration), utilisé par les modèles sans amsthm
\providecommand{\qed}{\hfill$\square$}
% \barre : double barre des valeurs interdites dans un tableau de variations (tkz-tab)
\providecommand{\barre}{\ensuremath{\|}}
% environnement proof (amsthm non chargé) utilisé par les modèles
\ifdefined\proof\else\newenvironment{proof}[1][Démonstration]{\par\noindent\textit{#1.}\ }{\qed\par}\fi
\usepackage{lastpage}
\usepackage{hyperref}
\hypersetup{hidelinks}

% ============================================================
% Palette « Pop » OnBuch+ — mêmes hex que la classe P (plus_ui.dart)
% ============================================================
\definecolor{poporange}{HTML}{F59321}
\definecolor{popdark}{HTML}{E07A0C}
\definecolor{popcream}{HTML}{FFF6E8}
\definecolor{popink}{HTML}{1C1714}
\definecolor{poppurple}{HTML}{7A5AE0}
\definecolor{popgreen}{HTML}{2E9E6A}
\definecolor{poppink}{HTML}{E0457B}
\definecolor{popblue}{HTML}{2F7DE1}
\definecolor{popgold}{HTML}{C98A12}
\definecolor{popmuted}{HTML}{8A7F76}
\definecolor{popline}{HTML}{EADFD2}
\definecolor{popgray}{HTML}{8A7F76}
\colorlet{popgrey}{popgray}
% Alias tolérants (le modèle nomme parfois les couleurs différemment)
\colorlet{popwhite}{white}
\colorlet{popblack}{popink}
\colorlet{popred}{poppink}
\colorlet{popyellow}{popgold}
% Teintes claires (fonds de tableaux/encadrés) — le modèle nomme parfois
% les couleurs avec un suffixe L (ex. popblueL) sans les avoir définies.
\colorlet{poporangeL}{poporange!20}
\colorlet{popdarkL}{popdark!20}
\colorlet{poppurpleL}{poppurple!20}
\colorlet{popgreenL}{popgreen!20}
\colorlet{poppinkL}{poppink!20}
\colorlet{popblueL}{popblue!20}
\colorlet{popgoldL}{popgold!20}
\colorlet{popredL}{popred!20}
\colorlet{popgrayL}{popgray!20}
% Teintes claires pour fonds de tableaux / zones
\colorlet{popblueL}{popblue!10!white}
\colorlet{poporangeL}{poporange!14!white}
\colorlet{popgreenL}{popgreen!12!white}
\colorlet{poppurpleL}{poppurple!12!white}
\colorlet{poppinkL}{poppink!12!white}
\colorlet{popgoldL}{popgold!14!white}

\pagecolor{popcream}

% ============================================================
% Typographie
% ============================================================
\defaultfontfeatures{Ligatures=TeX}
\setmainfont{PlusJakartaSans-Medium.ttf}[Path=../../../fonts/,BoldFont=PlusJakartaSans-Bold.ttf]
% Maths en sans-serif (Fira Math) avec chiffres et lettres droites de Plus
% Jakarta, pour que les formules se fondent dans le texte.
\usepackage[math-style=ISO,bold-style=ISO]{unicode-math}
\setmathfont{FiraMath-Regular.otf}[Path=../../../fonts/]
\setmathfont{PlusJakartaSans-Medium.ttf}[Path=../../../fonts/,range={up/num,up/{Latin,latin},\mathup}]
\usepackage{icomma} % virgule décimale sans espace en mode math
% Caractères mathématiques Unicode absents des polices de texte (Plus Jakarta,
% Archivo Black) : rendus par leur équivalent mathématique, en texte comme en
% formule — sinon ils s'affichent en carré vide (ex. « Arithmétique dans ℤ »).
\usepackage{newunicodechar}
\newunicodechar{ℝ}{\ensuremath{\mathbb{R}}}
\newunicodechar{ℤ}{\ensuremath{\mathbb{Z}}}
\newunicodechar{ℂ}{\ensuremath{\mathbb{C}}}
\newunicodechar{ℕ}{\ensuremath{\mathbb{N}}}
\newunicodechar{ℚ}{\ensuremath{\mathbb{Q}}}
\newunicodechar{𝔻}{\ensuremath{\mathbb{D}}}
\newunicodechar{α}{\ensuremath{\alpha}}
\newunicodechar{β}{\ensuremath{\beta}}
\newunicodechar{γ}{\ensuremath{\gamma}}
\newunicodechar{δ}{\ensuremath{\delta}}
\newunicodechar{ε}{\ensuremath{\varepsilon}}
\newunicodechar{θ}{\ensuremath{\theta}}
\newunicodechar{λ}{\ensuremath{\lambda}}
\newunicodechar{μ}{\ensuremath{\mu}}
\newunicodechar{σ}{\ensuremath{\sigma}}
\newunicodechar{τ}{\ensuremath{\tau}}
\newunicodechar{φ}{\ensuremath{\varphi}}
\newunicodechar{ψ}{\ensuremath{\psi}}
\newunicodechar{ω}{\ensuremath{\omega}}
\newunicodechar{Σ}{\ensuremath{\Sigma}}
% majuscules produites par \MakeUppercase dans les titres (α-aminés -> Α)
\newunicodechar{Α}{\ensuremath{\alpha}}
\newunicodechar{Β}{\ensuremath{\beta}}
\newunicodechar{Γ}{\ensuremath{\Gamma}}
\newunicodechar{Δ}{\ensuremath{\Delta}}
\newunicodechar{Ε}{\ensuremath{\varepsilon}}
\newunicodechar{Θ}{\ensuremath{\Theta}}
\newunicodechar{Λ}{\ensuremath{\Lambda}}
\newunicodechar{Μ}{\ensuremath{\mu}}
\newunicodechar{Τ}{\ensuremath{\tau}}
\newunicodechar{Φ}{\ensuremath{\Phi}}
\newunicodechar{Ψ}{\ensuremath{\Psi}}
\newunicodechar{Ω}{\ensuremath{\Omega}}
\newunicodechar{↦}{\ensuremath{\mapsto}}
\newunicodechar{∈}{\ensuremath{\in}}
\newunicodechar{∉}{\ensuremath{\notin}}
\newunicodechar{∧}{\ensuremath{\wedge}}
\newunicodechar{⇔}{\ensuremath{\Leftrightarrow}}
\newunicodechar{⟺}{\ensuremath{\Longleftrightarrow}}
\newunicodechar{⇒}{\ensuremath{\Rightarrow}}
\newunicodechar{⟹}{\ensuremath{\Longrightarrow}}
\newunicodechar{∘}{\ensuremath{\circ}}
\newunicodechar{∪}{\ensuremath{\cup}}
\newunicodechar{∩}{\ensuremath{\cap}}
\newunicodechar{≡}{\ensuremath{\equiv}}
\newunicodechar{⊂}{\ensuremath{\subset}}
\newunicodechar{⊥}{\ensuremath{\perp}}
\newunicodechar{∃}{\ensuremath{\exists}}
\newunicodechar{∀}{\ensuremath{\forall}}
\newunicodechar{∛}{\ensuremath{\sqrt[3]{\,}}}
\newunicodechar{∜}{\ensuremath{\sqrt[4]{\,}}}
\newunicodechar{⃗}{\ensuremath{{}^{\to}}}
\newunicodechar{ⁿ}{\ifmmode^{n}\else\textsuperscript{\ensuremath{n}}\fi}
\newunicodechar{ˣ}{\ifmmode^{x}\else\textsuperscript{\ensuremath{x}}\fi}
\newunicodechar{ᵇ}{\ifmmode^{b}\else\textsuperscript{\ensuremath{b}}\fi}
\newunicodechar{ʳ}{\ifmmode^{r}\else\textsuperscript{\ensuremath{r}}\fi}
\newunicodechar{ᵘ}{\ifmmode^{u}\else\textsuperscript{\ensuremath{u}}\fi}
\newunicodechar{ʸ}{\ifmmode^{y}\else\textsuperscript{\ensuremath{y}}\fi}
\newunicodechar{ᵃ}{\ifmmode^{a}\else\textsuperscript{\ensuremath{a}}\fi}
\newunicodechar{ᵉ}{\ifmmode^{e}\else\textsuperscript{\ensuremath{e}}\fi}
\newunicodechar{ᵏ}{\ifmmode^{k}\else\textsuperscript{\ensuremath{k}}\fi}
\newunicodechar{ᵖ}{\ifmmode^{p}\else\textsuperscript{\ensuremath{p}}\fi}
\newunicodechar{ᴺ}{\ifmmode^{N}\else\textsuperscript{\ensuremath{N}}\fi}
\newunicodechar{ᶜ}{\ifmmode^{c}\else\textsuperscript{\ensuremath{c}}\fi}
\newunicodechar{ʰ}{\ifmmode^{h}\else\textsuperscript{\ensuremath{h}}\fi}
\newunicodechar{ᵠ}{\ifmmode^{q}\else\textsuperscript{\ensuremath{q}}\fi}
\newunicodechar{ₙ}{\ifmmode_{n}\else\textsubscript{\ensuremath{n}}\fi}
\newunicodechar{ₐ}{\ifmmode_{a}\else\textsubscript{\ensuremath{a}}\fi}
\newunicodechar{ᵢ}{\ifmmode_{i}\else\textsubscript{\ensuremath{i}}\fi}
\newunicodechar{ᵣ}{\ifmmode_{r}\else\textsubscript{\ensuremath{r}}\fi}
\newunicodechar{ₖ}{\ifmmode_{k}\else\textsubscript{\ensuremath{k}}\fi}
\newunicodechar{ᵦ}{\ifmmode_{\beta}\else\textsubscript{\ensuremath{\beta}}\fi}
\newunicodechar{ₓ}{\ifmmode_{x}\else\textsubscript{\ensuremath{x}}\fi}
\newfontfamily\popdisplay{ArchivoBlack-Regular.ttf}[Path=../../../fonts/]
\newfontfamily\poplogo{SpaceGrotesk-Bold.ttf}[Path=../../../fonts/]
\newfontfamily\popsemi{PlusJakartaSans-SemiBold.ttf}[Path=../../../fonts/]
\newfontfamily\popxbold{PlusJakartaSans-ExtraBold.ttf}[Path=../../../fonts/]
\newfontfamily\popxboldls{PlusJakartaSans-ExtraBold.ttf}[Path=../../../fonts/,LetterSpace=3.0]

\setlist[enumerate]{leftmargin=1.7em,itemsep=5pt,topsep=4pt}
\setlist[itemize]{leftmargin=1.7em,itemsep=4pt,topsep=4pt,label={\color{poporange}\textbullet}}
\setlength{\parindent}{0pt}
\setlength{\parskip}{5pt}
\renewcommand{\arraystretch}{1.25}

% pgfplots : style Pop par défaut
\pgfplotsset{
  every axis/.append style={axis line style={popink,line width=1pt},tick style={popink},
    grid style={popline,line width=0.5pt},label style={font=\small},tick label style={font=\footnotesize}},
  popcurve/.style={poporange,line width=1.8pt,no markers,smooth},
}
% Même style disponible dans un \draw TikZ simple (hors pgfplots)
\tikzset{popcurve/.style={poporange,line width=1.8pt}}
\tikzset{popgrid/.style={popline,very thin}}
\colorlet{popcurve}{poporange} % le nom du style est parfois utilisé comme couleur

% ============================================================
% Pill / squiggle
% ============================================================
\newcommand{\poppill}[3]{%
  \tikz[baseline=-0.55ex]\node[draw=popink,line width=1.2pt,rounded corners=2.6pt,%
    fill=#2,inner xsep=6.5pt,inner ysep=3pt]{%
    \popxboldls\fontsize{7}{7}\selectfont\color{#3}\MakeUppercase{#1}};%
}
\newcommand{\popsquiggle}[1]{%
  \tikz[baseline=-0.4ex]{
    \draw[#1,line width=2.6pt,line cap=round]
      plot [smooth] coordinates {(0,0.09) (0.34,0.01) (0.68,0.21) (1.02,0.01) (1.36,0.10)};
  }%
}
% Mot-clé surligné (marqueur orange sous le texte)
\newcommand{\cle}[1]{{\popxbold\color{popdark}#1}}

% ============================================================
% En-tête / pied
% ============================================================
\pagestyle{fancy}
\fancyhf{}
\renewcommand{\headrulewidth}{0pt}
\renewcommand{\footrulewidth}{0pt}
\lhead{\raisebox{-0.15ex}{\poplogo\fontsize{13}{13}\selectfont\color{popink}OnBuch\color{poporange}+}}
\rhead{\poppill{\DOCMATIERE\ \textbullet\ \DOCNIVEAU\ \textbullet\ Leçon \DOCLECON}{white}{popink}}
\lfoot{\popxbold\fontsize{7.6}{7.6}\selectfont\color{popmuted}\MakeUppercase{Cours signé OnBuch+}\ \textbullet\ {\popsemi\DOCTITRE}}
\rfoot{\tikz[baseline=-0.6ex]\node[rounded corners=8.5pt,fill=popink,minimum height=17pt,minimum width=17pt,inner xsep=5pt,inner sep=0]{\popxbold\fontsize{8}{8}\selectfont\color{popcream}\thepage\,/\,\pageref*{LastPage}};}

% ============================================================
% Titres
% ============================================================
\newcommand{\chaptitre}[2]{%
  \begin{tcolorbox}[enhanced,colback=poporange,colframe=popink,boxrule=2.6pt,arc=11pt,
      drop shadow=popink,left=15pt,right=15pt,top=12pt,bottom=11pt,before skip=3pt,after skip=18pt]
    \poppill{#2}{popink}{popcream}\par\vspace{9pt}
    {\raggedright\hyphenpenalty=10000\exhyphenpenalty=10000\popdisplay\fontsize{22}{24}\selectfont\color{popcream}\MakeUppercase{#1}\par}
  \end{tcolorbox}
}
% Section numérotée : pastille orange avec le numéro + titre Archivo
\newcounter{popsec}
\newcommand{\coursec}[1]{%
  \stepcounter{popsec}\setcounter{popsub}{0}%
  \par\vspace{16pt}\noindent
  \begin{minipage}{\linewidth}
  \tikz[baseline=-0.7ex]\node[draw=popink,line width=1.6pt,fill=poporange,rounded corners=4pt,minimum size=22pt,inner sep=0]{\popdisplay\fontsize{12}{12}\selectfont\color{popcream}\thepopsec};%
  \hspace{8pt}{\popdisplay\fontsize{15}{17}\selectfont\color{popink}\MakeUppercase{#1}}\par
  \vspace{4pt}\noindent{\color{poporange}\rule{36pt}{3.6pt}}
  \end{minipage}\par\nopagebreak\vspace{8pt}
}
\newcounter{popsub}
\newcommand{\courssub}[1]{%
  \stepcounter{popsub}%
  \par\vspace{9pt}\noindent{\popxbold\fontsize{11.5}{13}\selectfont\color{popink}{\color{poporange}\thepopsec.\thepopsub}\hspace{6pt}#1}\par\nopagebreak\vspace{4pt}
}

% ============================================================
% Boîtes pédagogiques — même recette : fond blanc, bordure épaisse colorée,
% ombre dure encre, badge plein. Titre optionnel [#1].
% ============================================================
\tcbset{popbase/.style={breakable,enhanced,colback=white,boxrule=2.3pt,arc=9pt,
  shadow={3pt}{-3pt}{0pt}{popink},left=14pt,right=14pt,top=11pt,bottom=11pt,before skip=11pt,after skip=15pt,
  fontupper=\popsemi\fontsize{10.6}{14.5}\selectfont\color{popink},
  fontlower=\popsemi\fontsize{10.6}{14.5}\selectfont\color{popink}}}
% Coche dessinée (le glyphe ✓ n'existe pas dans les polices du document)
\newcommand{\popcheck}[1]{\tikz[baseline=-0.5ex]\draw[#1,line width=1.6pt,line cap=round,line join=round](-0.13,0.0)--(-0.04,-0.09)--(0.13,0.1);}
\providecommand{\coche}{\popcheck{popgreen}}
\providecommand{\resultat}[1]{\par\smallskip\noindent\fcolorbox{popgreen}{popgreenL}{\parbox{\dimexpr\linewidth-2\fboxsep-2\fboxrule\relax}{\textbf{Résultat.} #1}}\par\smallskip}
\newcommand{\popboxhead}[3]{\poppill{#1}{#2}{white}\ifx\relax#3\relax\else\hspace{6pt}{\popxbold\fontsize{10.6}{12}\selectfont\color{popink}#3}\fi\par\vspace{7pt}}

\newtcolorbox{aretenir}[1][]{popbase,colframe=popgreen,before upper={\popboxhead{À retenir}{popgreen}{#1}}}
% Texte en allemand (document de lecture, dialogue, extrait) : cadre
% d'étude. Un titre optionnel (source, auteur) s'affiche dans l'en-tête.
\newtcolorbox{textebox}[1][]{popbase,colframe=popblue,before upper={\popboxhead{Text}{popblue}{#1}\begin{otherlanguage}{ngerman}},after upper={\end{otherlanguage}}}
% Marques ✓ / ✗ (forme correcte / fautive) souvent utilisées par les modèles
\providecommand{\cross}{\textcolor{poppink}{\ensuremath{\times}}}
\providecommand{\xmark}{\cross}\providecommand{\crossmark}{\cross}\providecommand{\cmark}{\ensuremath{\checkmark}}
% Ligne de réponse / justification à compléter (inventée par les modèles)
\providecommand{\justification}[1]{\par\noindent\textit{Justification~:} \dotfill\par}
% \textipa (package tipa non chargé) : la police ne couvre pas l'API -> simple passage
\providecommand{\textipa}[1]{#1}
% Soulignement (ulem non chargé) : \uline, \ul
\providecommand{\uline}[1]{\underline{#1}}\providecommand{\ul}[1]{\underline{#1}}
% \glossaire{\item ...} inventée par les modèles : liste de glossaire
\providecommand{\glossaire}[1]{\begin{itemize}#1\end{itemize}}
% \alert (beamer) inventée par les modèles : texte mis en évidence
\providecommand{\alert}[1]{\textbf{\textcolor{poppink}{#1}}}
% variantes ASCII d'ordinaux écrites en macro
\providecommand{\iere}{\textsuperscript{re}}\providecommand{\ieme}{\textsuperscript{e}}\providecommand{\ere}{\textsuperscript{re}}\providecommand{\eme}{\textsuperscript{e}}\providecommand{\er}{\textsuperscript{er}}\providecommand{\ie}{\textsuperscript{e}}
% Ordinaux français écrits en macro par les modèles : 1\ière, 3\ième
\newcommand{\ière}{\textsuperscript{re}}\newcommand{\ième}{\textsuperscript{e}}
% \exemple (inventée par les modèles) : étiquette « Exemple : »
\providecommand{\exemple}{\textbf{Exemple~:}\ }
% \square utilisable aussi hors mode math (cases à cocher)
\AtBeginDocument{\let\squareMath\square\renewcommand{\square}{\ensuremath{\squareMath}}}
% Mot ou phrase en allemand à l'intérieur d'un paragraphe français
\newcommand{\all}[1]{\ifmmode\text{\textit{#1}}\else\begingroup\selectlanguage{ngerman}\itshape #1\endgroup\fi}
\let\esp\all % alias toléré
\newtcolorbox{exemplebox}[1][]{popbase,colframe=poporange,before upper={\popboxhead{Exemple}{poporange}{#1}}}
\newtcolorbox{definition}[1][]{popbase,colframe=popblue,before upper={\popboxhead{Définition}{popblue}{#1}}}
\newtcolorbox{propriete}[1][]{popbase,colframe=poppurple,before upper={\popboxhead{Propriété}{poppurple}{#1}}}
\newtcolorbox{methode}[1][]{popbase,colframe=popgold,before upper={\popboxhead{Méthode}{popgold}{#1}}}
\newtcolorbox{attention}[1][]{popbase,colframe=poppink,before upper={\popboxhead{Attention}{poppink}{#1}}}
\newtcolorbox{experience}[1][]{popbase,colframe=popblue,colback=popblueL,before upper={\popboxhead{Expérience}{popblue}{#1}}}
% « Le savais-tu ? » : carte violette pleine (contexte, culture, Cameroun)
\newtcolorbox{savaistu}[1][]{popbase,colback=poppurple,colframe=popink,
  fontupper=\popsemi\fontsize{10.4}{14.2}\selectfont\color{white},
  before upper={\poppill{Le savais-tu\,?}{white}{poppurple}\ifx\relax#1\relax\else\hspace{6pt}{\popxbold\color{white}#1}\fi\par\vspace{7pt}}}
% Exercice résolu : énoncé en haut, \tcblower puis la solution sur fond crème
\newcounter{popexo}\newcounter{popexr}
\newtcolorbox{exoresolu}[1][]{popbase,colframe=poporange,colbacklower=popcream,
  segmentation style={popink,line width=1.2pt,dashed},
  before upper={\stepcounter{popexr}\popboxhead{Exercice résolu \thepopexr}{poporange}{#1}},
  before lower={\poppill{Solution}{popink}{popcream}\par\vspace{6pt}}}
% Exercice (non résolu en place) — la correction est dans la partie Corrigés
\newtcolorbox{exercice}[1][]{popbase,colframe=popink,
  before upper={\stepcounter{popexo}\popboxhead{Exercice \thepopexo}{popink}{#1}}}
% Corrigé
\newtcolorbox{corrige}[1][]{popbase,colframe=popgreen,colback=popgreenL,
  before upper={\popboxhead{Corrigé}{popgreen}{#1}}}
% Objectifs / prérequis / bilan
\newtcolorbox{objectifs}{popbase,colframe=popink,colback=poporangeL,
  before upper={\popboxhead{Objectifs de la leçon}{popink}{}}}
\newtcolorbox{prerequis}{popbase,colframe=popink,colback=popblueL,
  before upper={\popboxhead{Prérequis}{popblue}{}}}
\newtcolorbox{bilan}{popbase,colframe=popink,colback=white,boxrule=2.8pt,
  before upper={\popboxhead{Fiche bilan}{poporange}{L'essentiel à connaître pour l'examen}}}
% Figure encadrée avec légende
\newtcolorbox{popfigure}[1][]{enhanced,breakable=false,colback=white,colframe=popline,boxrule=1.4pt,arc=8pt,
  left=10pt,right=10pt,top=10pt,bottom=8pt,before skip=10pt,after skip=12pt,halign=center,
  lower separated=false,halign lower=center,
  fontlower=\popsemi\fontsize{9}{11.5}\selectfont\color{popmuted},
  #1}
\newcommand{\legende}[1]{\par\vspace{6pt}{\popsemi\fontsize{9}{11.5}\selectfont\color{popmuted}#1}}

% ============================================================
% Signature OnBuch+
% ============================================================
\newcommand{\poplogoinline}[1]{{\poplogo\fontsize{#1}{#1}\selectfont\color{popink}OnBuch\color{poporange}+}}
\newcommand{\signatureonbuch}{%
  \begin{tcolorbox}[enhanced,colback=popink,colframe=popink,boxrule=2.2pt,arc=10pt,
      drop shadow=poporange,left=14pt,right=14pt,top=10pt,bottom=10pt,before skip=0pt,after skip=0pt]
    \tikz[baseline=-0.7ex]\node[circle,fill=popgreen,draw=popcream,line width=1.3pt,minimum size=22pt,inner sep=0]{\popcheck{white}};%
    \hspace{9pt}\begin{minipage}[c]{0.62\linewidth}
      {\popxbold\fontsize{12}{13}\selectfont\color{popcream}Signé\ \textbullet\ {\poplogo OnBuch\color{poporange}+}}\par\vspace{2pt}
      {\raggedright\popsemi\fontsize{8}{10.5}\selectfont\color{popline}\MakeUppercase{Équipe pédagogique OnBuch+}\\\MakeUppercase{Conforme au programme officiel MINESEC}\par}
    \end{minipage}\hfill
    {\poplogo\fontsize{22}{22}\selectfont\color{popcream}OnBuch\color{poporange}+}
  \end{tcolorbox}%
}

% ============================================================
% Page de garde
% #1 titre · #2 matière · #3 niveau · #4 module · #5 n° leçon · #6 durée ·
% #7 description / objectifs (texte ou itemize)
% ============================================================
\newcommand{\pagegarde}[7]{%
  \thispagestyle{empty}
  \newgeometry{a4paper,margin=1.7cm,top=1.6cm,bottom=1.4cm}%
  \begin{tikzpicture}[remember picture,overlay]
    \fill[poporange!18] ([xshift=-3.2cm,yshift=-3.6cm]current page.north east) circle (4.6cm);
    \fill[poppurple!14] ([xshift=2.4cm,yshift=7.6cm]current page.south west) circle (3.8cm);
  \end{tikzpicture}%
  {\poplogo\fontsize{30}{30}\selectfont\color{popink}OnBuch\color{poporange}+}\hfill
  \raisebox{0.6ex}{\poppill{Cours signé \textbullet\ Premium}{popink}{popcream}}\par
  \vspace{1pt}\popsquiggle{poppurple}\par
  \vspace{8pt}
  \begin{tcolorbox}[enhanced,colback=poporange,colframe=popink,boxrule=2.8pt,arc=13pt,
      drop shadow=popink,left=18pt,right=18pt,top=12pt,bottom=13pt,after skip=10pt]
    \poppill{#2 \textbullet\ #3}{popink}{popcream}\hspace{5pt}\poppill{#4}{white}{popink}\par\vspace{8pt}
    {\popdisplay\fontsize{14}{14}\selectfont\color{popink}LEÇON #5}\par\vspace{4pt}
    {\raggedright\hyphenpenalty=10000\exhyphenpenalty=10000\popdisplay\fontsize{25}{27}\selectfont\color{popcream}\MakeUppercase{#1}\par}
    \vspace{8pt}
    {\popxbold\fontsize{9.5}{9.5}\selectfont\color{popink}\MakeUppercase{Durée conseillée : #6}}
  \end{tcolorbox}
  \begin{tcolorbox}[enhanced,colback=white,colframe=popink,boxrule=2.2pt,arc=10pt,
      drop shadow=popink,left=16pt,right=16pt,top=10pt,bottom=10pt,after skip=8pt]
    \poppill{Dans cette leçon}{poppurple}{white}\par\vspace{5pt}
    \popsemi\fontsize{9.3}{12.6}\selectfont\color{popink}#7
  \end{tcolorbox}
  \noindent\begin{minipage}[t]{0.487\linewidth}
    \begin{tcolorbox}[enhanced,colback=popgreen,colframe=popink,boxrule=2.2pt,arc=10pt,
        drop shadow=popink,left=11pt,right=11pt,top=10pt,bottom=10pt,height=3.1cm,valign=center]
      \tcbox[colback=white,colframe=popink,boxrule=1.3pt,arc=5pt,boxsep=0pt,nobeforeafter,
        left=3pt,right=3pt,top=3pt,bottom=3pt,valign=center]{\qrcode[height=1.9cm]{https://play.google.com/store/apps/details?id=com.luvvix.onbuch&hl=fr}}%
      \hspace{8pt}\begin{minipage}[c]{0.5\linewidth}
        {\popdisplay\fontsize{11}{12.5}\selectfont\color{white}\MakeUppercase{Télécharge OnBuch}}\par\vspace{4pt}
        {\raggedright\popsemi\fontsize{8.4}{11}\selectfont\color{white}Gratuit \textbullet\ Google Play\\Tuteur IA, annales, concours, résultats\par}
      \end{minipage}
    \end{tcolorbox}
  \end{minipage}\hfill
  \begin{minipage}[t]{0.487\linewidth}
    \begin{tcolorbox}[enhanced,colback=poppurple,colframe=popink,boxrule=2.2pt,arc=10pt,
        drop shadow=popink,left=11pt,right=11pt,top=10pt,bottom=10pt,height=3.1cm,valign=center]
      \tikz[baseline=-0.7ex]\node[circle,fill=white,draw=popink,line width=1.3pt,minimum size=1.6cm,inner sep=0]{\popdisplay\fontsize{24}{24}\selectfont\color{poppurple}+};%
      \hspace{8pt}\begin{minipage}[c]{0.6\linewidth}
        {\popdisplay\fontsize{11}{12.5}\selectfont\color{white}\MakeUppercase{Passe à OnBuch+}}\par\vspace{4pt}
        {\raggedright\popsemi\fontsize{8.4}{11}\selectfont\color{white}Tous les cours signés, Tuteur IA illimité, fiches PDF — onglet OnBuch+ de l'app\par}
      \end{minipage}
    \end{tcolorbox}
  \end{minipage}
  \vspace{8pt}
  \popcontenu
  \vfill
  \signatureonbuch
  \restoregeometry
  \newpage
}

% Rangée « Ce cours contient » (page de garde)
\newcommand{\poptile}[3]{% #1 couleur #2 gros texte #3 légende
  \begin{tcolorbox}[enhanced,colback=#1,colframe=popink,boxrule=2pt,arc=9pt,shadow={2.5pt}{-2.5pt}{0pt}{popink},
      left=6pt,right=6pt,top=9pt,bottom=9pt,halign=center,valign=center,height=1.75cm,width=\linewidth,nobeforeafter]
    {\popdisplay\fontsize{12.5}{13}\selectfont\color{white}\MakeUppercase{#2}}\par\vspace{3pt}
    {\centering\popsemi\fontsize{7.8}{9.5}\selectfont\color{white}#3\par}
  \end{tcolorbox}}
\newcommand{\popcontenu}{%
  \par\noindent{\popxboldls\fontsize{7.5}{7.5}\selectfont\color{popmuted}CE COURS SIGNÉ CONTIENT}\par\vspace{7pt}
  \noindent\begin{minipage}[t]{0.235\linewidth}\poptile{popblue}{Cours}{complet, illustré, pas à pas}\end{minipage}\hfill
  \begin{minipage}[t]{0.235\linewidth}\poptile{poporange}{Méthodes}{et exercices résolus}\end{minipage}\hfill
  \begin{minipage}[t]{0.235\linewidth}\poptile{poppink}{Exercices}{type examen + corrigés}\end{minipage}\hfill
  \begin{minipage}[t]{0.235\linewidth}\poptile{popgreen}{Fiche bilan}{l'essentiel à retenir}\end{minipage}\par}

% Sommaire
\newtcolorbox{sommaire}{enhanced,colback=white,colframe=popink,boxrule=2.2pt,arc=10pt,
  drop shadow=popink,left=18pt,right=18pt,top=13pt,bottom=13pt,before skip=6pt,after skip=20pt,
  before upper={\poppill{Sommaire}{poppurple}{white}\par\vspace{9pt}\popsemi\fontsize{10.6}{14.5}\selectfont\color{popink}}}

% Signature de fin de cours — poussée tout en bas de la dernière page
\newcommand{\findecours}{%
  \par\vfill\nopagebreak
  \begin{center}\popsquiggle{poporange}\end{center}\vspace{4pt}
  \signatureonbuch
}
\AtBeginDocument{\def\llbracket{\mathopen{[\mkern-3.5mu[}}\def\rrbracket{\mathclose{]\mkern-3.5mu]}}} % intervalles d'entiers (glyphe ⟦ absent de la police)
% Glyphes absents de Fira Math : redéfinis à partir de glyphes disponibles
\AtBeginDocument{%
  \def\setminus{\mathbin{\backslash}}\let\smallsetminus\setminus
  \def\vdots{\mathord{\vbox{\baselineskip3.2pt\lineskiplimit0pt\kern4pt\hbox{.}\hbox{.}\hbox{.}}}}%
  \def\ddots{\mathinner{\mkern1mu\raise6.5pt\hbox{.}\mkern2mu\raise3.6pt\hbox{.}\mkern2mu\raise0.7pt\hbox{.}\mkern1mu}}%
  \def\triangle{\mathord{\scalebox{0.9}{$\symup{\Delta}$}}}%
  \def\nabla{\mathord{\raisebox{\depth}{\scalebox{1}[-1]{$\symup{\Delta}$}}}}%
  \def\checkmark{\popcheck{popdark}}%
  \def\star{\mathbin{\ast}}\def\bigstar{\mathord{\ast}}%
  \def\diamond{\mathbin{\scalebox{0.6}{$\Diamond$}}}%
  \def\bigcup{\mathop{\mathchoice{\raisebox{-0.3em}{\scalebox{1.7}{$\cup$}}}{\scalebox{1.25}{$\cup$}}{\cup}{\cup}}\displaylimits}%
  \def\bigcap{\mathop{\mathchoice{\raisebox{-0.3em}{\scalebox{1.7}{$\cap$}}}{\scalebox{1.25}{$\cap$}}{\cap}{\cap}}\displaylimits}%
  \def\lesseqqgtr{\mathrel{\lessgtr}}\def\bigoplus{\mathop{\oplus}\displaylimits}%
}
\providecommand{\ssi}{si et seulement si} % abréviation fréquente
\AtBeginDocument{\providecommand{\wideparen}{\overparen}} % arc de cercle

\begin{document}

\pagegarde{Les ingérences en politique}{Allemand}{Tle A}{Chapitre 5 : Citoyenneté}{AL26}{2 h}{Leçon de type « texte » : lecture et commentaire d'un texte allemand sur les ingérences en politique, enrichissement du lexique thématique (Einmischung, Macht, Souveränität...) et révision du passif, du plus-que-parfait et du subjonctif II. Production guidée : résumé et court commentaire argumenté.
\vspace{4pt}
\begin{multicols}{2}\small
\begin{itemize}
\item À la fin de cette leçon, je sais comprendre globalement puis en détail un texte allemand sur les ingérences en politique.
\item Je sais employer correctement le vocabulaire du thème (die Einmischung, die Macht, das Land, die Souveränität, der Einfluss, der Staat) avec articles et pluriels.
\item Je sais former et utiliser le passif (Präsens, Präteritum, Perfekt) dans des phrases.
\item Je sais conjuguer et employer le plus-que-parfait pour exprimer l'antériorité.
\item Je sais utiliser le subjonctif II pour exprimer l'hypothèse et le souhait.
\item Je sais répondre à des questions de compréhension en phrases complètes.
\end{itemize}\end{multicols}}

\begin{sommaire}
\begin{multicols}{2}
\begin{enumerate}
\item Texte support et première compréhension
\item Lexique thématique par champs
\item Grammaire 1 : le passif
\item Grammaire 2 : plus-que-parfait et subjonctif II
\item Méthode du commentaire et commentaire modèle
\item Production guidée et entraînement au Bac
\item Activité d'intégration
\item Méthodes et exercices résolus
\item Exercices
\item Corrigés des exercices
\item Fiche bilan
\end{enumerate}
\end{multicols}
\end{sommaire}

\begin{objectifs}
\begin{itemize}
\item À la fin de cette leçon, je sais comprendre globalement puis en détail un texte allemand sur les ingérences en politique.
\item Je sais employer correctement le vocabulaire du thème (die Einmischung, die Macht, das Land, die Souveränität, der Einfluss, der Staat) avec articles et pluriels.
\item Je sais former et utiliser le passif (Präsens, Präteritum, Perfekt) dans des phrases.
\item Je sais conjuguer et employer le plus-que-parfait pour exprimer l'antériorité.
\item Je sais utiliser le subjonctif II pour exprimer l'hypothèse et le souhait.
\item Je sais répondre à des questions de compréhension en phrases complètes.
\item Je sais rédiger un résumé et un court commentaire structuré d'un texte.
\item Je sais traduire des phrases du français vers l'allemand en appliquant les règles étudiées.
\end{itemize}
\end{objectifs}

\begin{prerequis}
\begin{itemize}
\item Le présent et le Präteritum des verbes forts et faibles (Seconde/Première)
\item Le Perfekt avec haben/sein (Première)
\item Les articles définis et indéfinis et les déclinaisons de base
\item Le vocabulaire général de la politique vu en Première (die Regierung, wählen, die Demokratie)
\item La construction de la phrase allemande : place du verbe (V2, verbe final en subordonnée)
\end{itemize}
\end{prerequis}

\newpage

\coursec{Texte support et première compréhension}

\noindent
\textbf{Situation d'accroche.} En 2019, l'Allemagne a critiqué publiquement certaines décisions politiques de pays africains, et aussitôt des voix se sont élevées : „Einmischung!“ (ingérence !). Au Cameroun aussi, on débat souvent du rôle des puissances étrangères dans nos affaires intérieures — pensons aux déclarations de l'Union européenne ou de la France sur le processus électoral. Quand un État intervient-il légitimement pour protéger des droits fondamentaux ? Quand viole-t-il la \cle{Souveränität} (souveraineté) d'un autre ? Ce débat passionne aussi l'Autriche et la Suisse, pays neutres par tradition. Aujourd'hui, nous lisons un texte allemand actuel sur cette question brûlante, afin de construire le vocabulaire, d'observer la grammaire en contexte et de préparer l'expression argumentée attendue au Baccalauréat.

\courssub{Lecture silencieuse et découverte du texte}

\noindent
Lisez le texte ci-dessous une première fois \emph{sans dictionnaire}. Repérez les mots transparents (internationalismes), la structure par paragraphes et les connecteurs logiques (\all{Immer wieder}, \all{Manche ... Andere}, \all{Nachdem}, \all{Doch}). Ensuite, lisez-le à voix haute en respectant la prononciation des Umlaute (ä, ö, ü) et du \all{ß} (Eszett, prononcé « ss »).

\begin{textebox}[Einmischung in die Politik?]
Immer wieder mischen sich große Staaten in die inneren Angelegenheiten anderer Länder ein. Manche Politiker und Juristen sagen: Das ist eine klare Verletzung der Souveränität! Jeder Staat hat das Recht, seine eigene Politik ohne äußeren Druck zu bestimmen. Andere hingegen meinen: Wenn Menschenrechte schwer verletzt werden, darf die Weltgemeinschaft nicht tatenlos zusehen. Die Verantwortung zu schützen — „Responsibility to Protect“ — gebietet manchmal ein Eingreifen.

Ein konkretes Beispiel: Nachdem ein Land jahrelang unter massivem internationalen Druck gestanden hatte, wurden endlich freie und faire Wahlen organisiert. Beobachter sind sich einig: Wäre die Weltgemeinschaft früher und entschlossener eingeschritten, hätte man viel Leid und Instabilität verhindern können. Doch wo liegt die Grenze zwischen legitimer Hilfe und unzulässiger Einmischung? Diese Frage bleibt schwer zu beantworten, denn jedes Land beansprucht seine politische Unabhängigkeit.
\end{textebox}

\noindent
\textbf{Glossaire de base} (mots du texte à mémoriser avec article, pluriel et traduction) :

\begin{center}
\begin{tabularx}{\linewidth}{|>{\bfseries}l|X|X|}\hline
\rowcolor{popblueL}
\textbf{Allemand} & \textbf{Français} & \textbf{Remarque / Famille} \\ \hline
die Einmischung, -en & l'ingérence & de \all{sich einmischen} (s'ingérer) \\ \hline
die Souveränität, -en & la souveraineté & \all{souverän} (adj.) = souverain \\ \hline
der Einfluss, -¨sse & l'influence & \all{einflussreich} = influent \\ \hline
sich einmischen in (+ Akk.) & s'ingérer dans & \textbf{Préposition \cle{in} + accusatif} \\ \hline
verletzen & violer, blesser & \all{die Verletzung, -en} (la violation) \\ \hline
der Druck, -¨e & la pression & \all{unter Druck stehen} = subir des pressions \\ \hline
die Menschenrechte (Pl.) & les droits de l'homme & \all{das Menschenrecht, -e} (sing.) \\ \hline
die Weltgemeinschaft, -en & la communauté internationale & \all{die Welt} + \all{die Gemeinschaft} \\ \hline
das Leid & la souffrance, le malheur & pas de pluriel usuel \\ \hline
die Grenze, -n & la frontière, la limite & \all{grenzenlos} = illimité \\ \hline
die Unabhängigkeit & l'indépendance & \all{unabhängig} = indépendant \\ \hline
die Angelegenheit, -en & l'affaire & \all{innere Angelegenheiten} = affaires intérieures \\ \hline
einschreiten (trat ein, ist eingeschritten) & intervenir & verbe à particule séparable \\ \hline
verhindern & empêcher, prévenir & \all{die Verhinderung} (l'empêchement) \\ \hline
beanspruchen & revendiquer & \all{der Anspruch} (la revendication) \\ \hline
\end{tabularx}
\end{center}

\begin{definition}[die Einmischung]
L'\textbf{Einmischung} (féminin, pluriel \all{die Einmischungen}) désigne l'intervention — souvent non désirée — d'un État, d'une organisation internationale ou d'un acteur puissant dans les affaires intérieures (\all{innere Angelegenheiten}) d'un État souverain. En droit international, on distingue l'ingérence interdite (\all{verbotene Einmischung}) de l'intervention humanitaire légitime (\all{humanitäre Intervention}).
\end{definition}

\courssub{Compréhension globale : stratégie et repérages}

\begin{methode}[Comment aborder la compréhension globale ?]
\begin{enumerate}
\item \textbf{Lecture rapide} (2--3 min) : ne cherchez pas chaque mot ; visez le \emph{sujet}, le \emph{ton} et la \emph{thèse}.
\item \textbf{Soulignez les mots-clés} du champ lexical politique : \all{Staat, Souveränität, Menschenrechte, Weltgemeinschaft, Druck, Wahlen, Unabhängigkeit}.
\item \textbf{Identifiez la structure} : paragraphe 1 = thèse / antithèse ; paragraphe 2 = exemple + argument contrefactuel + question ouverte.
\item \textbf{Formulez la thèse de l'auteur en une phrase} en français : « L'auteur présente le dilemme entre respect de la souveraineté et devoir de protéger les droits de l'homme, sans trancher définitivement. »
\end{enumerate}
\end{methode}

\noindent
\textbf{Questions de repérage global} (répondez en allemand, phrases courtes) :

\begin{enumerate}
\item \all{Worum geht es im Text?} (De quoi traite le texte ?)
\item \all{Welche zwei Positionen werden gegenübergestellt?} (Quelles deux positions sont opposées ?)
\item \all{Was ist die These des Autors?} (Quelle est la thèse / la conclusion de l'auteur ?)
\end{enumerate}

\begin{exoresolu}[Réponses guidées à la compréhension globale]
\textbf{Énoncé :} Répondez aux trois questions ci-dessus en vous appuyant sur le texte.

\tcblower
\textbf{Solution détaillée :}
\begin{enumerate}
\item \all{Es geht um die Frage, ob die Einmischung in die Politik anderer Länder erlaubt ist oder die Souveränität verletzt.} (Il s'agit de la question de savoir si l'ingérence dans la politique d'autres pays est permise ou si elle viole la souveraineté.)
\item \all{Position 1: Die Einmischung verletzt die Souveränität. Position 2: Bei Menschenrechtsverletzungen muss die Weltgemeinschaft eingreifen.} (Position 1 : l'ingérence viole la souveraineté. Position 2 : en cas de violation des droits de l'homme, la communauté internationale doit intervenir.)
\item \all{Der Autor zeigt das Dilemma auf: Es gibt keine einfache Grenze zwischen Hilfe und Einmischung. Er ergreift nicht eindeutig Partei.} (L'auteur met en lumière le dilemme : il n'y a pas de frontière simple entre aide et ingérence. Il ne prend pas parti clairement.)
\end{enumerate}
\end{exoresolu}

\courssub{Compréhension détaillée : arguments et structure}

\noindent
Relisez le texte paragraphe par paragraphe. Complétez le tableau suivant en allemand (mots du texte ou phrases courtes).

\begin{center}
\begin{tabularx}{\linewidth}{|>{\bfseries}X|X|X|}\hline
\rowcolor{popblueL}
\textbf{Paragraphe} & \textbf{Idée principale (Stichworte)} & \textbf{Arguments / Exemples (Argumente / Beispiele)} \\ \hline
1 & \all{Pro und Kontra der Einmischung; Souveränität vs. Menschenrechte} & \all{Arg. 1: Verletzung der Souveränität. Arg. 2: Verantwortung zu schützen (R2P).} \\ \hline
2 & \all{Beispiel: Druck $\rightarrow$ Wahlen; Konjunktiv II: verpasstes Eingreifen; offene Frage} & \all{Nach dem Druck: freie Wahlen. Konjunktiv II: früheres Eingreifen $\rightarrow$ Leid verhindert. Grenze unklar.} \\ \hline
\end{tabularx}
\end{center}

\noindent
\textbf{Questions de compréhension détaillée} (réponses attendues en allemand) :

\begin{enumerate}
\item \all{Was bedeutet „Verletzung der Souveränität“ im Kontext des Textes?}
\item \all{Welches konkrete Beispiel nennt der Autor für erfolgreichen Druck?}
\item \all{Warum benutzt der Autor den Konjunktiv II im Satz „Wäre die Weltgemeinschaft früher eingeschritten, hätte man viel Leid verhindern können“?}
\item \all{Wie lautet die Schlussfrage des Textes und warum ist sie „schwer zu beantworten“?}
\end{enumerate}

\begin{exoresolu}[Réponses détaillées]
\textbf{Énoncé :} Répondez aux quatre questions en allemand, en phrases complètes.

\tcblower
\textbf{Solution :}
\begin{enumerate}
\item \all{Es bedeutet, dass ein fremder Staat in die inneren Angelegenheiten eines souveränen Staates eingreift, ohne dazu berechtigt zu sein.} (Cela signifie qu'un État étranger intervient dans les affaires intérieures d'un État souverain sans y être autorisé.)
\item \all{Ein Land stand jahrelang unter internationalem Druck; daraufhin wurden freie Wahlen organisiert.} (Un pays a subi des pressions internationales pendant des années ; ensuite, des élections libres ont été organisées.)
\item \all{Der Konjunktiv II drückt eine irreale Bedingung in der Vergangenheit aus: Die Weltgemeinschaft ist NICHT früher eingeschritten, deshalb konnte man das Leid NICHT verhindern. Es ist eine verpasste Möglichkeit.} (Le subjonctif II exprime une condition irréelle au passé : la communauté internationale N'EST PAS intervenue plus tôt, donc on n'a PAS pu empêcher la souffrance. C'est une occasion manquée.)
\item \all{Die Schlussfrage lautet: „Wo liegt die Grenze zwischen Hilfe und Einmischung?“ Sie ist schwer zu beantworten, weil jedes Land seine Unabhängigkeit beansprucht und „Hilfe“ oft subjektiv definiert wird.} (La question finale est : « Où se situe la frontière entre aide et ingérence ? » Elle est difficile car chaque pays revendique son indépendance et l'« aide » est souvent définie de manière subjective.)
\end{enumerate}
\end{exoresolu}

\courssub{Reformulation et structuration du texte}

\noindent
Exercez-vous à résumer et à structurer — compétence clé pour l'épreuve de \emph{résumé} et de \emph{commentaire} au Bac.

\begin{exemplebox}[Titre possible et idées principales par paragraphe]
\textbf{Titre proposé :} \all{Zwischen Souveränität und Schutzverantwortung — das Dilemma der Einmischung} (Entre souveraineté et responsabilité de protéger — le dilemme de l'ingérence).

\textbf{Idée principale §1 :} L'auteur expose le conflit classique : l'ingérence viole le droit à l'autodétermination (\all{Souveränität}), mais l'inaction face aux violations graves des droits de l'homme (\all{Menschenrechtsverletzungen}) est moralement inacceptable (\all{Responsibility to Protect}).

\textbf{Idée principale §2 :} Un exemple (pression $\rightarrow$ élections libres) et une hypothèse contrefactuelle (subjonctif II) montrent qu'une intervention précoce peut éviter le malheur. La conclusion ouvre le débat : la frontière reste floue.
\end{exemplebox}

\courssub{Observation linguistique : formes grammaticales repérées dans le texte}

\noindent
Avant d'étudier les règles en profondeur (sections 3 et 4 de la leçon), identifions trois structures clés présentes dans le texte.

\begin{center}
\begin{tabularx}{\linewidth}{|>{\bfseries}l|X|X|}\hline
\rowcolor{popblueL}
\textbf{Structure} & \textbf{Exemple du texte} & \textbf{Traduction / Fonction} \\ \hline
\cle{Passiv (Vorgangspassiv)} & \all{wurden ... Wahlen organisiert} & « des élections furent organisées » — Präteritum passif ; l'accent est mis sur l'action, pas sur l'agent \\ \hline
\cle{Plusquamperfekt} & \all{hatte ... gestanden} & « avait subi » — antériorité par rapport à un passé (\all{nachdem} + Plusquamperfekt) \\ \hline
\cle{Konjunktiv II (Vergangenheit)} & \all{Wäre ... eingeschritten, hätte ... verhindern können} & « si elle était intervenue, on aurait pu empêcher » — irréel du passé \\ \hline
\end{tabularx}
\end{center}

\begin{attention}[Piège fréquent : \all{sich einmischen} vs. \all{sich mischen}]
Ne confondez pas :
\begin{itemize}
\item \all{sich \textbf{ein}mischen in + Akkusativ} = \textbf{s'ingérer dans}, intervenir de manière indue. Ex. : \all{Er mischt sich in meine Angelegenheiten ein.} (Il s'ingère dans mes affaires.)
\item \all{sich mischen} (sans \textbf{ein}) = \textbf{se mélanger}. Ex. : \all{Öl mischt sich nicht mit Wasser.} (L'huile ne se mélange pas à l'eau.)
\end{itemize}
\textbf{Erreur type du francophone :} *\all{sich einmischen an} ou *\all{sich einmischen bei}. \textbf{Correct :} \all{in + Akkusativ} (\all{in die Politik}, \all{in die inneren Angelegenheiten}). Attention aussi à la particule séparable : \all{ein} va en fin de phrase au présent et au prétérit.
\end{attention}

\begin{exemplebox}[Exemples traduits : structures observées]
\begin{enumerate}
\item \all{Nachdem ein Land jahrelang unter Druck \textbf{gestanden hatte}, wurden freie Wahlen \textbf{organisiert}.}
« Après qu'un pays \textbf{eut subi} des pressions pendant des années, des élections libres \textbf{furent organisées}. » (Plus-que-parfait après \all{nachdem} + passif au prétérit)
\item \all{\textbf{Wäre} die Weltgemeinschaft früher \textbf{eingeschritten}, \textbf{hätte} man viel Leid \textbf{verhindern können}.}
« \textbf{Si} la communauté internationale \textbf{était intervenue} plus tôt, \textbf{on aurait pu} empêcher beaucoup de souffrance. » (Konjunktiv II passé : \all{wäre} + Partizip II ; \all{hätte} + participe + infinitif du modal)
\item \all{Die Souveränität \textbf{wird} oft \textbf{verletzt}.}
« La souveraineté \textbf{est} souvent \textbf{violée}. » (Passif présent : \all{werden} + Partizip II)
\end{enumerate}
\end{exemplebox}

\courssub{Carte mentale lexicale : le champ de l'ingérence}

\noindent
La figure ci-dessous visualise les liens sémantiques autour du mot-clé \all{Einmischung}. Utilisez-la pour réviser le vocabulaire et préparer une prise de parole.

\begin{popfigure}
\begin{tikzpicture}[
  mindmap,
  concept color=popblue,
  level 1/.append style={level distance=3.5cm, sibling angle=90},
  every node/.style={font=\small, align=center},
  edge from parent/.style={draw=popline, thick}
]
\node[concept, text=white, scale=1.1] {Einmischung}
  child[concept color=poporange] { node[concept, text=white, align=center]{Souveränität\\ Verletzung} }
  child[concept color=popgreen] { node[concept, text=white, align=center]{Einfluss\\ Druck} }
  child[concept color=poppurple] { node[concept, text=white, align=center]{Menschenrechte\\ Schutzpflicht (R2P)} }
  child[concept color=poppink] { node[concept, text=white, align=center]{Staat\\ Unabhängigkeit} };
\end{tikzpicture}
\legende{Carte mentale : champs lexicaux liés à \all{Einmischung} (ingérence).}
\end{popfigure}

\courssub{Entraînement autonome : exercice de transfert}

\begin{exercice}[Vocabulaire et structures du texte]
Complétez les phrases suivantes en allemand en utilisant le vocabulaire du glossaire et les formes grammaticales observées (passif, plus-que-parfait, Konjunktiv II). Traduisez ensuite en français.
\begin{enumerate}
\item Nachdem die Regierung jahrelang unter \_\_\_\_\_\_\_\_\_\_ (Druck) \_\_\_\_\_\_\_\_\_\_ (stehen, Plusquamperfekt), \_\_\_\_\_\_\_\_\_\_ (werden, Präteritum) politische Gefangene \_\_\_\_\_\_\_\_\_\_ (freilassen, Partizip II).
\item \_\_\_\_\_\_\_\_\_\_ (sein, Konjunktiv II) die UNO früher \_\_\_\_\_\_\_\_\_\_ (einschreiten, Partizip II), \_\_\_\_\_\_\_\_\_\_ (haben, Konjunktiv II) man den Krieg \_\_\_\_\_\_\_\_\_\_ (verhindern können).
\item Eine \_\_\_\_\_\_\_\_\_\_ (Einmischung) in die inneren \_\_\_\_\_\_\_\_\_\_ (Angelegenheiten) eines anderen Staates \_\_\_\_\_\_\_\_\_\_ (verletzen, Präsens) oft die \_\_\_\_\_\_\_\_\_\_ (Souveränität).
\end{enumerate}
\end{exercice}

\begin{corrige}[Exercice 1]
\begin{enumerate}
\item \all{Nachdem die Regierung jahrelang unter \textbf{Druck gestanden hatte}, \textbf{wurden} politische Gefangene \textbf{freigelassen}.} — « Après que le gouvernement \textbf{eut subi} des pressions pendant des années, des prisonniers politiques \textbf{furent libérés}. » (Verbe fort : \all{stehen -- stand -- hat gestanden} ; \all{freilassen} est séparable : \all{freigelassen}.)
\item \all{\textbf{Wäre} die UNO früher \textbf{eingeschritten}, \textbf{hätte} man den Krieg \textbf{verhindern können}.} — « \textbf{Si} l'ONU \textbf{était intervenue} plus tôt, \textbf{on aurait pu} empêcher la guerre. » (Verbe fort : \all{einschreiten -- schritt ein -- ist eingeschritten} ; auxiliaire \all{sein} car mouvement/changement.)
\item \all{Eine \textbf{Einmischung} in die inneren \textbf{Angelegenheiten} eines anderen Staates \textbf{verletzt} oft die \textbf{Souveränität}.} — « Une \textbf{ingérence} dans les \textbf{affaires intérieures} d'un autre État \textbf{viole} souvent la \textbf{souveraineté}. »
\end{enumerate}
\end{corrige}

\begin{savaistu}[Le Cameroun et le principe de non-ingérence]
La Charte de l'Union africaine (article 4) consacre le respect des frontières existantes et la non-ingérence. Pourtant, son paragraphe 4(h) prévoit le droit d'intervention de l'Union dans un État membre « en cas de crimes de guerre, de génocide et de crimes contre l'humanité ». Ce dilemme juridique reflète exactement le débat de notre texte : \all{Souveränität} contre \all{Schutzverantwortung} (responsabilité de protéger). Le Cameroun a parfois rejeté certaines résolutions d'instances étrangères comme une \all{Einmischung in die inneren Angelegenheiten}, tout en acceptant l'appui technique de partenaires pour l'organisation d'élections. La nuance entre \all{Hilfe} (aide) et \all{Einmischung} (ingérence) se joue souvent sur le terrain du \emph{consentement} de l'État concerné.
\end{savaistu}

\coursec{Lexique thématique par champs}

Dans la section précédente, nous avons lu et compris globalement le texte support sur les ingérences en politique. Vous avez rencontré des mots comme \all{die Einmischung}, \all{die Souveränität} ou \all{der Einfluss}. Pour commenter ce texte, le traduire et en débattre à l'oral, il faut maintenant organiser ce vocabulaire de manière systématique. Nous allons le classer en quatre champs lexicaux, puis construire des familles de mots et mémoriser des expressions figées. Retenez le principe d'or de l'apprentissage du vocabulaire allemand : on n'apprend jamais un nom seul, mais toujours avec son \cle{article} et son \cle{pluriel}.

\courssub{Champ 1 — L'ingérence et l'influence}

\begin{definition}[Le champ lexical de l'ingérence]
Une \cle{ingérence} (\all{die Einmischung}) est l'action d'intervenir dans les affaires d'un autre État ou d'une autre personne sans y être invité. Le champ lexical regroupe aussi l'idée d'\cle{influence} (\all{der Einfluss}), de \cle{pression} (\all{der Druck}) et de \cle{sanction} (\all{die Sanktion}).
\end{definition}

\begin{center}
\begin{tabularx}{\linewidth}{|l|l|X|}
\hline
\rowcolor{popblueL} Mot allemand (article + pluriel) & Traduction & Exemple \\
\hline
\all{die Einmischung, die Einmischungen} & l'ingérence, l'immixtion & \all{Die Einmischung in die Wahlen ist verboten.} \\
\hline
\all{der Einfluss, die Einflüsse} & l'influence & \all{Der Einfluss der Medien wächst.} \\
\hline
\all{die Intervention, die Interventionen} & l'intervention & \all{Die Intervention der Armee war umstritten.} \\
\hline
\all{der Druck (nur Sing.)} & la pression & \all{Der internationale Druck nimmt zu.} \\
\hline
\all{die Sanktion, die Sanktionen} & la sanction & \all{Die Sanktionen treffen die Wirtschaft.} \\
\hline
\all{sich einmischen (in + Akk)} & s'immiscer (dans) & \all{Man soll sich nicht in die Wahlen einmischen.} \\
\hline
\all{beeinflussen} & influencer & \all{Große Staaten beeinflussen kleine Länder.} \\
\hline
\all{intervenieren} & intervenir & \all{Die UNO will nicht militärisch intervenieren.} \\
\hline
\end{tabularx}
\end{center}

\begin{attention}[Verbes à particule et verbes en -ieren]
\all{sich einmischen} est un verbe \cle{réfléchi} à \cle{particule séparable} : \all{Er mischt sich in alles ein.} « Il se mêle de tout. » Au Perfekt : \all{Er hat sich eingemischt.} En revanche, \all{beeinflussen} et \all{intervenieren} sont inséparables et forment leur Perfekt \textbf{sans} \all{ge-} : \all{Sie hat die Wahl beeinflusst.} / \all{Er hat interveniert.} Ne dites jamais \emph{gebeeinflusst} !
\end{attention}

\courssub{Champ 2 — L'État et le pouvoir}

\begin{center}
\begin{tabularx}{\linewidth}{|l|l|X|}
\hline
\rowcolor{popblueL} Mot allemand (article + pluriel) & Traduction & Exemple \\
\hline
\all{der Staat, die Staaten} & l'État & \all{Der Staat schützt seine Bürger.} \\
\hline
\all{das Land, die Länder} & le pays & \all{Kamerun ist ein Land in Zentralafrika.} \\
\hline
\all{die Macht, die Mächte} & le pouvoir & \all{Die Macht gehört dem Volk.} \\
\hline
\all{die Regierung, die Regierungen} & le gouvernement & \all{Die Regierung lehnt jede Einmischung ab.} \\
\hline
\all{die Souveränität (nur Sing.)} & la souveraineté & \all{Die Souveränität des Landes muss respektiert werden.} \\
\hline
\all{die Unabhängigkeit (nur Sing.)} & l'indépendance & \all{Kamerun feiert jedes Jahr seine Unabhängigkeit.} \\
\hline
\all{regieren} & gouverner & \all{Der Präsident regiert das Land.} \\
\hline
\all{herrschen (über + Akk)} & régner (sur), dominer & \all{Früher herrschten Könige über viele Länder.} \\
\hline
\end{tabularx}
\end{center}

\begin{definition}[die Souveränität]
\all{die Souveränität} (la souveraineté) désigne le pouvoir suprême d'un État de décider librement de ses affaires intérieures et extérieures, sans intervention extérieure. C'est le concept central du débat sur l'ingérence : celui qui s'immisce dans les affaires d'un État viole sa souveraineté. Adjectif correspondant : \all{souverän} — \all{Kamerun ist ein souveräner Staat.} « Le Cameroun est un État souverain. »
\end{definition}

\begin{attention}[Genre et pluriel : deux pièges classiques]
\all{das Land} (pluriel \all{die Länder}) signifie « le pays » : ne confondez pas avec le français « la terre ». Et \all{die Macht} est \textbf{féminin} (pluriel \all{die Mächte} = les puissances), alors que « le pouvoir » est masculin en français. Erreur fréquente : \emph{der Macht}. Apprenez toujours article + pluriel : \all{der Staat, die Staaten} ; \all{das Land, die Länder} ; \all{die Macht, die Mächte}.
\end{attention}

\courssub{Champ 3 — Les relations internationales}

\begin{center}
\begin{tabularx}{\linewidth}{|l|l|X|}
\hline
\rowcolor{popblueL} Mot allemand (article + pluriel) & Traduction & Exemple \\
\hline
\all{die Diplomatie (nur Sing.)} & la diplomatie & \all{Die Diplomatie kann Kriege verhindern.} \\
\hline
\all{der Botschafter, die Botschafter} & l'ambassadeur & \all{Der Botschafter wurde einbestellt.} \\
\hline
\all{die Vereinten Nationen (Pl.)} & les Nations unies & \all{Die Vereinten Nationen tagen in New York.} \\
\hline
\all{der Vertrag, die Verträge} & le traité, le contrat & \all{Beide Staaten unterzeichneten einen Vertrag.} \\
\hline
\all{die Hilfe, die Hilfen} & l'aide & \all{Die humanitäre Hilfe kommt an.} \\
\hline
\all{verhandeln (mit + Dat)} & négocier (avec) & \all{Die Regierungen verhandeln über Frieden.} \\
\hline
\all{kritisieren} & critiquer & \all{Viele Länder kritisieren die Sanktionen.} \\
\hline
\end{tabularx}
\end{center}

Notez l'abréviation : \all{die UNO} (\all{die Organisation der Vereinten Nationen}) = l'ONU. Elle est féminine, comme toutes les abréviations dont le mot principal est féminin (\all{die Organisation}).

\courssub{Champ 4 — Droits et conflits}

\begin{center}
\begin{tabularx}{\linewidth}{|l|l|X|}
\hline
\rowcolor{popblueL} Mot allemand (article + pluriel) & Traduction & Exemple \\
\hline
\all{die Menschenrechte (Pl.)} & les droits de l'homme & \all{Die Menschenrechte müssen geschützt werden.} \\
\hline
\all{der Konflikt, die Konflikte} & le conflit & \all{Der Konflikt dauert seit Jahren.} \\
\hline
\all{der Krieg, die Kriege} & la guerre & \all{Der Krieg zerstört das Land.} \\
\hline
\all{der Frieden (nur Sing.)} & la paix & \all{Alle Menschen wünschen sich Frieden.} \\
\hline
\all{verletzen} & violer, blesser & \all{Das Regime verletzt die Menschenrechte.} \\
\hline
\all{schützen} & protéger & \all{Die UNO will die Zivilbevölkerung schützen.} \\
\hline
\all{verhindern} & empêcher & \all{Diplomatie kann einen Krieg verhindern.} \\
\hline
\end{tabularx}
\end{center}

\begin{attention}[Faux amis et prépositions]
\all{verletzen} signifie « violer » (un droit) ou « blesser » (une personne), jamais « léser » au sens financier. \all{der Frieden} est masculin alors que « la paix » est féminin. Attention aussi à la préposition : \all{verhandeln \textbf{mit} jemandem \textbf{über} etwas (Akk)} = négocier avec quelqu'un au sujet de quelque chose.
\end{attention}

\courssub{Familles de mots : construire son vocabulaire en réseau}

Observer les familles de mots permet de multiplier son lexique sans effort. Voici les quatre familles essentielles de notre thème :

\begin{center}
\begin{tabularx}{\linewidth}{|l|X|X|}
\hline
\rowcolor{popblueL} Base & Dérivés (article + pluriel) & Exemple \\
\hline
\all{mischen} (mélanger) & \all{die Einmischung, -en}, \all{sich einmischen} & \all{Er mischt sich in die Politik ein.} \\
\hline
\all{die Macht, die Mächte} (le pouvoir) & \all{mächtig} (puissant), \all{der Machthaber, die Machthaber} (le détenteur du pouvoir) & \all{Die Machthaber fürchten Kritik.} \\
\hline
\all{der Staat, die Staaten} (l'État) & \all{staatlich} (étatique), \all{die Staatsangehörigkeit, -en} (la nationalité) & \all{Er hat die kamerunische Staatsangehörigkeit.} \\
\hline
\all{souverän} (souverain) & \all{die Souveränität (nur Sing.)} & \all{Der Staat verteidigt seine Souveränität.} \\
\hline
\end{tabularx}
\end{center}

\begin{popfigure}
\begin{center}
\begin{tikzpicture}[mindmap, grow cyclic, every node/.style={concept, align=center, font=\small},
root concept/.append style={concept color=popblue, font=\bfseries},
level 1/.append style={level distance=3.4cm, sibling angle=90},
level 2/.append style={level distance=2.4cm, sibling angle=45}]
\node[root concept]{die Einmischung}
  child[concept color=poporange] { node {der Staat}
    child { node {die Regierung} }
    child { node {die Macht} } }
  child[concept color=popgreen] { node {der Einfluss}
    child { node {der Druck} }
    child { node {die Sanktion} } }
  child[concept color=poppurple] { node {die Diplomatie}
    child { node {der Vertrag} }
    child { node {die UNO} } }
  child[concept color=poppink] { node {der Konflikt}
    child { node {der Krieg} }
    child { node {der Frieden} } };
\end{tikzpicture}
\end{center}
\legende{Carte mentale : le champ lexical de l'ingérence organisé en quatre branches.}
\end{popfigure}

\courssub{Expressions utiles à réemployer}

Ces expressions figées sont indispensables pour le commentaire de texte et l'expression écrite au Bac :

\begin{exemplebox}[Trois expressions clés]
\begin{itemize}
\item \all{Einfluss ausüben auf (+ Akk)} = exercer une influence sur.\\
\all{Der Staat übt Einfluss auf die Medien aus.} « L'État exerce une influence sur les médias. » (Attention : verbe à particule séparable, \all{aus} va en fin de phrase.)
\item \all{unter Druck stehen} = subir des pressions.\\
\all{Die Regierung steht unter internationalem Druck.} « Le gouvernement subit des pressions internationales. »
\item \all{in die inneren Angelegenheiten eingreifen} = intervenir dans les affaires intérieures.\\
\all{Kein Staat darf in die inneren Angelegenheiten eines anderen Landes eingreifen.} « Aucun État ne doit intervenir dans les affaires intérieures d'un autre pays. »
\end{itemize}
Et un exemple de presse : \all{Die UNO kritisierte die Einmischung.} « L'ONU a critiqué l'ingérence. »
\end{exemplebox}

\begin{savaistu}[La neutralité : refuser l'ingérence par principe]
La Suisse et l'Autriche sont des États neutres : \all{die Neutralität} (la neutralité) signifie qu'ils refusent toute ingérence militaire dans les conflits étrangers et ne participent à aucune alliance militaire. La Suisse, par exemple, n'est membre de l'ONU que depuis 2002, précisément pour préserver sa neutralité. Cela n'empêche pas ces pays d'accueillir des organisations internationales : Genève est un grand centre de la diplomatie mondiale.
\end{savaistu}

\begin{methode}[Comment mémoriser durablement le lexique]
\begin{enumerate}
\item Classez chaque mot nouveau dans un champ lexical (ingérence, État, diplomatie, conflits).
\item Notez-le toujours avec article et pluriel : \all{der Vertrag, die Verträge}.
\item Construisez la famille de mots : verbe, nom, adjectif (\all{mischen} → \all{die Einmischung} → \all{sich einmischen}).
\item Rédigez une phrase personnelle par mot, liée à votre réalité : \all{Kamerun ist ein souveränes Land.}
\item Réutilisez le mot à l'oral dans un débat ou un résumé : un mot employé trois fois est un mot acquis.
\end{enumerate}
\end{methode}

\begin{exoresolu}[Complétez avec le mot qui convient]
Énoncé — Complétez les phrases avec : \all{die Souveränität}, \all{der Einfluss}, \all{die Sanktionen}, \all{verhandeln}, \all{die Menschenrechte}.
\begin{enumerate}
\item \all{Die UNO verhängt \dots\ gegen das Land.}
\item \all{Jeder Staat verteidigt seine \dots .}
\item \all{Die beiden Regierungen \dots\ über einen Friedensvertrag.}
\item \all{Die Medien haben einen großen \dots\ auf die Wähler.}
\item \all{Die Organisation kämpft für \dots\ in der ganzen Welt.}
\end{enumerate}
\tcblower
Solution détaillée :
\begin{enumerate}
\item \all{die Sanktionen} — après \all{verhängen} (imposer), on attend un pluriel à l'accusatif : « L'ONU impose des sanctions au pays. »
\item \all{Souveränität} — possessif \all{seine} + nom féminin abstrait : « Chaque État défend sa souveraineté. »
\item \all{verhandeln} — sujet pluriel (\all{die beiden Regierungen}), verbe conjugué en 2\up{e} position : « Les deux gouvernements négocient un traité de paix. »
\item \all{Einfluss} — \all{einen großen} + accusatif masculin : « Les médias ont une grande influence sur les électeurs. »
\item \all{die Menschenrechte} — nom usuellement au pluriel : « L'organisation lutte pour les droits de l'homme dans le monde entier. »
\end{enumerate}
\end{exoresolu}

\begin{exercice}[À vous de jouer]
Traduisez en allemand : 1. « Le gouvernement critique l'ingérence dans ses affaires intérieures. » 2. « L'ONU exerce une influence sur les conflits. » 3. « Les deux pays négocient un traité de paix. » Puis rédigez deux phrases personnelles avec \all{die Souveränität} et \all{unter Druck stehen}.
\end{exercice}

\coursec{Grammaire 1 : le passif}

Dans la section précédente, nous avons construit le lexique du pouvoir et de la souveraineté : \all{die Einmischung}, \all{die Souveränität}, \all{der Einfluss}, \all{der Staat}. En relisant notre texte support, une forme verbale revient sans cesse : \all{Menschenrechte werden verletzt}, \all{freie Wahlen wurden organisiert}, \all{das Land ist von Sanktionen betroffen worden}. Cette forme, c'est le \cle{passif} (\all{das Passiv}), la structure grammaticale reine des textes politiques : elle permet de mettre en avant ce qui \emph{subit} l'action (les droits, le peuple, le pays) plutôt que celui qui la commet. C'est exactement l'outil qu'il faut pour parler d'ingérence.

\courssub{Observation : partir du texte}

\begin{textebox}[Extraits du texte support]
„In vielen Ländern \textbf{werden} Menschenrechte \textbf{verletzt}. Freie Wahlen \textbf{wurden} von der internationalen Gemeinschaft \textbf{organisiert}, aber die Ergebnisse \textbf{wurden} oft nicht \textbf{anerkannt}. Das Land \textbf{ist} von Sanktionen \textbf{betroffen worden}. Über die Einmischung \textbf{wurde} lange \textbf{diskutiert}.“

\emph{« Dans de nombreux pays, les droits de l'homme sont bafoués. Des élections libres ont été organisées par la communauté internationale, mais les résultats n'ont souvent pas été reconnus. Le pays a été touché par des sanctions. On a longuement discuté de l'ingérence. »}
\end{textebox}

Observons attentivement :

\begin{itemize}
\item Chaque forme passive contient \textbf{deux éléments} : une forme du verbe \all{werden} (\all{werden}, \all{wurden}, \all{ist ... worden}) et un \cle{participe II} (\all{verletzt}, \all{organisiert}, \all{anerkannt}, \all{betroffen}, \all{diskutiert}).
\item Le participe II se trouve \textbf{en fin de phrase} (ou en fin de proposition) : c'est la fameuse « cadre » allemande (\all{Satzklammer}).
\item L'auteur de l'action, quand il est mentionné, est introduit par \all{von} : \all{von der internationalen Gemeinschaft}, \all{von Sanktionen}.
\item Le sujet grammatical n'est \textbf{pas} celui qui agit : dans \all{Menschenrechte werden verletzt}, ce sont les droits de l'homme qui \emph{subissent} la violation.
\end{itemize}

\begin{definition}[Le passif (\all{das Passiv})]
Le passif est une forme verbale dans laquelle le sujet grammatical \emph{subit} l'action au lieu de la faire. Il se forme avec l'auxiliaire \all{werden} conjugué + le participe II du verbe principal placé en fin de phrase. L'agent (celui qui agit) est introduit par \all{von} (+ datif) ou \all{durch} (+ accusatif), ou bien il est simplement tu.
\end{definition}

\courssub{Formation du passif}

\begin{propriete}[Règle de formation]
\textbf{Actif :} sujet (agent) + verbe conjugué + objet accusatif.\\
\textbf{Passif :} objet devenu sujet (nominatif) + \all{werden} conjugué + participe II en fin de phrase.

\begin{itemize}
\item Präsens : \all{werden} conjugué au présent + Partizip II → \all{Die Souveränität \textbf{wird verletzt}.}
\item Präteritum : \all{wurden} + Partizip II → \all{Die Souveränität \textbf{wurde verletzt}.}
\item Perfekt : \all{sein} conjugué + Partizip II + \all{\textbf{worden}} → \all{Die Souveränität \textbf{ist verletzt worden}.}
\item Plusquamperfekt : \all{war} + Partizip II + \all{worden} → \all{Die Souveränität \textbf{war verletzt worden}.}
\end{itemize}

\textbf{Point capital :} au Perfekt et au Plusquamperfekt du passif, on utilise \all{sein} + participe II + \all{worden} (participe II « raccourci » de \all{werden}), \textbf{jamais} \all{geworden}.
\end{propriete}

\begin{attention}[L'erreur n° 1 des francophones]
Ne jamais écrire \emph{« ist verletzt geworden »}. Le participe II de \all{werden} employé comme auxiliaire du passif perd son \all{ge-} : la seule forme correcte est \all{\textbf{ist verletzt worden}}. De même : \all{war organisiert worden} (et non \emph{geworden}). Astuce : dès qu'il y a déjà un participe II dans la phrase, \all{werden} se contente de \all{worden}.
\end{attention}

\courssub{Tableau des temps du passif}

\begin{center}
\begin{tabularx}{\linewidth}{|l|X|X|}
\hline
\rowcolor{popblueL} \textbf{Temps} & \textbf{Formation} & \textbf{Exemple (traduction)} \\
\hline
Präsens & wird/werden + Partizip II & \all{Menschenrechte werden verletzt.} « Les droits de l'homme sont bafoués. » \\
\hline
Präteritum & wurde/wurden + Partizip II & \all{Freie Wahlen wurden organisiert.} « Des élections libres furent organisées. » \\
\hline
Perfekt & ist/sind + Partizip II + worden & \all{Das Land ist von Sanktionen betroffen worden.} « Le pays a été touché par des sanctions. » \\
\hline
Plusquamperfekt & war/waren + Partizip II + worden & \all{Der Vertrag war unterzeichnet worden.} « Le traité avait été signé. » \\
\hline
Futur I & wird/werden + Partizip II + werden & \all{Neue Wahlen werden organisiert werden.} « De nouvelles élections seront organisées. » \\
\hline
\end{tabularx}
\end{center}

\begin{exemplebox}[Exemples traduits]
\begin{itemize}
\item \all{Die Souveränität wurde verletzt.} « La souveraineté a été violée. »
\item \all{Das Land ist von Sanktionen betroffen worden.} « Le pays a été touché par des sanctions. »
\item \all{Der Minister wurde von der Presse kritisiert.} « Le ministre a été critiqué par la presse. »
\item \all{Die Grenzen werden durch Soldaten kontrolliert.} « Les frontières sont contrôlées par des soldats. »
\item \all{Viele Demonstranten waren verhaftet worden.} « De nombreux manifestants avaient été arrêtés. »
\end{itemize}
\end{exemplebox}

\courssub{L'agent : \all{von} ou \all{durch} ?}

\begin{propriete}[Introduire l'agent]
\begin{itemize}
\item \all{\textbf{von}} + \textbf{datif} : l'agent, l'auteur de l'action (personne, institution) : \all{Der Staat wird \textbf{von der UNO} kritisiert.} « L'État est critiqué par l'ONU. »
\item \all{\textbf{durch}} + \textbf{accusatif} : le moyen, l'instrument, la cause : \all{Die Stadt wurde \textbf{durch eine Flut} zerstört.} « La ville fut détruite par une inondation. » ; \all{Der Einfluss wird \textbf{durch die Medien} verstärkt.} « L'influence est renforcée par les médias. »
\end{itemize}
Repère pratique : \all{von} = « qui agit » ; \all{durch} = « par quel moyen ». En cas de doute au Bac, \all{von} est le choix le plus fréquent pour un agent humain ou institutionnel.
\end{propriete}

\courssub{Le passif sans agent : le fameux « on » français}

Très souvent, l'agent n'a pas d'importance ou est inconnu. L'allemand utilise alors le passif impersonnel, que le français rend presque toujours par « on » :

\begin{exemplebox}[Passif impersonnel]
\begin{itemize}
\item \all{Es wurde viel diskutiert.} « On a beaucoup discuté. »
\item \all{Es wird über die Einmischung gesprochen.} « On parle de l'ingérence. »
\item \all{In der Schule wird Deutsch gelernt.} « On apprend l'allemand à l'école. »
\item \all{Es wurde bis Mitternacht demonstriert.} « On a manifesté jusqu'à minuit. »
\end{itemize}
Remarque : ce passif impersonnel fonctionne même avec des verbes intransitifs (\all{diskutieren}, \all{sprechen}, \all{demonstrieren}), ce qui est impossible en français avec un vrai passif.
\end{exemplebox}

\courssub{Transformation actif $\leftrightarrow$ passif}

\begin{methode}[Passer de l'actif au passif en trois étapes]
Prenons : \all{Die UNO kritisiert den Staat.} « L'ONU critique l'État. »
\begin{enumerate}
\item L'\textbf{objet accusatif} de la phrase active devient le \textbf{sujet nominatif} de la phrase passive : \all{den Staat} → \all{der Staat}.
\item Le verbe conjugué devient \all{werden} (au même temps) + \textbf{participe II en fin de phrase} : \all{kritisiert} → \all{wird kritisiert}.
\item L'ancien sujet devient un complément d'agent introduit par \all{von} + datif : \all{die UNO} → \all{von der UNO}.
\end{enumerate}
Résultat : \all{\textbf{Der Staat wird von der UNO kritisiert.}} « L'État est critiqué par l'ONU. »

Sens inverse (passif → actif) : le complément \all{von} + datif redevient sujet nominatif, le sujet passif redevient objet accusatif, et le verbe reprend sa forme simple.
\end{methode}

\begin{popfigure}
\begin{tikzpicture}[font=\small, node distance=0.4cm]
\node[draw=popblue, fill=popblueL, rounded corners, align=center, text width=4.2cm] (aktiv) {\textbf{Aktiv}\\ Subjekt (Nominativ)\\ $\rightarrow$ Verb\\ $\rightarrow$ Objekt (Akkusativ)\\[2pt] \all{Die UNO kritisiert den Staat.}};
\node[draw=poporange, fill=poporangeL, rounded corners, align=center, text width=4.6cm, right=2.2cm of aktiv] (passiv) {\textbf{Passiv}\\ Objekt $\rightarrow$ Subjekt\\ werden + Partizip II (Ende)\\ von + Dativ (Agent)\\[2pt] \all{Der Staat wird von der UNO kritisiert.}};
\draw[-{Stealth[length=3mm]}, line width=1.5pt, popdark] (aktiv) -- (passiv) node[midway, above, align=center, text=popdark]{\textbf{Transformation}};
\draw[-{Stealth[length=2mm]}, popgreen, line width=1pt] ([yshift=-0.15cm]aktiv.south) .. controls +(0,-1.1) and +(0,-1.1) .. ([yshift=-0.15cm]passiv.south) node[midway, below, align=center, text=popgreen]{\footnotesize den Staat $\rightarrow$ der Staat};
\end{tikzpicture}
\legende{Schéma de transformation : de la phrase active à la phrase passive.}
\end{popfigure}

\courssub{Comparaison avec le français}

\begin{center}
\begin{tabularx}{\linewidth}{|X|X|X|}
\hline
\rowcolor{popblueL} \textbf{Français} & \textbf{Deutsch} & \textbf{Remarque} \\
\hline
On a beaucoup discuté. & \all{Es wurde viel diskutiert.} & Le français préfère « on » ; l'allemand préfère le passif. \\
\hline
On construit une nouvelle école à Yaoundé. & \all{In Yaoundé wird eine neue Schule gebaut.} & « On » + actif = passif allemand. \\
\hline
Les élections ont été organisées par l'ONU. & \all{Die Wahlen wurden von der UNO organisiert.} & Passif dans les deux langues : correspondance directe. \\
\hline
Le ministre a été critiqué. & \all{Der Minister ist kritisiert worden.} & Attention : \all{worden}, pas \emph{geworden}. \\
\hline
On m'a volé mon téléphone. & \all{Mir wurde mein Telefon gestohlen.} & L'allemand garde le datif (\all{mir}) : passif dit « datif ». \\
\hline
\end{tabularx}
\end{center}

Le passif allemand est \textbf{plus fréquent} que le passif français. En français, on évite souvent le passif (« on a organisé des élections ») ; en allemand, le passif est naturel et très employé, surtout dans la presse et les textes officiels : \all{Es wurden freie Wahlen organisiert.} En traduction (version), un passif allemand peut donc se rendre par « on » + actif, et réciproquement (thème) un « on » français appelle souvent un passif allemand.

\begin{attention}[Pièges de traduction]
\begin{itemize}
\item \all{werden} au passif n'est \textbf{pas} le futur : \all{Es wird diskutiert.} = « On discute. » (présent passif), et non « on discutera ». Le futur passif exige deux \all{werden} : \all{Es wird diskutiert werden.}
\item Ne pas confondre \all{sein} + participe II (passif d'action au Perfekt : \all{ist verletzt worden}) avec le passif d'état (\all{Zustandspassiv}) : \all{Die Tür ist geschlossen.} « La porte est fermée » (état, pas action). Au lycée, seul le passif d'action avec \all{werden} est exigé.
\item Le participe II reste \textbf{invariable} : pas d'accord comme en français (« les élections ont été organisé\emph{es} » mais \all{die Wahlen wurden organisiert}).
\end{itemize}
\end{attention}

\courssub{Texte d'application : le passif en contexte}

\begin{textebox}[Einmischung oder Hilfe?]
„Nach der Wahl wurde in der Hauptstadt viel demonstriert. Die Ergebnisse wurden von der Opposition nicht anerkannt, und es wurde behauptet, dass Stimmen gestohlen worden waren. Bald wurde das Land von der internationalen Gemeinschaft kritisiert: Die Souveränität des Volkes sei verletzt worden, hieß es in den Zeitungen. Daraufhin wurden Sanktionen beschlossen, und die Grenzen für bestimmte Politiker wurden geschlossen. In den Medien wurde heftig diskutiert: Ist das eine berechtigte Hilfe oder eine unerlaubte Einmischung? Ein Journalist schrieb: „Wenn Wahlen von außen kontrolliert werden, wird der Staat geschwächt.“ Andere meinten, dass Menschenrechte geschützt werden müssen, auch wenn dabei in die Politik eines Landes eingegriffen wird. Am Ende wurde ein Dialog zwischen den Parteien organisiert, und es wurde vereinbart, dass neue Wahlen von unabhängigen Beobachtern überwacht werden. So wurde eine Krise vermieden, aber die Frage bleibt: Wann wird Hilfe zur Einmischung?“

\emph{Glossaire :} die Wahl (-en) : l'élection ; anerkennen (trennbar) : reconnaître ; die Stimme (-n) : la voix (électorale) ; behaupten : affirmer ; beschließen (beschloss, beschlossen) : décider ; heftig : violemment ; berechtigt : légitime ; unerlaubt : illicite ; schwächen : affaiblir ; eingreifen (griff ein, eingegriffen) : intervenir ; vereinbaren : convenir ; unabhängig : indépendant ; der Beobachter (-) : l'observateur ; überwachen : surveiller ; vermeiden (vermied, vermieden) : éviter.
\end{textebox}

\textbf{Questions de compréhension :}
\begin{enumerate}
\item Wer hat die Wahlergebnisse nicht anerkannt ?
\item Welche Maßnahmen wurden von der internationalen Gemeinschaft beschlossen ?
\item Wie wurde die Krise am Ende vermieden ?
\item Relevez cinq formes passives du texte et indiquez leur temps.
\end{enumerate}

\courssub{Exercice résolu et entraînement}

\begin{exoresolu}[Transformation actif $\rightarrow$ passif]
Énoncé : Mettez les phrases suivantes au passif, au même temps. 1) \all{Die internationale Gemeinschaft kritisiert den Staat.} 2) \all{Die Opposition organisierte eine Demonstration.} 3) \all{Sanktionen haben die Wirtschaft geschwächt.}
\tcblower
Solution détaillée :
\begin{enumerate}
\item Objet accusatif : \all{den Staat} → sujet : \all{der Staat}. Verbe au Präsens : \all{kritisiert} → \all{wird kritisiert}. Ancien sujet : \all{von der internationalen Gemeinschaft}. Résultat : \all{\textbf{Der Staat wird von der internationalen Gemeinschaft kritisiert.}} « L'État est critiqué par la communauté internationale. »
\item \all{eine Demonstration} (accusatif) → \all{eine Demonstration} (nominatif, féminin, forme identique). Präteritum : \all{wurde organisiert}. Résultat : \all{\textbf{Eine Demonstration wurde von der Opposition organisiert.}} « Une manifestation fut organisée par l'opposition. »
\item Perfekt passif : \all{sein} + participe II + \all{worden}. Résultat : \all{\textbf{Die Wirtschaft ist durch Sanktionen geschwächt worden.}} « L'économie a été affaiblie par des sanctions. » (Ici \all{durch} convient car les sanctions sont le moyen.)
\end{enumerate}
\end{exoresolu}

\begin{exercice}[À vous de jouer]
\begin{enumerate}
\item Traduisez en allemand : a) « Les droits de l'homme sont bafoués. » b) « On a beaucoup discuté de la souveraineté. » c) « Le traité avait été signé par les deux États. »
\item Mettez au passif (Präsens) : \all{Die Medien beeinflussen die Wähler.}
\item Mettez au Perfekt passif : \all{Man organisierte freie Wahlen.}
\end{enumerate}
\end{exercice}

\begin{corrige}[Exercice « À vous de jouer »]
\begin{enumerate}
\item a) \all{Die Menschenrechte werden verletzt.} b) \all{Es wurde viel über die Souveränität diskutiert.} c) \all{Der Vertrag war von den beiden Staaten unterzeichnet worden.}
\item \all{Die Wähler werden von den Medien beeinflusst.}
\item \all{Freie Wahlen sind organisiert worden.} (passif impersonnel possible aussi : \all{Es wurden freie Wahlen organisiert.})
\end{enumerate}
\end{corrige}

\begin{savaistu}[Le passif au Bac]
Au baccalauréat, le passif est l'une des structures les plus testées, en version comme en thème. Deux réflexes gagnants : en version, dès que tu repères \all{von} ou \all{durch} suivi d'un nom, cherche un participe II en fin de phrase — c'est presque toujours un passif à traduire par « être + participe » ou par « on ». En thème, chaque fois que le français utilise « on », demande-toi si un passif allemand ne serait pas plus élégant : \all{On a organisé des élections} → \all{Es wurden Wahlen organisiert.} Et souviens-toi du mot magique : \all{worden}, jamais \emph{geworden} !
\end{savaistu}

\begin{aretenir}[L'essentiel du passif]
\begin{itemize}
\item Formation : \all{werden} conjugué + participe II en fin de phrase.
\item Perfekt/Plusquamperfekt : \all{ist/war} + participe II + \all{worden}.
\item Agent : \all{von} + datif (qui agit) ; \all{durch} + accusatif (le moyen).
\item Passif impersonnel : \all{Es wurde diskutiert.} = « On a discuté. »
\item Transformation : accusatif → nominatif ; verbe → \all{werden} + participe II ; sujet → \all{von} + datif.
\end{itemize}
\end{aretenir}

\coursec{Grammaire 2 : plus-que-parfait et subjonctif II}

Après avoir analysé la voix passive, qui permet de présenter l'ingérence comme une action subie par un État (\all{„Die Wahlen wurden manipuliert.“}), nous devons maintenant maîtriser deux outils grammaticaux indispensables pour \textbf{raconter le passé avec précision} (l'antériorité) et pour \textbf{exprimer l'hypothèse, le regret ou la critique} (l'irréel) : le \cle{plus-que-parfait} (\all{das Plusquamperfekt}) et le \cle{subjonctif II} (\all{der Konjunktiv II}).

\courssub{Le plus-que-parfait (das Plusquamperfekt) : forme et emploi}

\textbf{Observation dans le texte}\par

\begin{textebox}[Extrait du reportage sur l'ingérence électorale]
„Nachdem die Beobachter monatelang \textbf{berichtet hatten}, dass die Medien zensiert \textbf{wurden}, \textbf{entschied} die internationale Gemeinschaft, Sanktionen zu verhängen. Ein Land, das jahrelang unter Druck \textbf{gestanden hatte}, erhielt endlich freie Wahlen.“
\tcblower
\textbf{Glossaire :} \all{der Beobachter, -} (l'observateur) ; \all{monatelang} (pendant des mois) ; \all{berichten} (rapporter, signaler) ; \all{zensiert werden} (être censuré) ; \all{entscheiden, entschied, hat entschieden} (décider) ; \all{die Sanktion, -en} (la sanction) ; \all{verhängen} (imposer, décréter) ; \all{unter Druck stehen} (être sous pression) ; \all{endlich} (enfin) ; \all{die freie Wahl, -en} (l'élection libre).
\end{textebox}

\begin{propriete}[Formation du plus-que-parfait]
Le plus-que-parfait se forme avec l'auxiliaire \cle{hatte} (pour \all{haben}) ou \cle{war} (pour \all{sein}) au \textbf{Präteritum}, suivi du \cle{Partizip II} (participe passé) du verbe principal, placé en fin de proposition.
\begin{center}
\begin{tabular}{|l|l|l|}
\hline
\rowcolor{popblueL} \textbf{Personne} & \textbf{Aux. \all{haben} (\all{machen})} & \textbf{Aux. \all{sein} (\all{kommen})} \\ \hline
ich & hatte gemacht & war gekommen \\ \hline
du & hattest gemacht & warst gekommen \\ \hline
er/sie/es & hatte gemacht & war gekommen \\ \hline
wir & hatten gemacht & waren gekommen \\ \hline
ihr & hattet gemacht & wart gekommen \\ \hline
sie/Sie & hatten gemacht & waren gekommen \\ \hline
\end{tabular}
\end{center}
\textbf{Règle du choix de l'auxiliaire} (identique à celle du Perfekt) :
\begin{itemize}
    \item \textbf{\all{haben}} : verbes transitifs (\all{lesen, kaufen, berichten}), verbes d'état ou de durée (\all{stehen, liegen, dauern, leben}), verbes réfléchis (\all{sich entscheiden}), verbes modaux.
    \item \textbf{\all{sein}} : verbes de \textbf{mouvement} vers un lieu (\all{gehen, kommen, fahren, fliegen, laufen, reisen}), verbes de \textbf{changement d'état} (\all{sterben, einschlafen, aufwachen, werden, passieren, geschehen}), ainsi que \all{sein} et \all{bleiben}.
\end{itemize}
\end{propriete}

\begin{methode}[Choisir le bon auxiliaire en 3 étapes]
\begin{enumerate}
    \item Identifier le verbe principal sous sa forme de participe (ex. \all{gestanden}, \all{berichtet}, \all{gekommen}).
    \item Se poser la question : \emph{« Y a-t-il déplacement d'un lieu A vers un lieu B, ou changement d'état ? »}
    \item \textbf{OUI} $\rightarrow$ auxiliaire \all{sein} (\all{kommen, gehen, sterben, passieren}). \textbf{NON} $\rightarrow$ auxiliaire \all{haben} (\all{stehen, berichten, machen, können}).
\end{enumerate}
\textbf{Liste à mémoriser (auxiliaire \all{sein}) :} \all{gehen, kommen, fahren, fliegen, laufen, reisen, sterben, werden, passieren, geschehen, bleiben, sein, aufstehen, einschlafen, aufwachen}.
\end{methode}

\begin{propriete}[Emploi : l'antériorité par rapport au passé]
Le plus-que-parfait exprime un fait \textbf{antérieur} à un autre fait déjà situé dans le passé (exprimé au Präteritum ou au Perfekt). Il est \textbf{obligatoire} après la conjonction \cle{\all{nachdem}} (« après que ») lorsque la principale est au passé. Attention : dans la subordonnée introduite par \all{nachdem}, le verbe conjugué se place \textbf{en fin de proposition}.
\begin{center}
\begin{tabularx}{\linewidth}{|X|X|}
\hline
\rowcolor{popblueL} \textbf{Français} & \textbf{Deutsch} \\ \hline
Après que les élections \textbf{eurent eu lieu}, les observateurs partirent. & \all{Nachdem die Wahlen \textbf{stattgefunden hatten}, zogen die Beobachter ab.} \\ \hline
Le pays \textbf{avait subi} des pressions pendant des années. & \all{Das Land \textbf{hatte jahrelang unter Druck gestanden}.} \\ \hline
\end{tabularx}
\end{center}
\end{propriete}

\begin{attention}[Pièges francophones : \all{nachdem} et l'auxiliaire]
\begin{itemize}
    \item \textbf{Erreur :} \all{„Nachdem die Wahlen stattfanden, ...“} (Präteritum après \all{nachdem}). \textbf{Correct :} \all{„Nachdem die Wahlen \textbf{stattgefunden hatten}, ...“} L'antériorité doit être marquée grammaticalement en allemand.
    \item \textbf{Erreur :} \all{„Das Land hat unter Druck gestanden.“} dans un récit au passé. \textbf{Correct :} \all{„Das Land \textbf{hatte} unter Druck \textbf{gestanden}.“}
    \item \textbf{Verbes trompeurs :} \all{stehen} (rester debout, durer) prend \textbf{haben} : \all{hatte gestanden} ; mais \all{passieren} (se produire) prend \textbf{sein} : \all{war passiert}.
    \item \textbf{Faux ami :} \all{beobachten} signifie « observer, surveiller », et non « obtenir ».
\end{itemize}
\end{attention}

\begin{exemplebox}[Exemples traduits : le plus-que-parfait en contexte]
\begin{itemize}
    \item \all{„Nachdem der Diktator die Macht \textbf{übernommen hatte}, verbot er alle Oppositionsparteien.“} \\
    « Après que le dictateur \textbf{eut pris} le pouvoir, il interdit tous les partis d'opposition. »
    \item \all{„Die Journalisten \textbf{hatten} monatelang \textbf{recherchiert}, bevor sie den Bericht veröffentlichten.“} \\
    « Les journalistes \textbf{avaient enquêté} pendant des mois avant de publier le rapport. »
    \item \all{„Wir \textbf{waren} schon \textbf{abgereist}, als der Konflikt ausbrach.“} \\
    « Nous \textbf{étions déjà partis} quand le conflit éclata. » (\all{abreisen} = mouvement $\rightarrow$ \all{sein})
\end{itemize}
\end{exemplebox}

\courssub{Le subjonctif II au passé (Konjunktiv II Vergangenheit) : l'irréel du passé}

\textbf{Observation : le regret et l'hypothèse contrefactuelle}\par

\begin{exemplebox}[Phrases-clés du débat public]
\begin{itemize}
    \item \all{„\textbf{Wäre} die Welt früher \textbf{eingeschritten}, \textbf{hätte} man viel Leid \textbf{verhindern können}.“}
    \item \all{„\textbf{Hätte} die Opposition gemeinsam \textbf{kandidiert}, \textbf{wäre} sie \textbf{siegreich gewesen}.“}
\end{itemize}
\tcblower
\textbf{Traduction :} « Si le monde \textbf{était intervenu} plus tôt, on \textbf{aurait pu} éviter bien des souffrances. » / « Si l'opposition \textbf{s'était présentée} unie, elle \textbf{aurait été} victorieuse. »
\end{exemplebox}

\begin{propriete}[Formation du subjonctif II au passé]
On utilise le subjonctif II de l'auxiliaire (\cle{\all{hätte}} ou \cle{\all{wäre}}) + \cle{Partizip II} du verbe principal.
\begin{center}
\begin{tabular}{|l|l|l|}
\hline
\rowcolor{popblueL} \textbf{Personne} & \textbf{Aux. \all{haben} (\all{machen})} & \textbf{Aux. \all{sein} (\all{kommen})} \\ \hline
ich & hätte gemacht & wäre gekommen \\ \hline
du & hättest gemacht & wärst gekommen \\ \hline
er/sie/es & hätte gemacht & wäre gekommen \\ \hline
wir & hätten gemacht & wären gekommen \\ \hline
ihr & hättet gemacht & wärt gekommen \\ \hline
sie/Sie & hätten gemacht & wären gekommen \\ \hline
\end{tabular}
\end{center}
\textbf{Choix de l'auxiliaire :} \textbf{mêmes règles} qu'au Plusquamperfekt et au Perfekt (\all{sein} pour le mouvement et le changement d'état, sinon \all{haben}).
\end{propriete}

\begin{propriete}[Emploi : l'irréel du passé (contrefactuel)]
Le subjonctif II passé exprime ce qui \textbf{n'a pas eu lieu} mais \textbf{aurait pu avoir lieu} si une condition avait été réalisée. Il correspond au français « si + plus-que-parfait $\rightarrow$ conditionnel passé ». Deux structures sont possibles :
\begin{itemize}
    \item \textbf{Avec \all{wenn} :} \all{Wenn die UNO früher \textbf{eingegriffen hätte}, \textbf{wäre} der Krieg \textbf{verhindert worden}.} (Si l'ONU était intervenue plus tôt, la guerre aurait été évitée.)
    \item \textbf{Sans \all{wenn} (inversion, verbe en position 1) :} \all{\textbf{Hätte} die UNO früher \textbf{eingegriffen}, \textbf{wäre} der Krieg \textbf{verhindert worden}.}
\end{itemize}
\textbf{Passif au KII passé :} \all{wäre + Partizip II + worden} : \all{„Der Krieg wäre verhindert worden.“} (La guerre aurait été évitée.)
\end{propriete}

\begin{attention}[Trois erreurs fatales des francophones]
\begin{enumerate}
    \item \textbf{\all{würde} + Partizip II :} la forme \all{„er würde gekommen sein“} n'existe pas en allemand standard ! Au passé, on utilise obligatoirement \textbf{\all{hätte/wäre + Partizip II}} : \all{er wäre gekommen}, \all{er hätte kommen können}.
    \item \textbf{Oubli de l'inversion :} la suppression de \all{wenn} impose le verbe conjugué en \textbf{position 1} : \all{„Wäre er gekommen, ...“} = \all{„Wenn er gekommen wäre, ...“}. Ne jamais écrire \all{„Wenn wäre er gekommen...“}.
    \item \textbf{Auxiliaire incorrect au passif :} le passif exige \all{sein} : \all{„\textbf{Wäre} der Krieg verhindert worden“} (et non \all{hätte}).
\end{enumerate}
\end{attention}

\begin{exemplebox}[Comparaison Français / Allemand : « si + plus-que-parfait »]
\begin{center}
\begin{tabularx}{\linewidth}{|X|X|}
\hline
\rowcolor{popblueL} \textbf{Français} & \textbf{Deutsch} \\ \hline
Si l'État \textbf{avait protégé} les minorités, il n'y \textbf{aurait pas eu} de violences. & \all{Hätte der Staat die Minderheiten \textbf{geschützt}, \textbf{wäre} es \textbf{zu keiner Gewalt gekommen}.} \\ \hline
Si les médias \textbf{avaient été} libres, le peuple \textbf{aurait connu} la vérité. & \all{Wären die Medien \textbf{frei gewesen}, \textbf{hätte} das Volk die Wahrheit \textbf{erfahren}.} \\ \hline
Si la Cour pénale \textbf{avait agi} plus vite, les criminels \textbf{seraient} en prison. & \all{Hätte der Strafgerichtshof früher \textbf{gehandelt}, \textbf{wären} die Verbrecher \textbf{im Gefängnis}.} \\ \hline
\end{tabularx}
\end{center}
\end{exemplebox}

\courssub{Le subjonctif II au présent (Konjunktiv II Gegenwart) : hypothèse, souhait, politesse}

\begin{propriete}[Formation : \all{würde} + infinitif et formes irrégulières]
Pour la plupart des verbes, on utilise \cle{\all{würde}} + \cle{infinitif}. Mais \all{haben}, \all{sein} et les verbes modaux ont des formes propres, plus fréquentes, qu'il faut connaître par cœur.
\begin{center}
\begin{tabular}{|l|l|l|}
\hline
\rowcolor{popblueL} \textbf{Verbe} & \textbf{Präteritum} & \textbf{Konjunktiv II (présent)} \\ \hline
\all{haben} & hatte & hätte \\ \hline
\all{sein} & war & wäre \\ \hline
\all{werden} & wurde & würde \\ \hline
\all{können} & konnte & könnte \\ \hline
\all{müssen} & musste & müsste \\ \hline
\all{dürfen} & durfte & dürfte \\ \hline
\all{sollen} & sollte & sollte \\ \hline
\all{wollen} & wollte & wollte \\ \hline
\end{tabular}
\end{center}
\textit{Autres verbes :} \all{ich würde machen, du würdest gehen, er würde kommen...}
\end{propriete}

\begin{propriete}[Emplois principaux au niveau Terminale]
\begin{enumerate}
    \item \textbf{Hypothèse présente ou future peu réaliste :} \all{„Es \textbf{wäre} besser, wenn die Mächte sich nicht \textbf{einmischten}.“} (Il serait préférable que les puissances ne s'ingèrent pas.)
    \item \textbf{Souhait ou regret présent :} \all{„Wenn die Wahlen doch \textbf{fair wären}!“} (Si seulement les élections étaient équitables !)
    \item \textbf{Politesse (demande, question atténuée) :} \all{„\textbf{Könnten} Sie mir helfen? \textbf{Würde} das funktionieren?“} (Pourriez-vous m'aider ? Cela fonctionnerait-il ?)
    \item \textbf{Conseil (structure \all{an deiner Stelle}) :} \all{„An deiner Stelle \textbf{würde} ich protestieren.“} (À ta place, je protesterais.)
\end{enumerate}
\end{propriete}

\courssub{Synthèse visuelle : la frise des temps du récit et de l'hypothèse}

\begin{popfigure}
\begin{tikzpicture}[
    scale=0.95,
    every node/.style={font=\small},
    timeaxis/.style={->, thick, popdark},
    tick/.style={thick, popdark},
    event/.style={circle, fill=popblue, text=white, minimum width=1.6cm, align=center, font=\footnotesize\bfseries},
    labelbox/.style={rectangle, rounded corners, fill=popblueL, draw=popblue, inner sep=4pt, align=center, text width=3.6cm, font=\footnotesize}
]
% Axe du temps
\draw[timeaxis] (-0.5,0) -- (14.5,0) node[right, popdark] {\textbf{Temps}};

% Points temporels
\node[event, align=center] (PQP) at (1.2,0){Plusquam-\\perfekt};
\node[event, align=center] (PRAET) at (5.4,0){Präteritum /\\Perfekt};
\node[event] (PRES) at (9.6,0) {Präsens};
\node[event, fill=poppurple, align=center] (KII) at (13.2,0){KII\\Passé};

% Ticks
\draw[tick] (1.2,0.25) -- (1.2,-0.25);
\draw[tick] (5.4,0.25) -- (5.4,-0.25);
\draw[tick] (9.6,0.25) -- (9.6,-0.25);
\draw[tick] (13.2,0.25) -- (13.2,-0.25);

% Flèche d'antériorité
\draw[->, thick, poporange, bend left=25] (PRAET) to node[above, poporange, font=\tiny] {antériorité (\all{nachdem})} (PQP);

% Boîtes d'exemples
\node[labelbox, anchor=north, align=center] at (1.2,-0.6){
    \textbf{Antériorité}\\
    \all{„Nachdem der Druck \textbf{zugenommen hatte}...“}\\
    (Après que la pression eut augmenté...)
};

\node[labelbox, anchor=north, align=center] at (5.4,-0.6){
    \textbf{Point de référence}\\
    \all{„... \textbf{entschied} die UNO einzugreifen.“}\\
    (... l'ONU décida d'intervenir.)
};

\node[labelbox, anchor=north, align=center] at (9.6,-0.6){
    \textbf{Présent / généralités}\\
    \all{„Einmischung \textbf{verletzt} die Souveränität.“}\\
    (L'ingérence viole la souveraineté.)
};

\node[labelbox, anchor=north, align=center] at (13.2,-0.6){
    \textbf{Irréel du passé}\\
    \all{„\textbf{Wäre} die UNO früher \textbf{eingeschritten}, \textbf{wäre} Leid \textbf{erspart geblieben}.“}\\
    (Si l'ONU était intervenue plus tôt, la souffrance aurait été épargnée.)
};

% Connexion KII Passé vers Plusquamperfekt
\draw[dashed, thick, poppurple] (PQP.north) -- ++(0,1.3) -| node[near start, above, poppurple, font=\tiny, align=center] {base formelle du KII passé\\(\all{hätte/wäre} + Partizip II)} (KII.north);
\end{tikzpicture}
\legende{Frise chronologique : le Plusquamperfekt situe l'antériorité par rapport au Präteritum/Perfekt (base du récit). Le Konjunktiv II passé se construit sur la même base formelle (auxiliaire au KII + Partizip II) pour exprimer l'irréel par rapport à ce passé.}
\end{popfigure}

\courssub{Exercice résolu : du constat à l'hypothèse}

\begin{exoresolu}[Transformations : passif, temps du récit, irréel]
\textbf{Consigne :} transformez les phrases suivantes.
\begin{enumerate}
    \item (Actif $\rightarrow$ passif au Präteritum) \all{Die Mächte manipulierten die Wahlen.}
    \item (Relier avec \all{nachdem} + Plusquamperfekt passif) \all{Die Wahlen wurden manipuliert. Dann protestierte die Welt.}
    \item (KII passé : hypothèse contrefactuelle) \all{Die Welt protestierte nicht. Deshalb blieb der Diktator an der Macht.}
\end{enumerate}

\tcblower
\textbf{Solution détaillée :}
\begin{enumerate}
    \item \all{Die Wahlen \textbf{wurden} von den Mächten \textbf{manipuliert}.} \\
    \textit{Analyse :} passif au Präteritum = \all{werden} au Präteritum (\all{wurden}) + Partizip II (\all{manipuliert}) ; l'agent est introduit par \all{von} + datif (\all{von den Mächten}).
    \item \all{Nachdem die Wahlen \textbf{manipuliert worden waren}, protestierte die Welt.} \\
    \textit{Analyse :} passif au Plusquamperfekt = \all{sein} au Präteritum (\all{waren}) + Partizip II (\all{manipuliert}) + \all{worden}. \all{nachdem} impose le Plusquamperfekt et renvoie l'auxiliaire conjugué en fin de subordonnée ; la principale reste au Präteritum (\all{protestierte}).
    \item \all{\textbf{Hätte} die Welt \textbf{protestiert}, \textbf{wäre} der Diktator \textbf{nicht an der Macht geblieben}.} \\
    \textit{Analyse :}
    \begin{itemize}
        \item Protase (condition) : inversion du KII passé sans \all{wenn}. \all{protestieren} est intransitif, sans mouvement $\rightarrow$ auxiliaire \all{haben} $\rightarrow$ \all{hätte protestiert}.
        \item Apodose (conséquence) : KII passé. \all{bleiben} exige l'auxiliaire \all{sein} $\rightarrow$ \all{wäre ... geblieben}. La négation \all{nicht} se place avant le groupe verbal final.
    \end{itemize}
\end{enumerate}
\end{exoresolu}

\courssub{Entraînement : texte à trous et production écrite}

\begin{exercice}[Compléter : le bon temps, la bonne voix, le bon mode]
\textbf{Contexte :} un analyste commente l'échec d'une médiation internationale (situation fictive).
\textbf{Complétez avec la forme correcte du verbe entre parenthèses.}

„Nachdem die Vermittler monatelang \_\_\_\_\_\_\_\_\_\_\_\_ (\all{verhandeln} -- Plusquamperfekt actif), \_\_\_\_\_\_\_\_\_\_\_\_ die Gespräche (\all{scheitern} -- Präteritum). Die Rebellen \_\_\_\_\_\_\_\_\_\_\_\_ (\all{sich weigern} -- Präteritum), die Waffen niederzulegen. \_\_\_\_\_\_\_\_\_\_\_\_ die Regierung früher kompromissbereit \_\_\_\_\_\_\_\_\_\_\_\_ (\all{sein} -- KII passé, inversion sans \all{wenn}), \_\_\_\_\_\_\_\_\_\_\_\_ viele Leben \_\_\_\_\_\_\_\_\_\_\_\_ (\all{retten können} -- KII passé, passif avec modal). Heute \_\_\_\_\_\_\_\_\_\_\_\_ die internationale Gemeinschaft stärkere Sanktionen \_\_\_\_\_\_\_\_\_\_\_\_ (\all{müssen} + \all{verhängen} -- KII présent + infinitif), wenn sie glaubwürdig \_\_\_\_\_\_\_\_\_\_\_\_ (\all{bleiben wollen} -- infinitif).“
\end{exercice}

\begin{corrige}[Exercice : le bon temps, la bonne voix, le bon mode]
\begin{enumerate}
    \item \all{verhandelt hatten} : Plusquamperfekt actif, auxiliaire \all{haben}, obligatoire après \all{nachdem} ; auxiliaire conjugué en fin de subordonnée.
    \item \all{scheiterten} : Präteritum de \all{scheitern} (verbe faible). \textit{Remarque :} \all{scheitern} (échouer) est un verbe d'état/processus qui ne se met pas au passif ; on le conjugue à l'actif : \all{die Gespräche scheiterten} (les pourparlers échouèrent).
    \item \all{weigerten sich} : Präteritum du verbe réfléchi \all{sich weigern} (refuser).
    \item \all{Wäre ... gewesen} : KII passé de \all{sein} avec inversion (verbe en position 1, sans \all{wenn}) : \all{Wäre die Regierung früher kompromissbereit gewesen} (Si le gouvernement avait été plus tôt disposé au compromis).
    \item \all{hätten ... gerettet werden können} : KII passé du passif avec modal : auxiliaire \all{haben} (à cause du modal \all{können}) $\rightarrow$ \all{hätten viele Leben gerettet werden können} (de nombreuses vies auraient pu être sauvées).
    \item \all{müsste ... verhängen} : KII présent de \all{müssen} (conseil, nécessité atténuée) + infinitif \all{verhängen} en fin de phrase.
    \item \all{bleiben will} : après \all{wenn}, présent de \all{wollen} + infinitif : \all{wenn sie glaubwürdig bleiben will} (si elle veut rester crédible).
\end{enumerate}
\textbf{Texte complet :} \all{„Nachdem die Vermittler monatelang verhandelt hatten, scheiterten die Gespräche. Die Rebellen weigerten sich, die Waffen niederzulegen. Wäre die Regierung früher kompromissbereit gewesen, hätten viele Leben gerettet werden können. Heute müsste die internationale Gemeinschaft stärkere Sanktionen verhängen, wenn sie glaubwürdig bleiben will.“}
\end{corrige}

\courssub{Production guidée : argumenter sur l'ingérence (entraînement au Bac)}

\begin{textebox}[Analyse d'une situation fictive : « le cas de la République de N'Zassa »]
„In der Republik N'Zassa fanden Präsidentschaftswahlen statt. Monate vorher \textbf{hatte} die Regierung Oppositionsführer \textbf{verhaften lassen} und die Internetverbindung \textbf{unterbrochen}. Westliche Botschafter \textbf{erklärten} die Wahl später für „nicht glaubwürdig“. Die Regierung \textbf{reagierte} scharf: „Das ist Einmischung in innere Angelegenheiten!“

\textbf{Aufgabe :} Nehmen Sie Stellung. Schreiben Sie einen kurzen Kommentar (ca. 80--100 Wörter) auf Deutsch. Verwenden Sie unbedingt:
\begin{itemize}
    \item 1x \textbf{Plusquamperfekt} (antériorité d'un fait passé)
    \item 1x \textbf{Konjunktiv II Vergangenheit} (irréel du passé : ce qui aurait dû se passer)
    \item 1x \textbf{Konjunktiv II Gegenwart} (conseil ou hypothèse présente : \all{würde / könnte / müsste})
    \item 1x \textbf{Passiv} (Präteritum ou Perfekt)
\end{itemize}
\textbf{Wortschatz-Hilfe :} \all{die Glaubwürdigkeit} (la crédibilité) ; \all{unabhängig} (indépendant) ; \all{die Einmischung, -en} (l'ingérence) ; \all{verhängen} (imposer) ; \all{die Sanktion, -en} (la sanction) ; \all{respektieren} (respecter) ; \all{die Wahlkommission, -en} (la commission électorale).“
\end{textebox}

\begin{exemplebox}[Modèle de production attendue (corrigé type)]
„\textbf{Nachdem} die Regierung die Opposition \textbf{ausgeschaltet hatte}, \textbf{wurden} die Wahlen \textbf{manipuliert}. \textbf{Wäre} eine unabhängige Wahlkommission \textbf{zugelassen worden}, \textbf{hätte} das Volk frei \textbf{wählen können}. Heute \textbf{müsste} die internationale Gemeinschaft konsequenter reagieren: \textbf{würde} sie nur zuschauen, \textbf{würde} sie die Souveränität des Volkes verraten. Klare Sanktionen \textbf{wären} jetzt das richtige Signal.“
\tcblower
\textbf{Vérification des formes obligatoires :}
\begin{itemize}
    \item \all{hatte ... ausgeschaltet} : Plusquamperfekt actif après \all{nachdem} (\all{ausschalten} = éliminer, verbe à particule séparable).
    \item \all{wurden ... manipuliert} : passif au Präteritum.
    \item \all{Wäre ... zugelassen worden} : KII passé au passif (\all{wäre} + Partizip II + \all{worden}), avec inversion sans \all{wenn} ; \all{hätte ... wählen können} : KII passé avec modal.
    \item \all{müsste ... reagieren}, \all{würde ... zuschauen}, \all{würde ... verraten}, \all{wären} : KII présent.
\end{itemize}
\end{exemplebox}

\begin{aretenir}[Check-list finale : maîtriser le passé et l'irréel]
\begin{itemize}
    \item \cle{Plusquamperfekt} = \all{hatte/war + Partizip II}. \textbf{Signal :} \all{nachdem} + récit au passé. Choix de l'auxiliaire : mêmes règles qu'au Perfekt (\all{sein} = mouvement ou changement d'état).
    \item \cle{Konjunktiv II Vergangenheit} = \all{hätte/wäre + Partizip II}. \textbf{Sens :} « si j'avais fait, j'aurais fait ». \textbf{Structures :} \all{wenn} + KII passé $\rightarrow$ KII passé, \textbf{ou} inversion avec le verbe en position 1 (\all{Hätte..., wäre...}).
    \item \cle{Konjunktiv II Gegenwart} = \all{würde + infinitif}, ou formes propres : \all{hätte, wäre, könnte, müsste, dürfte, sollte, wollte}. \textbf{Sens :} hypothèse présente, souhait, conseil, politesse.
    \item \textbf{Passif au KII passé :} \all{wäre + Partizip II + worden} (ex. \all{wäre verhindert worden}).
    \item \textbf{Interdiction formelle :} \all{würde + Partizip II} n'existe pas en allemand standard.
\end{itemize}
\end{aretenir}

\coursec{Méthode du commentaire et commentaire modèle}

Dans la section précédente, nous avons étudié le plus-que-parfait et le subjonctif II : deux formes qui permettent d'exprimer l'hypothèse et le regret, par exemple \all{Wäre die Weltgemeinschaft untätig geblieben, hätten die Konflikte schlimmere Folgen gehabt.} „Si la communauté internationale était restée inactive, les conflits auraient eu des conséquences plus graves.“ Ces outils grammaticaux, ajoutés au passif et au lexique de la souveraineté, vont maintenant nous servir à rédiger l'exercice le plus attendu au Bac : le \cle{commentaire de texte} (\all{der Kommentar}). C'est l'aboutissement de la leçon : comprendre, analyser, puis prendre position par écrit, en allemand correct et structuré.

\courssub{Qu'est-ce qu'un commentaire de texte au Bac ?}

\begin{definition}[Der Kommentar]
Le \cle{commentaire} (\all{der Kommentar}) est une production écrite personnelle et argumentée à partir d'un texte support. Il ne s'agit ni de résumer le texte, ni de le traduire, mais de : présenter le texte, dégager la question qu'il pose (\all{die Problemstellung}), analyser les arguments, puis donner son propre avis (\all{die eigene Meinung}) de manière justifiée.
\end{definition}

Observez d'abord la différence entre un résumé et un commentaire :

\begin{center}
\begin{tabularx}{\linewidth}{|X|X|}\hline
\rowcolor{popblueL} \textbf{Résumé (Zusammenfassung)} & \textbf{Commentaire (Kommentar)} \\ \hline
On rapporte les idées de l'auteur, sans avis personnel. & On analyse et on prend position. \\ \hline
\all{Der Autor schreibt, dass Einmischung die Souveränität verletzt.} & \all{Meiner Meinung nach hat der Autor recht, denn die Souveränität ist die Basis der internationalen Ordnung.} \\ \hline
Temps : présent, style indirect. & Subjonctif II, passif, connecteurs d'opinion. \\ \hline
\end{tabularx}
\end{center}

\courssub{Les quatre étapes du commentaire}

\begin{methode}[Rédiger un commentaire en quatre étapes]
\begin{enumerate}
\item \textbf{Présenter le texte} (\all{Einleitung}) : indiquer le type de texte (\all{Zeitungsartikel, Rede, Interview}), le thème (\all{das Thema}) et la source si elle est connue. Phrase d'accroche modèle : \all{Der Text behandelt das Problem der Einmischung in die Politik anderer Staaten.} „Le texte traite le problème de l'ingérence dans la politique d'autres États.“
\item \textbf{Dégager la problématique} (\all{die Problemstellung}) : formuler la question centrale sous forme d'interrogation indirecte : \all{Es geht um die Frage, ob sich Staaten in die inneren Angelegenheiten anderer Länder einmischen dürfen.} „Il s'agit de la question de savoir si les États peuvent s'ingérer dans les affaires intérieures d'autres pays.“
\item \textbf{Analyser les arguments et le lexique} (\all{Hauptteil}) : présenter les arguments pour et contre avec les connecteurs \all{einerseits ... andererseits}, relever les mots-clés du champ lexical (\all{die Einmischung, die Macht, der Einfluss, die Souveränität, der Staat, das Land}).
\item \textbf{Donner son avis argumenté} (\all{Schluss}) : conclure avec \all{Meiner Meinung nach ...} suivi d'un argument personnel, éventuellement au subjonctif II pour nuancer.
\end{enumerate}
\end{methode}

\begin{popfigure}
\begin{tikzpicture}[
node distance=0.9cm,
etape/.style={rectangle, rounded corners, draw=popblue, fill=popblueL, thick, align=center, text width=6.2cm, minimum height=1cm},
fleche/.style={-{Stealth[length=3mm]}, thick, popdark}
]
\node[etape, align=center] (e1){\textbf{Einleitung}\\ \all{Der Text behandelt das Problem der Einmischung ...}};
\node[etape, below=of e1, fill=poppurpleL, draw=poppurple, align=center] (e2){\textbf{Hauptteil : Pro}\\ \all{Einerseits müssen Menschenrechte geschützt werden.}};
\node[etape, below=of e2, fill=poporangeL, draw=poporange, align=center] (e3){\textbf{Hauptteil : Contra}\\ \all{Andererseits wird die Souveränität oft verletzt.}};
\node[etape, below=of e3, fill=popgreenL, draw=popgreen, align=center] (e4){\textbf{Schluss : eigene Meinung}\\ \all{Meiner Meinung nach ist Hilfe nur legitim, wenn ...}};
\draw[fleche] (e1) -- (e2);
\draw[fleche] (e2) -- (e3);
\draw[fleche] (e3) -- (e4);
\end{tikzpicture}
\legende{Organigramme de la structure du commentaire : Einleitung, Hauptteil (Pro/Contra), Schluss.}
\end{popfigure}

\courssub{Les expressions de liaison (die Konnektoren)}

Un commentaire sans connecteurs est une liste de phrases ; avec des connecteurs, c'est un raisonnement. Voici la boîte à outils indispensable :

\begin{center}
\begin{tabularx}{\linewidth}{|l|X|X|}\hline
\rowcolor{popblueL} \textbf{Deutsch} & \textbf{Français} & \textbf{Exemple} \\ \hline
\all{zuerst} & d'abord, en premier lieu & \all{Zuerst muss man den Text vorstellen.} \\ \hline
\all{dann} & ensuite & \all{Dann analysiert man die Argumente.} \\ \hline
\all{außerdem} & de plus, en outre & \all{Außerdem spielt die Macht eine große Rolle.} \\ \hline
\all{jedoch} & cependant, toutefois & \all{Der Einfluss ist groß, jedoch nicht immer negativ.} \\ \hline
\all{einerseits ... andererseits} & d'une part ... d'autre part & \all{Einerseits schützt Hilfe Menschen, andererseits verletzt sie die Souveränität.} \\ \hline
\all{schließlich} & finalement, en fin de compte & \all{Schließlich muss jeder Staat selbst entscheiden.} \\ \hline
\all{meiner Meinung nach} & à mon avis & \all{Meiner Meinung nach ist Einmischung selten legitim.} \\ \hline
\end{tabularx}
\end{center}

\begin{propriete}[Place des connecteurs dans la phrase]
\begin{itemize}
\item \all{Zuerst, dann, außerdem, jedoch, schließlich, meiner Meinung nach} occupent la \cle{première place} de la phrase : le verbe conjugué vient \textbf{immédiatement après} (inversion). \all{Außerdem spielt die Macht eine Rolle.} (et non \all{Außerdem die Macht spielt...}).
\item \all{Einerseits} et \all{andererseits} introduisent chacun une phrase complète, souvent séparées par une virgule : \all{Einerseits wird die Souveränität verletzt, andererseits müssen Menschenrechte geschützt werden.}
\end{itemize}
\end{propriete}

\begin{attention}[Meiner Meinung nach : datif et place]
\begin{itemize}
\item \all{Meiner Meinung nach} : ici \all{nach} est \textbf{postposé} et exige le \cle{datif} : \all{mein\textbf{er} Meinung} (et non \all{meine Meinung nach}).
\item En copie soignée, on préfère \all{meiner Meinung nach} à \all{nach meiner Meinung}, qui est plus oral et relâché.
\item Placé en tête, il provoque l'inversion : \all{Meiner Meinung nach \textbf{sollte} die Souveränität jedes Staates respektiert \textbf{werden}.} „À mon avis, la souveraineté de chaque État devrait être respectée.“
\item Autres variantes utiles : \all{Ich bin der Meinung, dass ...} „Je suis d'avis que...“ ; \all{Meines Erachtens ...} „À mon sens...“.
\end{itemize}
\end{attention}

\begin{exemplebox}[Phrases modèles à réutiliser]
\begin{itemize}
\item \all{Der Text behandelt das Problem der Einmischung in die Politik anderer Länder.} „Le texte traite le problème de l'ingérence dans la politique d'autres pays.“
\item \all{Einerseits wird die Souveränität der Staaten oft verletzt.} „D'une part, la souveraineté des États est souvent violée.“ (passif)
\item \all{Andererseits müssen Menschenrechte überall geschützt werden.} „D'autre part, les droits de l'homme doivent être protégés partout.“ (passif + modal)
\item \all{Wäre die Weltgemeinschaft immer untätig geblieben, hätten viele Konflikte schlimmere Folgen gehabt.} „Si la communauté internationale était toujours restée inactive, de nombreux conflits auraient eu des conséquences plus graves.“ (subjonctif II passé)
\item \all{Meiner Meinung nach sollte die Souveränität jedes Staates respektiert werden.} „À mon avis, la souveraineté de chaque État devrait être respectée.“
\end{itemize}
\end{exemplebox}

\courssub{Commentaire modèle guidé}

Lisez le commentaire suivant, rédigé à partir du texte support de la leçon. Repérez : la structure en trois parties, les cinq mots du lexique thématique (\all{Einmischung, Souveränität, Staaten, Länder, Volk}), la phrase au passif et la phrase au subjonctif II.

\begin{textebox}[Kommentar : Modelltext]
Der Text behandelt die Frage, ob sich Staaten in die Politik anderer Länder einmischen dürfen. Zuerst stellt der Autor das Problem der Einmischung vor, dann zeigt er seine Folgen für kleine Länder wie viele afrikanische Länder.

Einerseits wird die Souveränität oft verletzt, wenn große Mächte ihren Einfluss missbrauchen. Andererseits müssen Menschenrechte geschützt werden, auch über die Grenzen hinweg. Außerdem hat die Geschichte gezeigt, dass Schweigen manchmal schlimmer ist als Handeln. Wäre die Weltgemeinschaft immer untätig geblieben, hätten viele Konflikte schlimmere Folgen gehabt.

Meiner Meinung nach ist Hilfe nur dann legitim, wenn sie vom Volk gewünscht wird. Schließlich gehört die Zukunft eines Landes seinen Bürgern und nicht fremden Staaten.
\end{textebox}

\textbf{Glossaire :} \all{sich einmischen in (+ Akk.)} „s'ingérer dans“ ; \all{die Folge, -n} „la conséquence“ ; \all{missbrauchen} „abuser de“ ; \all{über die Grenzen hinweg} „au-delà des frontières“ ; \all{das Schweigen} „le silence“ ; \all{untätig} „inactif“ ; \all{legitim} „légitime“ ; \all{der Bürger, -} „le citoyen“ (pl. \all{die Bürger}).

\begin{aretenir}[La checklist du commentaire réussi]
Avant de rendre votre copie, vérifiez :
\begin{enumerate}
\item Introduction avec \all{Der Text behandelt ...} (type, thème, problématique) ;
\item au moins deux arguments équilibrés avec \all{einerseits ... andererseits} ;
\item au moins cinq mots du lexique thématique ;
\item au moins une phrase au passif et une au subjonctif II ;
\item une conclusion personnelle avec \all{Meiner Meinung nach ...} ;
\item verbe conjugué en deuxième position, noms avec majuscule, articles corrects.
\end{enumerate}
\end{aretenir}

\courssub{Entraînement : corriger et compléter un commentaire}

\begin{exoresolu}[Un commentaire à améliorer]
Un élève a écrit : \all{Nach meiner Meinung die Einmischung ist immer schlecht. Die Staaten müssen nicht helfen.} Relevez trois fautes et réécrivez ce passage en commentaire correct de trois phrases, avec un connecteur, le passif et le subjonctif II.
\tcblower
\textbf{Fautes relevées :} 1) \all{nach meiner Meinung} est relâché et l'ordre des mots est faux (pas d'inversion) ; 2) absence totale de connecteurs ; 3) argumentation pauvre, aucune nuance, aucune forme étudiée réemployée.

\textbf{Version corrigée :} \all{Meiner Meinung nach ist die Einmischung nicht immer schlecht. Einerseits sollte die Souveränität jedes Staates respektiert werden, andererseits müssen Menschen in Not geschützt werden. Wäre niemand eingeschritten, hätte das Volk noch mehr gelitten.} „À mon avis, l'ingérence n'est pas toujours mauvaise. D'une part, la souveraineté de chaque État devrait être respectée, d'autre part, les personnes en détresse doivent être protégées. Si personne n'était intervenu, le peuple aurait encore plus souffert.“
\end{exoresolu}

\begin{exercice}[À vous de rédiger]
Rédigez un commentaire de 8 à 10 lignes sur le sujet suivant : \all{Sollen reiche Länder die Politik armer Länder beeinflussen?} „Les pays riches doivent-ils influencer la politique des pays pauvres ?“ Respectez la structure Einleitung -- Hauptteil (Pro/Contra) -- Schluss et la checklist de la boîte „À retenir“. Vous pouvez évoquer le contexte camerounais ou africain.
\end{exercice}

\courssub{Critères d'évaluation au Bac}

\begin{center}
\begin{tabularx}{\linewidth}{|l|X|}\hline
\rowcolor{popblueL} \textbf{Critère} & \textbf{Ce que le correcteur attend} \\ \hline
Correction grammaticale & verbe en 2\up{e} position, passif et subjonctif II corrects, cas et articles justes \\ \hline
Richesse du lexique & au moins 5 mots du champ thématique, connecteurs variés \\ \hline
Structure logique & Einleitung -- Hauptteil (Pro/Contra) -- Schluss clairement visibles \\ \hline
Pertinence de l'argumentation & avis personnel justifié, lié au texte, sans hors-sujet \\ \hline
\end{tabularx}
\end{center}

\begin{savaistu}[Le commentaire, l'exercice qui rapporte]
Au Bac allemand de Terminale A, le commentaire personnel bien structuré (Einleitung -- Hauptteil -- Schluss) rapporte des points même avec un allemand simple mais correct. Les correcteurs valorisent une argumentation claire avec des connecteurs bien placés plus qu'un allemand compliqué plein de fautes. Mieux vaut écrire \all{Meiner Meinung nach ist Hilfe wichtig, aber sie muss respektvoll sein.} qu'une longue phrase bancale. La clarté est une stratégie gagnante !
\end{savaistu}

\coursec{Production guidée et entraînement au Bac}

Après avoir analysé la structure du texte support et modélisé un commentaire composé dans la section précédente, nous passons maintenant à la phase active : \textbf{mobiliser le lexique, la grammaire (passif, plus-que-parfait, subjonctif II) et les stratégies de rédaction} pour réussir les épreuves types du Baccalauréat (résumé, questions de compréhension, thème, version, production écrite, expression orale). Chaque exercice est accompagné de sa méthode, d'exemples corrigés et de pièges à éviter.

\courssub{Résumé du texte (Zusammenfassung)}

\begin{methode}[Rédiger un résumé en allemand (die Zusammenfassung)]
    \begin{enumerate}
        \item \textbf{Lire} le texte une première fois pour saisir l'idée générale (\all{das Thema}).
        \item \textbf{Découper} le texte en séquences logiques (un paragraphe = une idée principale).
        \item \textbf{Rédiger} une phrase par paragraphe \textbf{au présent} (Präsens), \textbf{sans citation} directe, \textbf{avec ses propres mots}.
        \item \textbf{Relier} les phrases par des connecteurs logiques : \all{zuerst} (d'abord), \all{dann} / \all{anschließend} (ensuite), \all{außerdem} (de plus), \all{schließlich} / \all{zusammenfassend} (enfin / en résumé).
        \item \textbf{Vérifier} : 3 à 4 phrases maximum, pas de détails superflus, pas d'opinion personnelle, vocabulaire du thème réemployé.
    \end{enumerate}
\end{methode}

\begin{exemplebox}[Modèle de résumé (Zusammenfassung) pour le texte support de la leçon]
    \all{Zuerst wird das Thema der politischen Einmischung in die inneren Angelegenheiten souveräner Staaten vorgestellt. Dann wird erklärt, dass solche Eingriffe das Völkerrecht verletzen und die Souveränität untergraben. Anschließend wird am Beispiel wirtschaftlicher Sanktionen gezeigt, wie Druck ausgeübt wird. Schließlich fordert der Text die Achtung der Unabhängigkeit und warnt vor den Folgen unerlaubter Interventionen.}
    \par\smallskip
    \textit{Traduction : D'abord, le thème de l'ingérence politique dans les affaires intérieures d'États souverains est présenté. Ensuite, il est expliqué que de telles interventions violent le droit international et sapent la souveraineté. Par la suite, on montre par l'exemple des sanctions économiques comment la pression est exercée. Enfin, le texte exige le respect de l'indépendance et met en garde contre les conséquences d'interventions illégales.}
\end{exemplebox}

\begin{attention}[Pièges fréquents au résumé]
    \begin{itemize}
        \item \textbf{Temps :} En français, le résumé est souvent au passé. En allemand, \textbf{le résumé se rédige au Präsens} (présent de narration).
        \item \textbf{Citations :} Ne recopiez jamais des morceaux de phrases du texte. Reformulez : \all{die Einmischung} au lieu de reprendre mot pour mot \all{die Einmischung ausländischer Mächte}.
        \item \textbf{Connecteurs :} N'oubliez pas \all{zuerst, dann, schließlich}. Ils structurent le résumé et rapportent des points au Bac.
        \item \textbf{Passif au résumé :} Le passif présent est idéal pour un résumé neutre : \all{es wird erklärt, dass...} (il est expliqué que...), \all{es wird gezeigt, dass...} (il est montré que...). Attention au verbe en fin de subordonnée : \all{..., dass Druck ausgeübt wird.}
    \end{itemize}
\end{attention}

\begin{exoresolu}[Entraînement : résumé]
Énoncé : Rédigez en allemand un résumé de 3 à 4 phrases du texte suivant (extrait du support de la leçon 26).
\tcblower
\textbf{Texte source :}
„Die Einmischung in die Politik eines anderen Landes ist ein sensibles Thema. Historisch gesehen haben Großmächte oft ihre Interessen durchgesetzt, ohne die Souveränität der betroffenen Staaten zu achten. Heute verbietet die UN-Charta die Androhung oder Anwendung von Gewalt. Dennoch gibt es verdeckte Einmischung durch Cyberangriffe, Desinformation oder wirtschaftliche Erpressung. Wenn die internationale Gemeinschaft nicht reagiert, leidet die Glaubwürdigkeit des Völkerrechts.“

\textbf{Solution détaillée :}
\all{Zuerst wird die Einmischung in die Politik anderer Länder als sensibles Thema bezeichnet. Dann wird darauf hingewiesen, dass Großmächte historisch oft die Souveränität missachtet haben. Anschließend wird beschrieben, dass trotz der UN-Charta verdeckte Methoden wie Cyberangriffe oder Desinformation genutzt werden. Schließlich warnt der Text vor einem Verlust der Glaubwürdigkeit des Völkerrechts bei fehlender Reaktion.}

\textbf{Analyse :} 4 phrases = 4 idées principales. Présent (et Perfekt dans la subordonnée pour l'antériorité : \all{missachtet haben}). Connecteurs : \all{Zuerst, Dann, Anschließend, Schließlich}. Vocabulaire réemployé : \all{Souveränität, UN-Charta, verdeckte Methoden, Völkerrecht}.
\end{exoresolu}

\courssub{Réponses aux questions de compréhension (Leseverstehen)}

\begin{methode}[Répondre aux questions de compréhension]
    \begin{enumerate}
        \item \textbf{Lire} la question et souligner le mot interrogatif (\all{Wer, Was, Warum, Wie, Welche Folgen}).
        \item \textbf{Repérer} l'information dans le texte (lecture ciblée, \emph{scanning}).
        \item \textbf{Formuler} la réponse \textbf{en phrase complète}, \textbf{au présent} (sauf si la question porte sur le passé), \textbf{avec ses propres mots} (pas de copier-coller).
        \item \textbf{Citer} la ligne si demandé : \all{In Zeile 5 heißt es: „...“} (à la ligne 5, il est dit : « ... »).
    \end{enumerate}
\end{methode}

\begin{exoresolu}[Questions types Bac sur le texte support (ingérences)]
Énoncé : Répondez en allemand, en phrases complètes, aux questions suivantes.
\tcblower
\textbf{Frage 1 :} Was versteht man unter „verdeckter Einmischung“? Nennen Sie zwei Beispiele aus dem Text.

\textbf{Antwort 1 :} \all{Unter verdeckter Einmischung versteht man eine geheime Einflussnahme, die nicht offen als militärische Aktion erkennbar ist. Beispiele sind Cyberangriffe und Desinformation.}

\textit{Traduction : Par ingérence cachée, on entend une prise d'influence secrète, qui n'est pas reconnaissable ouvertement comme une action militaire. Les exemples sont les cyberattaques et la désinformation.}

\textbf{Frage 2 :} Warum ist die Souveränität der Staaten nach dem Text gefährdet?

\textbf{Antwort 2 :} \all{Die Souveränität ist gefährdet, weil Großmächte ihre Interessen durchsetzen, ohne die Unabhängigkeit anderer Staaten zu respektieren, und weil moderne Methoden wie wirtschaftliche Erpressung schwer zu kontrollieren sind.}

\textit{Traduction : La souveraineté est menacée parce que les grandes puissances imposent leurs intérêts sans respecter l'indépendance des autres États, et parce que les méthodes modernes comme le chantage économique sont difficiles à contrôler.}

\textbf{Frage 3 :} Welche Rolle spielt die UN-Charta laut Text?

\textbf{Antwort 3 :} \all{Die UN-Charta verbietet die Androhung oder Anwendung von Gewalt und schützt damit theoretisch die Souveränität der Staaten.}

\textit{Traduction : La Charte de l'ONU interdit la menace ou l'emploi de la force et protège ainsi théoriquement la souveraineté des États.}

\textbf{Frage 4 :} Welche Konsequenz zieht der Autor, wenn die Weltgemeinschaft nicht reagiert?

\textbf{Antwort 4 :} \all{Wenn die internationale Gemeinschaft nicht reagiert, leidet die Glaubwürdigkeit des Völkerrechts.}

\textit{Traduction : Si la communauté internationale ne réagit pas, la crédibilité du droit international en souffre.}
\end{exoresolu}

\courssub{Thème : traduction Français $\rightarrow$ Allemand}

\begin{methode}[Réussir le thème (Übersetzung ins Deutsche)]
    \begin{enumerate}
        \item \textbf{Analyser} la phrase française : identifier le \textbf{sujet}, le \textbf{verbe} (temps, mode, voix), les \textbf{compléments}.
        \item \textbf{Appliquer l'ordre des mots allemand} : verbe conjugué en 2\up{e} position (principale), verbe conjugué à la fin (subordonnée).
        \item \textbf{Traiter la grammaire ciblée} :
            \begin{itemize}
                \item \textbf{Passif :} \all{werden} + Partizip II (Présent : \all{wird verletzt} ; Präteritum : \all{wurde verletzt} ; Perfekt : \all{ist verletzt worden}).
                \item \textbf{Plus-que-parfait :} \all{hatte} / \all{war} + Partizip II (\all{hatte verletzt, war eingedrungen}).
                \item \textbf{Subjonctif II (irréel présent/passé) :} \all{wäre} / \all{hätte} / \all{würde} + Infinitif (\all{wäre eingetreten, hätte verhindert, würde sich einmischen}).
            \end{itemize}
        \item \textbf{Vérifier} : majuscules aux noms, Umlaute (ä, ö, ü), ß, guillemets „ “ si dialogue.
    \end{enumerate}
\end{methode}

\begin{exemplebox}[Exemples corrigés (thème) — passif, plus-que-parfait, subjonctif II]
    \begin{tabularx}{\linewidth}{|X|X|}
        \hline
        \rowcolor{popblueL} \textbf{Français} & \textbf{Deutsch} \\
        \hline
        La souveraineté de ce pays \textbf{a été violée}. (passif Präteritum) & \all{Die Souveränität dieses Landes \textbf{wurde verletzt}.} \\
        \hline
        \textbf{Si} la communauté internationale \textbf{était intervenue} plus tôt, \textbf{on aurait pu éviter} la guerre. (irréel passé) & \all{\textbf{Wäre} die Weltgemeinschaft früher \textbf{eingeschritten}, \textbf{hätte man den Krieg verhindern können}.} \\
        \hline
        Le gouvernement \textbf{avait déjà pris} des mesures avant l'arrivée des observateurs. (plus-que-parfait) & \all{Die Regierung \textbf{hatte} bereits Maßnahmen \textbf{ergriffen}, bevor die Beobachter eintrafen.} \\
        \hline
        La démocratie \textbf{est menacée} par l'influence étrangère. (passif présent) & \all{Die Demokratie \textbf{wird} durch ausländischen Einfluss \textbf{bedroht}.} \\
        \hline
    \end{tabularx}
\end{exemplebox}

\begin{attention}[Erreurs typiques des francophones au thème]
    \begin{itemize}
        \item \textbf{Ordre des mots (subordonnée) :} \all{..., dass die Demokratie bedroht wird.} (verbe conjugué à la fin !) — Faux : \all{..., dass die Demokratie wird bedroht.}
        \item \textbf{Passif :} Ne pas confondre \all{werden} (passif d'action) et \all{sein} (passif d'état). \all{Die Tür wird geöffnet} (on ouvre la porte, action) vs \all{Die Tür ist geöffnet} (la porte est ouverte, état).
        \item \textbf{Subjonctif II :} Ne pas mettre l'indicatif dans la subordonnée en \all{wenn} exprimant l'irréel. Faux : \all{Wenn die Weltgemeinschaft früher eingriff...} Juste : \all{Wäre die Weltgemeinschaft früher eingeschritten, ...} (inversion, sans \all{wenn}) ou \all{Wenn die Weltgemeinschaft früher eingeschritten wäre, ...}.
        \item \textbf{Plus-que-parfait :} Ne pas oublier l'auxiliaire \all{hatte/war}. Faux : \all{Die Regierung bereits Maßnahmen ergriffen.} Juste : \all{Die Regierung \textbf{hatte} bereits Maßnahmen ergriffen.}
        \item \textbf{Verbe réfléchi :} \all{sich einmischen} (s'ingérer) garde son pronom réfléchi : \all{Der Staat \textbf{mischt sich} in die Politik ein.} (particule \all{ein-} séparée et renvoyée en fin de phrase).
    \end{itemize}
\end{attention}

\begin{exoresolu}[Thème : 4 phrases à traduire]
Énoncé : Traduisez en allemand en mobilisant le passif, le plus-que-parfait et le subjonctif II.
\tcblower
\textbf{1. Les droits de l'homme ont été ignorés par le régime.} (passif Präteritum)

\all{Die Menschenrechte \textbf{wurden} vom Regime \textbf{ignoriert}.}

\textbf{2. Si les Nations Unies avaient agi plus vite, la catastrophe n'aurait pas eu lieu.} (irréel passé)

\all{\textbf{Hätten} die Vereinten Nationen schneller \textbf{gehandelt}, \textbf{wäre} die Katastrophe \textbf{nicht eingetreten}.}

\textit{Remarque : inversion sujet-auxiliaire dans la subordonnée sans \all{wenn} ; \all{eintreten} (avoir lieu, survenir) se conjugue avec \all{sein}.}

\textbf{3. Les observateurs avaient déjà signalé les fraudes avant les élections.} (plus-que-parfait)

\all{Die Beobachter \textbf{hatten} die Betrügereien bereits vor den Wahlen \textbf{gemeldet}.}

\textbf{4. On s'attend à ce que le gouvernement respecte la souveraineté.} (tournure passive impersonnelle)

\all{Es wird erwartet, dass die Regierung die Souveränität \textbf{respektiert}.}

\textit{Variante active impersonnelle :} \all{Man erwartet, dass die Regierung die Souveränität respektiert.}
\end{exoresolu}

\courssub{Version : traduction Allemand $\rightarrow$ Français}

\begin{textebox}[Extrait du texte support (version — environ 5 lignes)]
    \all{Die Einmischung ausländischer Mächte in die inneren Angelegenheiten eines souveränen Staates verletzte das Völkerrecht. Hätte die internationale Gemeinschaft früher reagiert, wäre der Konflikt vielleicht vermeidbar gewesen. Die Souveränität wurde durch wirtschaftliche Sanktionen und politische Druckmittel systematisch untergraben. Man muss die Unabhängigkeit jedes Landes respektieren, um Frieden zu sichern.}
    \par\medskip
    \textbf{Glossaire :} \all{die Einmischung, -en} (l'ingérence) ; \all{die ausländische Macht, -¨e} (la puissance étrangère) ; \all{die inneren Angelegenheiten} (pl.) (les affaires intérieures) ; \all{der souveräne Staat, -en} (l'État souverain) ; \all{verletzen, verletzte, hat verletzt} (violer) ; \all{das Völkerrecht} (le droit international) ; \all{reagieren} (réagir) ; \all{vermeidbar} (évitable) ; \all{die Sanktion, -en} (la sanction) ; \all{das Druckmittel, -} (le moyen de pression) ; \all{systematisch} (systématiquement) ; \all{untergraben, untergrub, hat untergraben} (saper, miner) ; \all{die Unabhängigkeit} (l'indépendance) ; \all{respektieren} (respecter) ; \all{sichern} (assurer, garantir).
\end{textebox}

\begin{attention}[Méthode version : repérer les temps AVANT de traduire]
    \begin{enumerate}
        \item \all{verletzte} $\rightarrow$ \textbf{Präteritum actif} : « violait / a violé ».
        \item \all{Hätte ... reagiert} $\rightarrow$ \textbf{Konjunktiv II Plusquamperfekt} (irréel du passé) : « Si ... avait réagi ».
        \item \all{wäre ... gewesen} $\rightarrow$ \textbf{Konjunktiv II Plusquamperfekt} (auxiliaire \all{sein}) : « aurait été ».
        \item \all{wurde ... untergraben} $\rightarrow$ \textbf{passif Präteritum} (\all{werden} + Partizip II) : « a été sapée / était sapée ».
        \item \all{muss ... respektieren} $\rightarrow$ \textbf{présent + modal + infinitif} : « doit respecter ».
    \end{enumerate}
    \textbf{Ne traduisez jamais mot à mot sans avoir identifié la structure verbale !}
\end{attention}

\begin{exoresolu}[Version corrigée]
Énoncé : Traduisez en français l'extrait ci-dessus.
\tcblower
\textbf{Traduction :}
« L'ingérence de puissances étrangères dans les affaires intérieures d'un État souverain \textbf{violait / a violé} le droit international. \textbf{Si} la communauté internationale \textbf{avait réagi} plus tôt, le conflit \textbf{aurait peut-être pu être évité}. La souveraineté \textbf{a été / était} systématiquement \textbf{sapée} par des sanctions économiques et des moyens de pression politiques. \textbf{Il faut / On doit} respecter l'indépendance de chaque pays pour \textbf{assurer} la paix. »

\textbf{Notes de traduction :}
\begin{itemize}
    \item \all{verletzte} : Präteritum de narration $\rightarrow$ passé simple ou passé composé selon le contexte.
    \item \all{Hätte ... reagiert / wäre ... gewesen} : structure \all{hätte/wäre + Partizip II} = irréel du passé (si + plus-que-parfait $\rightarrow$ conditionnel passé).
    \item \all{wurde ... untergraben} : passif d'action au Präteritum (\all{werden} conjugué + Partizip II).
    \item \all{Man muss} : impersonnel $\rightarrow$ « Il faut / On doit ».
    \item \all{um Frieden zu sichern} : \all{um ... zu} + infinitif = « pour + infinitif » (but).
\end{itemize}
\end{exoresolu}

\courssub{Production écrite : essai argumenté (Schreibaufgabe)}

\begin{propriete}[Règle d'or : la place du verbe]
    En allemand, \textbf{la position du verbe structure la phrase} :
    \begin{itemize}
        \item \textbf{Phrase principale (Hauptsatz) :} verbe conjugué en \textbf{2\up{e} position} (le sujet ou un complément occupe la 1\up{re} position). Ex. : \all{\textbf{Heute verletzt} die Einmischung das Recht.}
        \item \textbf{Phrase subordonnée (Nebensatz) :} verbe conjugué en \textbf{position finale}. Ex. : \all{..., \textbf{weil} die Einmischung das Recht \textbf{verletzt}.}
        \item \textbf{Infinitif, Partizip II, particule séparable :} en fin de phrase. Ex. : \all{Die Einmischung \textbf{muss} respektiert \textbf{werden}.} / \all{..., \textbf{dass} die Einmischung respektiert \textbf{werden muss}.}
    \end{itemize}
    \textbf{Toute phrase de votre production écrite doit respecter cette règle.}
\end{propriete}

\begin{methode}[Plan pour l'essai : „Soll sich ein Staat in die Politik eines anderen Landes einmischen?“]
    \textbf{Introduction (3-4 lignes) :}
    \begin{enumerate}
        \item \textbf{Accroche :} actualité ou définition (\all{Die Frage der Einmischung ist aktueller denn je...}).
        \item \textbf{Définition du sujet :} \all{Einmischung bedeutet den Eingriff in die Souveränität eines anderen Staates.}
        \item \textbf{Annonce du plan :} \all{Ich werde Argumente für und gegen eine Intervention prüfen.}
    \end{enumerate}
    \textbf{Développement (6-8 lignes) :}
    \begin{enumerate}
        \item \textbf{Argument POUR (la responsabilité de protéger) :} \all{Einerseits: Schutz der Menschenrechte, Verhinderung von Völkermord.} Connecteurs : \all{Zunächst, Ein starkes Argument ist...}
        \item \textbf{Argument CONTRE (souveraineté, abus) :} \all{Andererseits: Verletzung der Souveränität, Missbrauch als Vorwand für Machtinteressen, Instabilität.} Connecteurs : \all{Auf der anderen Seite, Allerdings, Kritiker warnen davor, dass...}
        \item \textbf{Synthèse / nuance :} \all{nur unter UN-Mandat, letzte Option, Diplomatie vor Militär.}
    \end{enumerate}
    \textbf{Conclusion (2-3 lignes) :}
    \begin{enumerate}
        \item \textbf{Bilan :} \all{Zusammenfassend lässt sich sagen...}
        \item \textbf{Opinion personnelle nuancée :} \all{Meiner Meinung nach sollte eine Einmischung nur durch die UN legitimiert und als letzte Option erfolgen.}
    \end{enumerate}
\end{methode}

\begin{exemplebox}[Modèle de production écrite (environ 11 lignes)]
    \all{Die Frage, ob sich ein Staat in die Politik eines anderen Landes einmischen soll, ist aktueller denn je. Einmischung bedeutet den Eingriff in die inneren Angelegenheiten eines souveränen Staates. Ich werde Argumente für und gegen eine solche Intervention prüfen.}
    \par\smallskip
    \all{Einerseits gibt es die Verantwortung, Menschen zu schützen. Wenn ein Regime Verbrechen gegen die Menschlichkeit begeht, hat die internationale Gemeinschaft die Pflicht einzugreifen, um Menschenleben zu retten. Ein starkes Argument ist die Verhinderung von Völkermord: Das Nichteingreifen in Ruanda 1994 hatte fatale Folgen.}
    \par\smallskip
    \all{Auf der anderen Seite verletzt jede Einmischung die Souveränität und das Völkerrecht. Kritiker warnen davor, dass Menschenrechte oft nur als Vorwand für geopolitische Machtinteressen dienen. Historisch gesehen haben Interventionen oft zu mehr Instabilität geführt, zum Beispiel im Irak 2003.}
    \par\smallskip
    \all{Zusammenfassend lässt sich sagen, dass eine Einmischung nur legitim ist, wenn sie durch ein UN-Mandat gedeckt ist und als letzte Option nach gescheiterter Diplomatie erfolgt. Meiner Meinung nach muss der Schutz der Souveränität Vorrang haben, solange keine Massenverbrechen drohen.}
    \par\smallskip
    \textit{Traduction : La question de savoir si un État doit s'ingérer dans la politique d'un autre pays est plus actuelle que jamais. L'ingérence signifie l'intervention dans les affaires intérieures d'un État souverain. J'examinerai les arguments pour et contre une telle intervention. D'un côté, il existe la responsabilité de protéger les hommes. Si un régime commet des crimes contre l'humanité, la communauté internationale a le devoir d'intervenir pour sauver des vies. Un argument fort est la prévention du génocide : la non-intervention au Rwanda en 1994 a eu des conséquences fatales. De l'autre côté, toute ingérence viole la souveraineté et le droit international. Les critiques avertissent que les droits de l'homme servent souvent de prétexte à des intérêts de puissance géopolitiques. Historiquement, les interventions ont souvent mené à plus d'instabilité, par exemple en Irak en 2003. En résumé, on peut dire qu'une ingérence n'est légitime que si elle est couverte par un mandat de l'ONU et intervient en dernier recours après l'échec de la diplomatie. À mon avis, la protection de la souveraineté doit avoir la priorité tant qu'aucun crime de masse ne menace.}
\end{exemplebox}

\begin{exercice}[Sujet d'entraînement Bac — production écrite]
    \textbf{Consigne :} Traitez le sujet suivant en 10 à 12 lignes (environ 100-120 mots). Respectez le plan ci-dessus. Utilisez le vocabulaire de la leçon (\cle{die Einmischung, die Souveränität, das Völkerrecht, die Menschenrechte, die Sanktionen, das UN-Mandat}) et au moins \textbf{un passif, un plus-que-parfait et un subjonctif II}.
    \par\medskip
    \textbf{Thema :} \all{Soll sich ein Staat in die Politik eines anderen Landes einmischen?}
\end{exercice}

\courssub{Expression orale : débat en binômes (Pro / Contra)}

\begin{experience}[Débat structuré : „Einmischung — Ja oder Nein?“]
    \textbf{Organisation (15 min) :}
    \begin{enumerate}
        \item \textbf{Préparation (5 min) :} en binômes. Chaque binôme tire un rôle : \textbf{groupe A (pour l'ingérence)} ou \textbf{groupe B (contre l'ingérence)}. Préparer 3 arguments solides avec le lexique de la leçon. Noter des mots-clés, pas des phrases toutes faites.
        \item \textbf{Confrontation (5 min) :} binôme A contre binôme B. Tour de parole : 1 minute par argument. L'écouteur note les arguments adverses pour la réfutation.
        \item \textbf{Synthèse orale (5 min) :} 2-3 binômes volontaires présentent leur débat devant la classe. Le professeur note au tableau les arguments clés dans un tableau Pro/Contra.
    \end{enumerate}
    \textbf{Aide lexicale (Sprachhilfe) :}
    \begin{center}
    \begin{tabularx}{\linewidth}{|X|X|}
        \hline
        \rowcolor{popgreenL} \textbf{Pro (pour l'ingérence)} & \textbf{Contra (contre l'ingérence)} \\
        \hline
        \all{der Schutz der Menschenrechte} & \all{die Verletzung der Souveränität} \\
        \all{die Verhinderung von Völkermord} & \all{das Nichteinmischungsgebot (Völkerrecht)} \\
        \all{die Schutzverantwortung} & \all{der Missbrauch als Vorwand (Machtinteressen)} \\
        \all{die Legitimation durch ein UN-Mandat} & \all{das Risiko der Destabilisierung / des Bürgerkriegs} \\
        \all{die letzte Option} & \all{Diplomatie / Sanktionen wirksamer} \\
        \hline
    \end{tabularx}
    \end{center}
    \textbf{Phrases types (Redemittel) :}
    \all{Meiner Meinung nach muss man eingreifen, wenn...} (à mon avis, il faut intervenir quand...) \\
    \all{Das stimmt, aber man darf nicht vergessen, dass...} (c'est vrai, mais il ne faut pas oublier que...) \\
    \all{Ein Gegenargument ist, dass...} (un contre-argument est que...) \\
    \all{Ich sehe das anders, weil...} (je vois cela autrement, parce que...)
\end{experience}

\begin{savaistu}[L'Allemagne et le débat sur les interventions extérieures]
    L'Allemagne débat régulièrement au \textbf{Bundestag} de ses interventions militaires à l'étranger (\all{die Auslandseinsätze der Bundeswehr}). Depuis les années 1990 (Kosovo 1999, Afghanistan à partir de 2001, Mali), la question de l'\all{Einmischung} divise : le pacifisme hérité de l'histoire s'oppose à la \all{Verantwortung} (responsabilité) de protéger les droits de l'homme. Tout déploiement armé nécessite un mandat du Parlement (\all{ein Mandat des Bundestages}), ce qui illustre le contrôle démocratique de la politique étrangère (\all{die Außenpolitik}). C'est un exemple réel et vivant du dilemme « souveraineté contre droits de l'homme » traité dans cette leçon.
\end{savaistu}

\courssub{Auto-évaluation et grille de révision (Bac)}

\begin{center}
\begin{tabularx}{\linewidth}{|l|X|}
    \hline
    \rowcolor{popblueL} \textbf{Critère} & \textbf{Indicateurs de réussite (à cocher)} \\
    \hline
    \textbf{Lexique thématique} /5 & \all{Einmischung, Souveränität, Macht, Einfluss, Staat, Land, Völkerrecht, Sanktionen, Mandat} : articles et pluriels corrects. \\
    \hline
    \textbf{Grammaire} /5 & Passif (\all{werden} + Partizip II) aux trois temps ; plus-que-parfait (\all{hatte/war} + Partizip II) ; subjonctif II (\all{würde/hätte/wäre}) pour l'irréel. Formes verbales exactes. \\
    \hline
    \textbf{Structure et cohérence} /5 & Plan respecté (introduction / développement / conclusion) ; connecteurs (\all{zuerst, einerseits/andererseits, schließlich}) ; place du verbe (2\up{e} position / finale) respectée partout. \\
    \hline
    \textbf{Argumentation et contenu} /5 & Arguments pour et contre équilibrés ; exemples concrets (ONU, histoire) ; opinion personnelle nuancée en conclusion ; registre soutenu. \\
    \hline
    \rowcolor{popgoldL} \textbf{TOTAL /20} & 0-5 insuffisant ; 6-10 fragile ; 11-14 assez bien ; 15-17 bien ; 18-20 très bien. \\
    \hline
\end{tabularx}
\end{center}

\begin{aretenir}[Check-list finale avant l'épreuve]
    \begin{itemize}
        \item \textbf{Vocabulaire :} je connais le genre, le pluriel et un synonyme pour \cle{die Einmischung, die Souveränität, der Einfluss, die Macht, der Staat, das Land}.
        \item \textbf{Passif :} je sais former le \textbf{présent} (\all{wird gemacht}), le \textbf{Präteritum} (\all{wurde gemacht}), le \textbf{Perfekt} (\all{ist gemacht worden}) et la forme avec modal (\all{muss gemacht werden}).
        \item \textbf{Plus-que-parfait :} \all{hatte gemacht / war gegangen} (exprimer l'antériorité dans le passé).
        \item \textbf{Subjonctif II :} \all{würde machen / hätte gemacht / wäre gegangen} (irréel, politesse). Phrase en « si » : \all{Wenn ich Zeit hätte, käme ich.} ou \all{Hätte ich Zeit, käme ich.}
        \item \textbf{Ordre des mots :} \textbf{principale = verbe en 2\up{e} position ; subordonnée = verbe à la fin} (même avec \all{weil, dass, wenn, obwohl}).
        \item \textbf{Rédaction :} présent pour le résumé ; plan introduction / pour / contre / conclusion pour l'essai ; connecteurs logiques obligatoires.
    \end{itemize}
\end{aretenir}

\newpage

\coursec{Activité d'intégration}

\courssub{Objectifs de l'activité}

À l'issue de cette activité d'intégration, l'élève sera capable de :

\begin{itemize}
\item réinvestir le \cle{lexique thématique} de la leçon (\all{die Einmischung}, \all{die Souveränität}, \all{der Einfluss}, \all{die Macht}, \all{der Staat}, \all{das Völkerrecht}) dans un échange oral argumenté ;
\item employer correctement le \cle{passif} (\all{Die Menschenrechte werden verletzt.}), le \cle{plus-que-parfait} (\all{Die Regierung hatte die Proteste niedergeschlagen.}) et le \cle{subjonctif II} (\all{Man sollte eingreifen.}) dans des prises de parole spontanées et préparées ;
\item construire une \cle{argumentation} en allemand : thèse, arguments, exemples, concession, conclusion ;
\item rédiger un court \cle{article de presse} équilibré rendant compte d'un débat et exprimant une opinion personnelle ;
\item développer des compétences de citoyenneté : écouter, répondre, respecter la parole de l'autre, défendre une position avec des arguments et non avec des slogans.
\end{itemize}

\courssub{Situation}

Votre lycée de Yaoundé participe à la « Semaine de la citoyenneté » organisée en partenariat avec le club d'allemand. Le thème de cette année est : « Les États peuvent-ils se mêler des affaires des autres ? ». Pour clôturer la semaine, la classe de Terminale A organise un débat parlementaire en langue allemande, devant les élèves de Première et le professeur principal.

Voici le dossier de presse distribué à tous les participants. Il présente un pays fictif, \emph{Korania}, où la situation des droits de l'homme inquiète la communauté internationale.

\begin{textebox}[Dossier de presse : Krise in Korania]
Seit zwei Jahren gibt es in dem kleinen Staat Korania schwere politische Spannungen. Nach den umstrittenen Wahlen wurden Demonstrationen von der Armee gewaltsam niedergeschlagen. Viele Journalisten wurden verhaftet, mehrere Zeitungen wurden geschlossen. Die Opposition behauptet, dass Tausende Menschen ohne Prozess ins Gefängnis gebracht worden seien.

Die Nachbarländer und die Vereinten Nationen diskutieren jetzt über eine mögliche Intervention. Die einen sagen: „Die Menschenrechte werden verletzt, wir müssen handeln!“ Die anderen warnen: „Die Souveränität eines Staates muss respektiert werden. Eine Einmischung von außen würde die Lage nur verschlimmern.“

Die Regierung von Korania erklärt: „Kein fremder Staat hat das Recht, sich in unsere inneren Angelegenheiten einzumischen. Wir lösen unsere Probleme selbst.“

\emph{Glossaire :} umstritten = contesté ; niederschlagen = réprimer ; verhaften = arrêter ; das Gefängnis, -se = la prison ; die Lage = la situation ; sich einmischen in (+ Akk.) = s'immiscer dans ; die innere Angelegenheit, -en = l'affaire intérieure.
\end{textebox}

La classe est divisée en deux camps :

\begin{itemize}
\item le \textbf{Bundestag} : groupe favorable à une intervention internationale (sanctions, médiation, aide humanitaire) ;
\item le \textbf{Conseil de souveraineté} (\all{der Souveränitätsrat}) : groupe opposé à toute ingérence extérieure.
\end{itemize}

\courssub{Consignes}

\begin{methode}[Déroulement de l'activité d'intégration]
\textbf{Tâche 1 — Préparation écrite individuelle (15 min).} Chaque élève rédige \textbf{trois arguments} pour son camp. Contraintes obligatoires :
\begin{enumerate}
\item au moins \textbf{une phrase au passif} (ex. \all{Die Pressefreiheit wurde abgeschafft.}) ;
\item au moins \textbf{une phrase au subjonctif II} (ex. \all{Man sollte die Flüchtlinge aufnehmen.}) ;
\item au moins \textbf{une phrase au plus-que-parfait} pour rappeler les faits antérieurs (ex. \all{Die Regierung hatte die Wahlen manipuliert.}) ;
\item chaque argument utilise au moins un mot du lexique de la leçon.
\end{enumerate}

\textbf{Tâche 2 — Le débat oral (10 min).} Le professeur joue le rôle du président de séance (\all{der Sitzungspräsident}). Les deux camps alternent : un orateur du Bundestag, puis un orateur du Conseil de souveraineté. Chaque prise de parole dure au maximum une minute. Règles du débat :
\begin{enumerate}
\item commencer par une formule de politesse : \all{Sehr geehrte Damen und Herren, ...} ;
\item annoncer sa position : \all{Ich bin der Meinung, dass ...} / \all{Meiner Ansicht nach ...} ;
\item réagir à l'argument précédent : \all{Sie haben gesagt, dass ..., aber ...} / \all{Ich stimme Ihnen zu, dass ..., jedoch ...} ;
\item conclure : \all{Aus diesen Gründen bin ich für / gegen eine Einmischung.}
\end{enumerate}

\textbf{Tâche 3 — Le compte rendu écrit (20 min).} Après le débat, chaque élève rédige individuellement un court article de presse d'environ \textbf{dix lignes}, intitulé \all{„Einmischung oder Respekt der Souveränität?“}. L'article doit :
\begin{enumerate}
\item présenter brièvement la situation en Korania (2 lignes, avec un plus-que-parfait) ;
\item résumer les arguments des deux camps (4 lignes, avec au moins deux passifs) ;
\item donner son avis personnel justifié (3 lignes, avec un subjonctif II) ;
\item se terminer par une phrase de conclusion ou une question ouverte (1 ligne).
\end{enumerate}
\end{methode}

\textbf{Tâche 4 — Évaluation.} Le professeur évalue chaque production selon la grille suivante :

\begin{center}
\begin{tabularx}{\linewidth}{|X|c|}
\hline
\rowcolor{popblueL} \textbf{Critère} & \textbf{Points} \\
\hline
Lexique thématique (au moins 5 mots du champ lexical) & 5 \\
\hline
Grammaire (passif, plus-que-parfait, subjonctif II corrects) & 6 \\
\hline
Structure de l'article (titre, introduction, deux positions, avis, conclusion) & 4 \\
\hline
Argumentation (arguments clairs, exemples, connecteurs logiques) & 5 \\
\hline
\end{tabularx}
\end{center}

\courssub{Ressources à mobiliser}

\begin{aretenir}[Boîte à outils grammaticale]
\begin{itemize}
\item \textbf{Le passif (Vorgangspassiv)} : \all{werden} (conjugué) + participe II. \all{Die Zeitungen wurden geschlossen.} « Les journaux ont été fermés. » Agent introduit par \all{von} + datif : \all{Die Proteste wurden von der Armee niedergeschlagen.} Distinction : \all{sein} + participe II = passif d'état (\all{Die Zeitungen sind geschlossen} = « Les journaux sont fermés »).
\item \textbf{Le plus-que-parfait} : \all{hatte / war} + participe II. \all{Die Regierung hatte die Wahlen manipuliert.} « Le gouvernement avait manipulé les élections. » Il exprime l'antériorité par rapport à un moment du passé (Präteritum ou Perfekt).
\item \textbf{Le subjonctif II (Konjunktiv II)} : forme \all{würde} + infinitif (tous verbes) ou formes simples pour les verbes modaux et \all{sein/haben} : \all{sollte, könnte, müsste, wäre, hätte}. \all{Man sollte die Menschenrechte schützen.} « On devrait protéger les droits de l'homme. » Emplois : conseil, hypothèse, opinion nuancée, politesse. \textbf{Ordre des mots :} l'infinitif (ou \all{würde}) se place \textbf{en fin de proposition}.
\end{itemize}
\end{aretenir}

\begin{aretenir}[Lexique et connecteurs utiles]
\textbf{Noms (avec article et pluriel) :}
\begin{itemize}
\item \all{die Einmischung, -en} (l'ingérence) ; \all{die Souveränität} (la souveraineté, pas de pluriel usuel) ; \all{der Einfluss, -e} (l'influence) ; \all{die Macht, -en} (le pouvoir) ; \all{der Staat, -en} (l'État) ; \all{das Völkerrecht} (le droit international) ; \all{die Menschenrechte} (pl.) (les droits de l'homme).
\end{itemize}
\textbf{Verbes et locutions :}
\begin{itemize}
\item \all{sich einmischen in + Akk.} (s'immiscer dans) ; \all{eingreifen} (intervenir) ; \all{respektieren} (respecter) ; \all{verletzen} (violer) ; \all{niederschlagen} (réprimer) ; \all{Sanktionen beschließen} (décider des sanctions).
\end{itemize}
\textbf{Connecteurs logiques :}
\begin{center}
\begin{tabularx}{\linewidth}{|X|X|X|}
\hline
\rowcolor{popblueL} \textbf{Français} & \textbf{Deutsch} & \textbf{Niveau} \\
\hline
d'une part ... d'autre part & \all{einerseits ... andererseits} & B2 \\
\hline
de plus / en outre & \all{außerdem / ferner} & B1/B2 \\
\hline
cependant / toutefois & \all{jedoch / allerdings} & B1 \\
\hline
c'est pourquoi / par conséquent & \all{deshalb / daher / folglich} & B2 \\
\hline
à mon avis & \all{meiner Meinung nach / meines Erachtens} & B1 \\
\hline
au contraire & \all{im Gegenteil} & B1 \\
\hline
\end{tabularx}
\end{center}
\end{aretenir}

\begin{attention}[Pièges à éviter pour francophones]
\begin{itemize}
\item \textbf{Rection verbale :} \all{sich einmischen \textbf{in} + Akkusativ} (jamais \all{für}, \all{über}, \all{an}). Ex. \all{Er mischt sich in die Angelegenheit ein.}
\item \textbf{Auxiliaire du passif :} \all{werden} (processus), pas \all{sein} (état). \all{Die Journalisten \textbf{wurden} verhaftet} (action) vs \all{Die Journalisten \textbf{waren} verhaftet} (état résultant, rare).
\item \textbf{Place de l'infinitif (Subjonctif II / Modaux) :} \all{Man \textbf{sollte} die Flüchtlinge \textbf{aufnehmen.}} (verbe à la fin).
\item \textbf{Faux ami :} \all{die Intervention} existe (intervention médicale, militaire), mais en politique \all{die Einmischung} a une connotation péjorative (ingérence indue), \all{das Eingreifen} est neutre.
\item \textbf{Majuscules :} Tous les noms prennent une majuscule : \all{die Menschenrechte, die Souveränität, der Staat}.
\item \textbf{Passif au plus-que-parfait :} \all{waren verhaftet worden} (et non \all{hatten verhaftet worden}).
\end{itemize}
\end{attention}

\begin{savaistu}[Le Cameroun, l'Union africaine et le principe de non-ingérence]
L'article 4(h) de l'Acte constitutif de l'Union africaine (2000) prévoit le droit d'ingérence de l'Union dans un État membre en cas de crimes de guerre, de génocide ou de crimes contre l'humanité. C'est une évolution majeure par rapport à la charte de l'OUA (1963) qui sacralisait l'intangibilité des frontières et la non-ingérence absolue. Le Cameroun, en tant que membre actif de l'UA, participe à ces débats : la souveraineté n'est plus un bouclier pour l'impunité, mais une responsabilité de protéger (\emph{Responsibility to Protect} / \all{Schutzverantwortung}). En allemand, on parle de \all{Schutzpflicht} ou \all{Responsibility to Protect (R2P)}.
\end{savaistu}

\courssub{Proposition de production et corrigé commenté}

\begin{exoresolu}[Article modèle : compte rendu du débat]
Rédigez un article d'environ dix lignes intitulé \all{„Einmischung oder Respekt der Souveränität?“} selon les consignes de la tâche 3.
\tcblower
\textbf{Production modèle :}

\all{„Einmischung oder Respekt der Souveränität?“}

\all{In Korania waren nach den umstrittenen Wahlen Tausende Menschen verhaftet worden, und die Regierung hatte mehrere Zeitungen geschlossen. Darüber hat unsere Klasse eine lebhafte parlamentarische Debatte geführt.}

\all{Der „Bundestag“ war der Meinung, dass die internationale Gemeinschaft eingreifen müsse, weil die Menschenrechte verletzt werden. Die Demonstranten seien von der Armee gewaltsam niedergeschlagen worden; deshalb sollte man Sanktionen beschließen und den Flüchtlingen helfen.}

\all{Der „Souveränitätsrat“ war dagegen der Ansicht, dass die Souveränität jedes Staates respektiert werden muss. Eine Einmischung von außen würde die Lage nur verschlimmern; außerdem hatte Korania das Recht, seine inneren Probleme selbst zu lösen.}

\all{Meiner Meinung nach sollte man zuerst eine friedliche Vermittlung versuchen, bevor man militärisch eingreift. Die Menschenrechte sind wichtig, aber der Respekt der Souveränität darf nicht vergessen werden. Bleibt die Frage: Wo endet der Schutz der Menschen, und wo beginnt die Einmischung?}

\medskip
\textbf{Commentaire des choix de langue :}
\begin{itemize}
\item \textbf{Plus-que-parfait} : \all{hatte ... geschlossen} (actif) et \all{waren ... verhaftet worden} (passif : \all{sein} au plus-que-parfait + participe II + \all{worden}) situent les faits antérieurs au débat.
\item \textbf{Passifs variés} : \all{werden verletzt} (Präsens passif), \all{niedergeschlagen worden seien} (Perfekt passif au subjonctif I pour le discours rapporté), \all{muss respektiert werden} (passif modal avec \all{müssen}).
\item \textbf{Subjonctif II} : \all{sollte man Sanktionen beschließen} (conseil), \all{würde ... verschlimmern} (hypothèse), \all{sollte man ... versuchen} (opinion nuancée). Note : \all{müsse eingreifen} est un subjonctif I (discours indirect), toléré et attendu en Terminale.
\item \textbf{Lexique de la leçon} : \all{Einmischung, Souveränität, Staat, eingreifen, Menschenrechte, inneren Probleme} — six occurrences du champ lexical, critère pleinement validé.
\item \textbf{Structure} : titre, situation (2 lignes, plus-que-parfait), position A (passifs), position B (passif modal), avis personnel (subjonctif II), question finale ouverte : les quatre exigences de la tâche 3 sont respectées.
\end{itemize}
\end{exoresolu}

\courssub{Bilan de l'activité}

Cette activité d'intégration a permis de mobiliser simultanément les quatre savoir-faire de la leçon :

\begin{itemize}
\item \textbf{Compréhension} : lecture et exploitation du dossier de presse sur Korania (inférence, repérage d'arguments contraires) ;
\item \textbf{Lexique} : emploi actif, à l'oral et à l'écrit, du vocabulaire de l'ingérence et de la souveraineté (noms, verbes à préposition, collocations) ;
\item \textbf{Grammaire} : production contrôlée puis libre du passif (Présent, Prétérit, Modal, Plus-que-parfait), du plus-que-parfait (antériorité narrative) et du subjonctif II (modalité épistémique/déontique) dans un contexte de communication réel ;
\item \textbf{Production} : prise de parole argumentée (interaction orale) et rédaction d'un article de presse structuré (expression écrite), exercice type des épreuves du Baccalauréat (commentaire / rédaction guidée).
\end{itemize}

Sur le plan citoyen, le débat a montré qu'une question politique complexe n'appelle pas une réponse simpliste : défendre les droits de l'homme (\all{die Menschenrechte schützen}) et respecter la souveraineté des États (\all{die Souveränität respektieren}) sont deux exigences légitimes qu'il faut savoir articuler (\all{abwägen}). C'est exactement la compétence attendue d'un candidat de Terminale A : comprendre, argumenter, nuancer — en allemand.

Pour prolonger l'activité : les meilleurs articles pourront être publiés dans le journal du club d'allemand du lycée ; le débat pourra être rejoué en version filmée pour travailler la prononciation (\all{Aussprache}, \all{Satzmelodie}) et la prise de parole en public (\all{freies Sprechen}).

\newpage

\coursec{Méthodes et exercices résolus}

\coursec{AL26 : Les ingérences en politique}

\courssub{Texte support et première compréhension}

\begin{textebox}[Texte original — Contexte camerounais et germanophone]
\all{Nach der Kolonialzeit behielten viele ehemalige Kolonialmächte großen Einfluss auf die Politik ihrer früheren Kolonien. In Kamerun zum Beispiel prägten Frankreich und Großbritannien die staatlichen Strukturen bis heute: die Amtssprachen, das Rechtssystem und die Verwaltungsorganisation tragen ihre Handschrift. Doch auch nach der formellen Unabhängigkeit 1960 blieb die Einmischung Realität. Wirtschaftsverträge, an Bedingungen geknüpfte Entwicklungshilfe und militärische Kooperationen schränken oft die Souveränität der jungen Staaten ein.}

\all{Die Bürgerinnen und Bürger fordern zunehmend eine Politik ohne äußere Einmischung. „Wir wollen unsere eigenen Entscheidungen treffen“, sagt ein Student in Yaoundé. In Deutschland und der Schweiz beobachtet man die Debatte um „Einmischung“ ebenfalls kritisch: Wahlbeeinflussung durch Cyberangriffe, Desinformation in sozialen Medien und wirtschaftlicher Druck großer Konzerne gelten als moderne Formen der Machtausübung. Die Vereinten Nationen betonen in ihrer Charta das Prinzip der Nichteinmischung in die inneren Angelegenheiten souveräner Staaten — ein Prinzip, das im Alltag oft verletzt wird.}
\end{textebox}

\begin{definition}[Glossaire du texte (mots difficiles)]
\begin{itemize}
\item \all{die Kolonialzeit, -en} — l’époque coloniale
\item \all{die Kolonialmacht, -¨e} — la puissance coloniale
\item \all{der Einfluss, -¨e} — l’influence
\item \all{die Amtssprache, -n} — la langue officielle
\item \all{das Rechtssystem, -e} — le système juridique
\item \all{die Verwaltungsorganisation} — l’organisation administrative
\item \all{die Handschrift} — la marque, l’empreinte (fig.)
\item \all{die formelle Unabhängigkeit} — l’indépendance formelle
\item \all{die Einmischung, -en} — l’ingérence
\item \all{an Bedingungen geknüpft} — assorti de conditions (part. II adjectivé)
\item \all{die Entwicklungshilfe} — l’aide au développement
\item \all{die militärische Kooperation} — la coopération militaire
\item \all{einschränken} — limiter, restreindre (verbe à particule séparable : \all{schränkt ... ein})
\item \all{die Souveränität} — la souveraineté (gén. Sg.)
\item \all{die Bürgerin, -nen / der Bürger, -} — la citoyenne / le citoyen
\item \all{treffen (Entscheidungen)} — prendre (des décisions)
\item \all{die Wahlbeeinflussung} — l’ingérence électorale / l’influence sur les élections
\item \all{der Cyberangriff, -e} — la cyberattaque
\item \all{die Desinformation} — la désinformation
\item \all{der wirtschaftliche Druck} — la pression économique
\item \all{die Machtausübung} — l’exercice du pouvoir
\item \all{die Vereinten Nationen (VN)} — les Nations Unies (ONU)
\item \all{die Charta} — la Charte
\item \all{das Prinzip der Nichteinmischung} — le principe de non-ingérence
\item \all{verletzen} — violer (un principe, une loi)
\end{itemize}
\end{definition}

\courssub{Lexique thématique : la politique et la souveraineté}

\begin{propriete}[Mots-clés : articles, pluriels, familles et collocations]
\begin{center}
\begin{tabularx}{\linewidth}{|X|X|X|}
\hline
\rowcolor{popblueL} \textbf{Nom (article + pluriel)} & \textbf{Famille verbale / adjectivale} & \textbf{Collocations fréquentes (à apprendre par cœur)} \\
\hline
\all{die Einmischung, -en} & \all{sich einmischen in + Akk.} (s’ingérer dans) ; \all{unparteiisch} (impartial) & \all{sich in die inneren Angelegenheiten eines Landes einmischen} ; \all{Einmischung von außen} ; \all{Einmischung verbieten / kritisieren} \\
\hline
\all{die Macht, -¨e} & \all{mächtig} (puissant) ; \all{die Machtausübung} & \all{Macht ausüben} ; \all{an der Macht sein} ; \all{Machtmissbrauch} ; \all{Machtwechsel} \\
\hline
\all{das Land, -¨er} & \all{ländlich} (rural) ; \all{das Ausland} (l’étranger) & \all{im In- und Ausland} ; \all{Entwicklungsland} ; \all{Industrieland} ; \all{Land verlassen / regieren} \\
\hline
\all{die Souveränität} (meist Sg.) & \all{souverän} (souverain) ; \all{die Souveränitätsrechte} (pl.) & \all{die Souveränität wahren / respektieren / verletzen / einschränken} ; \all{nationale Souveränität} \\
\hline
\all{der Einfluss, -¨e} & \all{beeinflussen} (influencer) ; \all{einflussreich} (influent) & \all{Einfluss ausüben auf + Akk.} ; \all{unter dem Einfluss stehen von + Dat.} ; \all{großen Einfluss haben} \\
\hline
\all{der Staat, -en} & \all{staatlich} (étatique) ; \all{der Staatsbürger, -} (le citoyen) & \all{der souveräne Staat} ; \all{Staatsgrenzen} ; \all{Staatschef} ; \all{Staatsangehörigkeit} ; \all{im Staate} (dat. formel) \\
\hline
\end{tabularx}
\end{center}
\end{propriete}

\begin{aretenir}[Faux amis et pièges de traduction]
\begin{itemize}
\item \all{die Macht} = le pouvoir, la puissance (jamais « la marche »).
\item \all{die Politik} = la politique (jamais « la politesse » → \all{die Höflichkeit}).
\item \all{sich einmischen in + Akkusativ} = s’ingérer dans / se mêler de (pas *\all{sich mischen in}).
\item \all{beeinflussen} = influencer (pas *\all{influenzieren}, qui n’existe pas).
\item \all{die Stimme} = la voix (vote) / la voix (son) ; \all{seine Stimme abgeben} = voter.
\item \all{die Wahl} = l’élection ; \all{die Wahlkampagne} / \all{der Wahlkampf} = la campagne électorale.
\end{itemize}
\end{aretenir}

\begin{exoresolu}[Vocabulaire — complétion et traduction]
\textbf{Aufgabe} : Complétez avec le nom convenable (\all{Einmischung, Macht, Souveränität, Einfluss, Staat, Land}), mettez-le à la bonne forme (cas, nombre), puis traduisez la phrase complète.
\begin{enumerate}
\item \all{Jeder \_\_\_\_\_\_\_\_ hat das Recht auf eigene Entscheidungen.}
\item \all{Die \_\_\_\_\_\_\_\_ fremder Mächte in die Wahlen wurde international kritisiert.}
\item \all{Der \_\_\_\_\_\_\_\_ der sozialen Medien auf die Jugend ist enorm.}
\item \all{Ohne \_\_\_\_\_\_\_\_ kann ein \_\_\_\_\_\_\_\_ nicht frei handeln.}
\item \all{Viele afrikanische \_\_\_\_\_\_\_\_ fordern ein Ende der wirtschaftlichen Abhängigkeit.}
\end{enumerate}
\tcblower
\textbf{Solutions et traductions} :
\begin{enumerate}
\item \all{Jeder \textbf{Staat} hat das Recht auf eigene Entscheidungen.} (Nominatif sing. masc. après \all{jeder}). — « Chaque État a le droit à ses propres décisions. »
\item \all{Die \textbf{Einmischung} fremder Mächte in die Wahlen wurde international kritisiert.} (Nominatif sing. fém. ; \all{fremder Mächte} = génitif pluriel). — « L’ingérence de puissances étrangères dans les élections a été critiquée internationalement. »
\item \all{Der \textbf{Einfluss} der sozialen Medien auf die Jugend ist enorm.} (Nominatif sing. masc.). — « L’influence des réseaux sociaux sur la jeunesse est énorme. »
\item \all{Ohne \textbf{Souveränität} kann ein \textbf{Staat} nicht frei handeln.} (Accusatif après \all{ohne} → \all{Souveränität} Sg. fém. ; \all{ein Staat} nominatif/accusatif sing. masc.). — « Sans souveraineté, un État ne peut pas agir librement. »
\item \all{Viele afrikanische \textbf{Länder} fordern ein Ende der wirtschaftlichen Abhängigkeit.} (Nominatif/accusatif pluriel neutre : \all{die Länder}). — « De nombreux pays africains exigent la fin de la dépendance économique. »
\end{enumerate}
\end{exoresolu}

\courssub{Grammaire 1 : Le passif (Vorgangspassiv) — formation et emplois}

\begin{propriete}[Règle de formation du Vorgangspassiv (passif de procès)]
Le passif de procès décrit une \textbf{action} ou un \textbf{processus}. Il se forme avec l’auxiliaire \cle{werden} (conjugué au temps de la phrase active) + \cle{Partizip II} du verbe lexical (en position finale).
\begin{itemize}
\item \textbf{Sujet passif} = COD (accusatif) de l’actif $\rightarrow$ devient nominatif.
\item \textbf{Agent} (facultatif) = \all{von + Dativ} (personne) ou \all{durch + Akkusatif} (moyen/instrument).
\item \textbf{Verbes exclus} : verbes intransitifs sans COD (sauf \all{es} passif impersonnel), verbes d’état (\all{sein, haben, werden, kennen, wissen}...).
\end{itemize}
\end{propriete}

\begin{center}
\begin{tabularx}{\linewidth}{|X|X|X|}
\hline
\rowcolor{popblueL} \textbf{Temps} & \textbf{Actif (Aktiv)} & \textbf{Passif (Vorgangspassiv)} \\
\hline
Präsens & \all{Die Regierung respektiert die Souveränität.} & \all{Die Souveränität \textbf{wird} respektiert.} \\
\hline
Präteritum & \all{Die Regierung respektierte die Souveränität.} & \all{Die Souveränität \textbf{wurde} respektiert.} \\
\hline
Perfekt & \all{Die Regierung hat die Souveränität respektiert.} & \all{Die Souveränität \textbf{ist} respektiert \textbf{worden}.} \\
\hline
Plusquamperfekt & \all{Die Regierung hatte die Souveränität respektiert.} & \all{Die Souveränität \textbf{war} respektiert \textbf{worden}.} \\
\hline
Futur I & \all{Die Regierung wird die Souveränität respektieren.} & \all{Die Souveränität \textbf{wird} respektiert \textbf{werden}.} \\
\hline
\end{tabularx}
\end{center}

\begin{attention}[Piège majeur : Perfekt du passif]
\all{worden} (sans \emph{ge-}) est le Partizip II de \all{werden} \textbf{au passif}.
\all{geworden} (avec \emph{ge-}) est le Partizip II de \all{werden} \textbf{à l’actif} (devenir).
\begin{itemize}
\item Correct : \all{Die Wahlen sind manipuliert \textbf{worden}.} (Les élections ont été manipulées.)
\item Faux : *\all{Die Wahlen sind manipuliert \textbf{geworden}.}
\item Actif : \all{Er ist zum Diktator \textbf{geworden}.} (Il est devenu dictateur.)
\end{itemize}
\end{attention}

\begin{methode}[Transformer une phrase active en passif (étapes)]
\begin{enumerate}
\item Identifier le \textbf{sujet actif}, le \textbf{verbe} (temps, mode) et le \textbf{COD (accusatif)}.
\item Le COD devient \textbf{sujet passif} (nominatif) $\rightarrow$ accorder \all{werden} avec ce nouveau sujet.
\item Conjuguer \all{werden} au \textbf{même temps} que le verbe actif.
\item Placer le \textbf{Partizip II} du verbe lexical en \textbf{fin de phrase}.
\item (Optionnel) Reconvertir l’ancien sujet en \textbf{complément d’agent} : \all{von + Dativ} (personne) / \all{durch + Akkusativ} (instrument).
\end{enumerate}
\end{methode}

\begin{exoresolu}[Transformation active $\rightarrow$ passif (même temps) — type Bac]
\textbf{Aufgabe} : Mettez au passif (conservez le temps) :
\begin{enumerate}
\item \all{Die Großmächte beeinflussen die Politik des Landes.} (Präsens)
\item \all{Fremde Staaten finanzierten die Wahlkampagne.} (Präteritum)
\item \all{Die Medien haben die Wahrheit verbreitet.} (Perfekt)
\item \all{Man wird die Verträge überprüfen.} (Futur I — sujet indéfini \all{man})
\item \all{Die Kolonialmächte hatten die Grenzen gezogen.} (Plusquamperfekt)
\end{enumerate}
\tcblower
\textbf{Solutions détaillées} :
\begin{enumerate}
\item \all{Die Politik des Landes \textbf{wird} von den Großmächten \textbf{beeinflusst}.} (Sujet fém. sg. $\rightarrow$ \all{wird} ; agent \all{von den Großmächten} dat. pl.)
\item \all{Die Wahlkampagne \textbf{wurde} von fremden Staaten \textbf{finanziert}.} (Präteritum ; \all{von fremden Staaten} dat. pl. adjectif \all{fremden}).
\item \all{Die Wahrheit \textbf{ist} von den Medien \textbf{verbreitet worden}.} (Perfekt passif : \all{ist} + Part. II + \all{worden} final).
\item \all{Die Verträge \textbf{werden} \textbf{überprüft werden}.} (Futur I : \all{werden} + Part. II + \all{werden} final. \all{Man} disparaît.)
\item \all{Die Grenzen \textbf{waren} von den Kolonialmächten \textbf{gezogen worden}.} (Plusquamperfekt passif : \all{waren} + Part. II + \all{worden}).
\end{enumerate}
\end{exoresolu}

\courssub{Grammaire 2 : Le plus-que-parfait (Plusquamperfekt) et le subjonctif II (Konjunktiv II)}

\begin{propriete}[Plusquamperfekt — antériorité dans le passé]
\textbf{Formation} : auxiliaire \cle{hatte} (haben) ou \cle{war} (sein) au Präteritum + \cle{Partizip II} (finale).
\textbf{Emploi} : action antérieure à un autre passé (souvent introduit par \all{nachdem}, \all{bevor}, \all{als}).
\begin{center}
\begin{tabular}{|l|l|l|}
\hline
\rowcolor{popblueL} \textbf{Verbe} & \textbf{Auxiliaire} & \textbf{Plusquamperfekt (ich/er/sie/es)} \\
\hline
\all{beeinflussen} (régulier) & haben & \all{hatte beeinflusst} \\
\hline
\all{einmischen} (réfléchi, mouvement) & sein & \all{war sich eingemischt} \\
\hline
\all{entscheiden} & haben & \all{hatte entschieden} \\
\hline
\all{kommen / gehen / werden} & sein & \all{war gekommen / gegangen / geworden} \\
\hline
\end{tabular}
\end{center}
\all{Beispiel: Nachdem die Mächte die Wahlen \textbf{beeinflusst hatten}, protestierte das Volk.}
\end{propriete}

\begin{propriete}[Konjunktiv II — irréel du présent / futur, souhait, politesse]
\textbf{Formation générale} : \cle{würde} + \cle{infinitif} (valable pour presque tous les verbes).
\textbf{Formes propres (à connaître par cœur)} :
\begin{center}
\begin{tabular}{|l|l|l|l|}
\hline
\rowcolor{popblueL} \textbf{Infinitif} & \textbf{Präteritum} & \textbf{Konjunktiv II (ich/er/sie/es)} & \textbf{Traduction} \\
\hline
\all{sein} & \all{war} & \all{\textbf{wäre}} & serait / fût \\
\hline
\all{haben} & \all{hatte} & \all{\textbf{hätte}} & aurait / eût \\
\hline
\all{können} & \all{konnte} & \all{\textbf{könnte}} & pourrait \\
\hline
\all{müssen} & \all{musste} & \all{\textbf{müsste}} & devrait \\
\hline
\all{sollen} & \all{sollte} & \all{\textbf{sollte}} & devrait (obligation morale) \\
\hline
\all{wollen} & \all{wollte} & \all{\textbf{wollte}} & voudrait \\
\hline
\all{geben} & \all{gab} & \all{\textbf{gäbe}} & il y aurait \\
\hline
\all{werden} & \all{wurde} & \all{\textbf{würde}} & deviendrait / auxiliaire KII \\
\hline
\end{tabular}
\end{center}
\end{propriete}

\begin{methode}[Emplois du Konjunktiv II — repères]
\begin{enumerate}
\item \cle{Hypothèse irréelle (si + KII, principale KII)} : \all{Wenn es keine Einmischung \textbf{gäbe}, \textbf{wäre} das Land freier.} (Si pas d’ingérence $\rightarrow$ pays plus libre.)
\item \cle{Souhait irréel (\all{doch})} : \all{Wenn die Staaten \textbf{doch} unabhängiger \textbf{wären}!} (Si seulement les États étaient plus indépendants !)
\item \cle{Conseil / politesse} : \all{Ich \textbf{wäre} Ihnen dankbar.} (Je vous serais reconnaissant.) / \all{Könnten Sie mir helfen?}
\item \cle{Style indirect (discours rapporté)} : \all{Er sagte, er \textbf{könne} nicht kommen.} (Il dit qu’il ne peut pas venir — KII remplace KI si ambiguïté).
\end{enumerate}
\end{methode}

\begin{attention}[Concordance des temps avec \all{nachdem} / \all{bevor}]
\begin{itemize}
\item \all{nachdem} + \textbf{Plusquamperfekt} $\rightarrow$ principale au \textbf{Präteritum} (ou Perfekt à l’oral).
  \all{Nachdem der Staat die Souveränität \textbf{verletzt hatte}, reagierte die UNO.}
\item \all{bevor} + \textbf{Präteritum/Perfekt} $\rightarrow$ principale au \textbf{Plusquamperfekt}.
  \all{Bevor die Wahl \textbf{stattfand}, hatte man die Medien \textbf{kontrolliert}.}
\item \textbf{Jamais} deux présents, jamais \all{nachdem} + Présent.
\end{itemize}
\end{attention}

\begin{exoresolu}[Thème grammatical — traduction vers l’allemand]
\textbf{Aufgabe} : Traduisez en allemand en mobilisant Plusquamperfekt et Konjunktiv II.
\begin{enumerate}
\item « Après que les puissances étrangères eurent influencé les élections, le peuple protesta. »
\item « Si l’État respectait la souveraineté des autres pays, il n’y aurait pas de conflits. »
\item « Si seulement il n’y avait pas d’ingérence ! »
\item « Avant que le gouvernement ne décide, les puissances étrangères s’étaient déjà ingérées. »
\end{enumerate}
\tcblower
\textbf{Solutions commentées} :
\begin{enumerate}
\item \all{Nachdem die ausländischen Mächte die Wahlen \textbf{beeinflusst hatten}, \textbf{protestierte} das Volk.}
  (\all{nachdem} + Plusquamperfekt \all{hatten beeinflusst} ; principale Präteritum \all{protestierte}. Verbe final en subordonnée, verbe 2\up{e} pos. en principale.)
\item \all{Wenn der Staat die Souveränität der anderen Länder \textbf{respektieren würde}, \textbf{gäbe} es keine Konflikte.}
  (Hypothèse irréelle présent : KII des deux côtés. \all{respektieren würde} ou \all{respektierte} ; \all{gäbe} = KII de \all{geben} « il y a ».)
\item \all{Wenn es \textbf{doch} keine Einmischung \textbf{gäbe}!} (Souhait irréel : \all{doch} + KII \all{gäbe}.)
\item \all{Bevor die Regierung \textbf{entschied}, hatten sich fremde Mächte \textbf{schon eingemischt}.}
  (\all{bevor} + Präteritum \all{entschied} ; principale Plusquamperfekt \all{hatten ... eingemischt} — \all{sich einmischen} verbe réfléchi mouvement $\rightarrow$ auxiliaire \all{sein}.)
\end{enumerate}
\end{exoresolu}

\courssub{Méthode : Comprendre, commenter et produire}

\begin{methode}[Lire et exploiter un texte politique (global $\rightarrow$ détaillé)]
\begin{enumerate}
\item \cle{Première lecture globale} (sans dictionnaire) : identifier le \textbf{thème} (\all{die Einmischung in die Politik}), la \textbf{source}, le \textbf{ton} (critique, neutre, engagé) et les \textbf{mots transparents} (\all{die Souveränität, der Staat, die Nation, die Unabhängigkeit}).
\item \cle{Repérage des acteurs} : souligner \textbf{qui agit} (\all{die Großmächte, die Regierung, die Kolonialmächte}) et \textbf{qui subit} (\all{das Volk, die Bürger, die jungen Staaten}). $\rightarrow$ prépare le passif.
\item \cle{Deuxième lecture détaillée} : pour chaque question, \textbf{localiser} le passage, \textbf{souligner} la phrase-clé, \textbf{reformuler} avec ses mots (allemand) — jamais recopier le paragraphe.
\item \cle{Exploitation du lexique} : déduire le sens par la \textbf{famille de mots} (\all{mischen $\rightarrow$ sich einmischen $\rightarrow$ die Einmischung} ; \all{der Einfluss $\rightarrow$ beeinflussen $\rightarrow$ einflussreich}).
\item \cle{Rédaction de la réponse} : \textbf{phrase complète}, verbe conjugué en \textbf{2\up{e} position}, reprise du \textbf{temps de la question} (Présens $\rightarrow$ Présens ; Präteritum $\rightarrow$ Präteritum/Perfekt).
\end{enumerate}
\end{methode}

\begin{methode}[Rédiger un résumé puis une prise de position (expression écrite guidée — 8–10 phrases)]
\begin{enumerate}
\item \cle{Résumé (3–5 phrases)} : idées principales uniquement (qui ? quoi ? conséquence ?) ; \textbf{Präsens} ; connecteurs : \all{zuerst, dann, deshalb, schließlich, insgesamt}.
\item \cle{Introduction de l’opinion} : \all{Meiner Meinung nach ...} (verbe en 2\up{e} pos.) / \all{Ich bin der Meinung, dass ...} (verbe en \textbf{fin} après \all{dass}).
\item \cle{Argumenter (2 arguments)} : \all{Erstens ... / Zweitens ...} + \textbf{exemple concret} : \all{Zum Beispiel ...} / \all{Ein Beispiel ist ...}
\item \cle{Nuancer avec le Konjunktiv II} : \all{Ohne Einmischung \textbf{könnten} die Länder freier entscheiden.} / \all{Wenn die Souveränität \textbf{respektiert würde}, \textbf{gäbe} es weniger Konflikte.}
\item \cle{Conclure} : \all{Zusammenfassend lässt sich sagen, dass die Souveränität jedes Staates \textbf{respektiert werden muss}.} (Passif + modal : Part. II + \all{werden} + modal conjugué final).
\end{enumerate}
\end{methode}

\begin{exoresolu}[Compréhension de texte — type Bac (d’après le texte support)]
\textbf{Fragen zum Text} :
\begin{enumerate}
\item \all{Worum geht es in dem Text? Nennen Sie das Hauptthema.}
\item \all{Welche historischen Beispiele für Einmischung nennt der Text für Kamerun?}
\item \all{Wodurch wird die Souveränität der jungen Staaten heute noch eingeschränkt? (Nennen Sie drei Punkte.)}
\item \all{Was fordern die Bürgerinnen und Bürger?}
\item \all{Welche modernen Formen der Einmischung werden in Deutschland und der Schweiz beobachtet?}
\end{enumerate}
\tcblower
\textbf{Réponses rédigées (modèles)} :
\begin{enumerate}
\item \all{Der Text handelt von der politischen Einmischung ehemaliger Kolonialmächte und moderner Mächte in die inneren Angelegenheiten souveräner Staaten, besonders in Afrika.}
\item \all{Für Kamerun nennt der Text Frankreich und Großbritannien als ehemalige Kolonialmächte, die bis heute die Amtssprachen, das Rechtssystem und die Verwaltung prägen.}
\item \all{Die Souveränität wird durch Wirtschaftsverträge, an Bedingungen geknüpfte Entwicklungshilfe und militärische Kooperationen eingeschränkt.}
\item \all{Die Bürger fordern eine Politik ohne äußere Einmischung und wollen ihre eigenen Entscheidungen treffen.}
\item \all{In Deutschland und der Schweiz beobachtet man Wahlbeeinflussung durch Cyberangriffe, Desinformation in sozialen Medien und wirtschaftlichen Druck großer Konzerne.}
\end{enumerate}
\end{exoresolu}

\begin{exoresolu}[Production écrite guidée — type Bac]
\textbf{Aufgabe} : \all{Schreiben Sie einen zusammenhängenden Text (8–10 Sätze) zum Thema: „Sollen sich fremde Staaten in die Politik anderer Länder einmischen?“ Begründen Sie Ihre Meinung mit zwei Argumenten und einem Beispiel. Nutzen Sie mindestens ein Passiv und einen Konjunktiv II.}
\tcblower
\textbf{Texte modèle (corrigé type)} :

\all{Die Einmischung fremder Staaten in die Politik anderer Länder ist ein zentrales Problem der internationalen Beziehungen. Meiner Meinung nach soll sich kein Staat in die inneren Angelegenheiten eines anderen Landes einmischen. Erstens hat jedes Volk das unveräußerliche Recht, seine Regierung frei und ohne Druck von außen zu wählen. Wenn Wahlen von ausländischen Mächten beeinflusst werden, verliert das Volk seine demokratische Stimme. Zweitens wird die Souveränität eines Staates durch jede Form der Einmischung verletzt. Zum Beispiel haben viele afrikanische Länder nach dem Kolonialismus jahrzehntelang unter dem politischen und wirtschaftlichen Einfluss der ehemaligen Kolonialmächte gelitten. Wenn es keine Einmischung gäbe, könnten diese Länder ihre Entwicklung eigenständiger gestalten. Zusammenfassend bin ich der Meinung, dass die Souveränität jedes Staates respektiert werden muss, damit Frieden und Gerechtigkeit möglich sind.}

\textbf{Analyse grammaticale du modèle} :
\begin{itemize}
\item Passif Présens : \all{werden beeinflusst} ; \all{wird ... verletzt}.
\item Passif + modal (infinitif) : \all{respektiert werden muss}.
\item Konjunktiv II : \all{gäbe} (hypothèse), \all{könnten} (conséquence).
\item Subordonnées : verbe final après \all{dass}, \all{wenn}, \all{damit}.
\item Connecteurs : \all{Erstens, Zweitens, Zum Beispiel, Zusammenfassend}.
\end{itemize}
\end{exoresolu}

\courssub{Entraînement : Exercices d’application (Bac blanc)}

\begin{exercice}[Grammaire \& Vocabulaire — Entraînement individuel]
\textbf{A. Passif :} Mettez au passif (même temps). Indiquez l’agent seulement si nécessaire.
\begin{enumerate}
\item \all{Die UNO verurteilt die Einmischung.} (Präsens)
\item \all{Man hat die Verträge gebrochen.} (Perfekt — sujet \all{man})
\item \all{Die Großmächte bestimmten die Politik.} (Präteritum)
\item \all{Man wird die Souveränität achten.} (Futur I)
\item \all{Die Kolonialmächte hatten die Ressourcen ausgebeutet.} (Plusquamperfekt)
\end{enumerate}

\textbf{B. Plusquamperfekt / Konjunktiv II :} Conjuguez le verbe entre parenthèses au temps indiqué.
\begin{enumerate}
\item \all{Nachdem der Staat die Unabhängigkeit (erlangen) \_\_\_\_\_\_\_\_, (beginnen) \_\_\_\_\_\_\_\_ eine neue Ära.} (Plusquamperfekt + Präteritum)
\item \all{Wenn die Mächte die Souveränität (respektieren) \_\_\_\_\_\_\_\_, (geben) \_\_\_\_\_\_\_\_ es weniger Kriege.} (Konjunktiv II)
\item \all{Ich wünschte, die Einmischung (sein) \_\_\_\_\_\_\_\_ vorbei!} (Konjunktiv II souhait)
\item \all{Bevor die Wahl (stattfinden) \_\_\_\_\_\_\_\_, die Medien (kontrollieren) \_\_\_\_\_\_\_\_ schon.} (Präteritum + Plusquamperfekt)
\end{enumerate}

\textbf{C. Vocabulaire :} Complétez avec le mot juste (forme correcte : cas, nombre, article).
\all{Einmischung, Macht, Souveränität, Einfluss, Staat, Land}
\begin{enumerate}
\item \all{Die \_\_\_\_\_\_\_\_ der Supermächte auf die UNO ist groß.}
\item \all{Nach dem Krieg verloren viele \_\_\_\_\_\_\_\_ ihre \_\_\_\_\_\_\_\_.}
\item \all{Ohne \_\_\_\_\_\_\_\_ ist kein \_\_\_\_\_\_\_\_ wirklich frei.}
\item \all{Die \_\_\_\_\_\_\_\_ in die inneren Angelegenheiten verstößt gegen das Völkerrecht.}
\end{enumerate}

\textbf{D. Traduction (Thème) :} Traduisez en allemand.
\begin{enumerate}
\item « L’ingérence de l’étranger a été critiquée par le gouvernement. »
\item « Si le pays était souverain, il n’accepterait pas cette pression. »
\item « Après que les troupes eurent quitté le pays, la paix revint. »
\end{enumerate}
\end{exercice}

\begin{corrige}[Exercices d’application]
\textbf{A. Passif}
\begin{enumerate}
\item \all{Die Einmischung wird von der UNO verurteilt.}
\item \all{Die Verträge sind gebrochen worden.} (\all{man} disparaît)
\item \all{Die Politik wurde von den Großmächten bestimmt.}
\item \all{Die Souveränität wird geachtet werden.}
\item \all{Die Ressourcen waren von den Kolonialmächten ausgebeutet worden.}
\end{enumerate}

\textbf{B. Temps composés / KII}
\begin{enumerate}
\item \all{Nachdem der Staat die Unabhängigkeit \textbf{erlangt hatte}, \textbf{begann} eine neue Ära.} (\all{erlangen} $\rightarrow$ \all{hatte erlangt} ; \all{beginnen} Präteritum \all{begann})
\item \all{Wenn die Mächte die Souveränität \textbf{respektieren würden} (ou \textbf{respektierten}), \textbf{gäbe} es weniger Kriege.}
\item \all{Ich wünschte, die Einmischung \textbf{wäre} vorbei!} (\all{sein} $\rightarrow$ KII \all{wäre})
\item \all{Bevor die Wahl \textbf{stattfand}, \textbf{hatten} die Medien \textbf{schon kontrolliert}.} (\all{bevor} + Präteritum ; principale Plusquamperfekt)
\end{enumerate}

\textbf{C. Vocabulaire}
\begin{enumerate}
\item \all{Der \textbf{Einfluss} der Supermächte auf die UNO ist groß.} (Nominatif masc. sg.)
\item \all{Nach dem Krieg verloren viele \textbf{Länder} ihre \textbf{Souveränität}.} (Accusatif pluriel neutre \all{Länder} ; Accusatif fém. sg. \all{Souveränität})
\item \all{Ohne \textbf{Souveränität} ist kein \textbf{Staat} wirklich frei.} (Accusatif après \all{ohne} : \all{Souveränität} fém., \all{Staat} masc. — prédicat au nominatif après \all{sein})
\item \all{Die \textbf{Einmischung} in die inneren Angelegenheiten verstößt gegen das Völkerrecht.} (Nominatif fém. sg.)
\end{enumerate}

\textbf{D. Thème}
\begin{enumerate}
\item \all{Die Einmischung aus dem Ausland wurde von der Regierung kritisiert.} (ou \all{Die ausländische Einmischung ...})
\item \all{Wenn das Land souverän wäre, würde es diesen Druck nicht akzeptieren.} (KII \all{wäre} / \all{würde ... akzeptieren})
\item \all{Nachdem die Truppen das Land verlassen hatten, kehrte der Frieden zurück.} (Plusquamperfekt \all{hatten ... verlassen} — \all{verlassen} aux. \all{haben} ; Präteritum \all{kehrte ... zurück})
\end{enumerate}
\end{corrige}

\begin{attention}[Les 5 erreurs qui coûtent des points au Bac — Récapitulatif]
\begin{enumerate}
\item \cle{Confondre \all{worden} et \all{geworden}} : Perfekt passif = \all{ist ... Part. II \textbf{worden}} ; Actif \all{werden} = \all{ist ... \textbf{geworden}}.
\item \cle{Oublier majuscule et article des noms} : *\all{die einmischung} $\rightarrow$ \all{die \textbf{E}inmischung}. Apprendre \textbf{chaque nom avec son article et son pluriel} : \all{das Land, die Länder} ; \all{der Staat, die Staaten}.
\item \cle{Calquer l’ordre des mots français} : Après \all{dass, weil, wenn, nachdem, bevor} $\rightarrow$ verbe conjugué en \textbf{fin de phrase}. *\all{weil die Souveränität wurde verletzt} $\rightarrow$ \all{weil die Souveränität \textbf{verletzt wurde}}.
\item \cle{Faux amis et calques} : \all{Macht} = pouvoir ; \all{sich einmischen in + Akk.} (pas *\all{in + Dat.} ni *\all{mischen}) ; \all{beeinflussen} (pas *\all{influenzieren}) ; \all{die Stimme abgeben} = voter.
\item \cle{Mauvaise concordance \all{nachdem} / \all{bevor}} : \all{nachdem} + Plusquamperfekt $\rightarrow$ principale Präteritum. \all{bevor} + Präteritum $\rightarrow$ principale Plusquamperfekt. Ne jamais mettre deux présents.
\end{enumerate}
\end{attention}

\begin{savaistu}[Contexte : Cameroun, Allemagne, Suisse — Histoire et actualité]
Le Cameroun, indépendant depuis 1960, porte l’héritage d’une double colonisation (France, Royaume-Uni) : deux systèmes juridiques, deux langues officielles, une administration bilingue. La \all{Einmischung} post-coloniale s’est exercée via la \all{Françafrique} (accords de défense, monnaie CFA, conseillers techniques). En Allemagne, le débat porte sur l’\all{Wahlbeeinflussung} (ingérence électorale) via cyberattaques (\all{Cyberangriffe}) et désinformation (\all{Desinformation}), attribuées à des acteurs étatiques russes ou chinois. La Suisse, pays neutre, veille à sa \all{Souveränität} face aux pressions de l’UE et des grandes puissances. La Charte de l’ONU (Art. 2, al. 7) consacre le \all{Prinzip der Nichteinmischung} — mais la \all{Realpolitik} montre souvent l’écart entre droit et pratique.
\end{savaistu}

\begin{experience}[Activité orale : Débat structuré — « Souveraineté vs Ingérence humanitaire »]
\textbf{Consigne} : En groupes de 4. Deux élèves défendent : « La souveraineté est absolue, aucune ingérence n’est légitime. » Deux élèves défendent : « L’ingérence humanitaire (R2P — Responsabilité de protéger) est un devoir moral face aux crimes de masse. »
\textbf{Préparation (10 min)} : Chaque camp note 3 arguments en allemand (vocabulaire de la leçon : \all{Einmischung, Souveränität, Menschenrechte, Völkerrecht, UN-Charta, militärische Intervention}).
\textbf{Débat (10 min)} : Tour de parole 1 min chacun. Utiliser : \all{Meiner Meinung nach ... ; Erstens ... ; Ein Gegenargument ist ... ; Ich sehe das anders, weil ... ; Zusammenfassend ...}
\textbf{Retour (5 min)} : La classe vote pour l’argumentation la plus cohérente linguistiquement (passif, KII, connecteurs) et logiquement.
\end{experience}

\newpage

\coursec{Exercices}

\courssub{Je vérifie mes connaissances}

\begin{exercice}[QCM : grammaire et vocabulaire]
Entourez la bonne réponse.
\begin{enumerate}
\item \all{Die Souveränität eines Staates muss \ldots\ werden.} \\
a) respektieren \quad b) respektiert \quad c) respektierte \quad d) respektierst
\item Le plus-que-parfait de \all{entscheiden} à la 3\ieme{} personne du singulier est : \\
a) \all{hat entschieden} \quad b) \all{hatte entschieden} \quad c) \all{ist entschieden} \quad d) \all{entschied}
\item Au passif Präteritum, \all{werden} devient : \\
a) \all{wurden} \quad b) \all{worden} \quad c) \all{geworden} \quad d) \all{werdeten}
\item \all{die Einmischung} signifie : \\
a) la puissance \quad b) l'ingérence \quad c) la souveraineté \quad d) l'influence
\item \all{Wenn die Menschenrechte geachtet würden, \ldots\ keine Einmischung nötig.} \\
a) wird \quad b) wäre \quad c) ist \quad d) wurde
\end{enumerate}
\end{exercice}

\begin{exercice}[Vrai ou faux ? Justifiez]
Répondez par \all{richtig} ou \all{falsch} et corrigez les phrases fausses.
\begin{enumerate}
\item \all{Das Land wurde von den Sanktionen betroffen.} est une phrase au passif.
\item Au plus-que-parfait, tous les verbes se conjuguent avec \all{haben}.
\item \all{Der Staat mischt sich in die inneren Angelegenheiten ein.} — le verbe \all{einmischen} est réfléchi.
\item Le subjonctif II exprime un fait réel et certain.
\item Le pluriel de \all{die Macht} est \all{die Mächte}.
\end{enumerate}
\end{exercice}

\begin{exercice}[Vocabulaire : articles et pluriels]
Complétez le tableau avec l'article et le pluriel de chaque nom, puis traduisez-le en français.

\begin{center}
\begin{tabularx}{\linewidth}{|X|l|l|X|}
\hline
\rowcolor{popblueL} Nom & Article & Pluriel & Traduction \\
\hline
Einmischung & \ldots & \ldots & \ldots \\
\hline
Macht & \ldots & \ldots & \ldots \\
\hline
Land & \ldots & \ldots & \ldots \\
\hline
Souveränität & \ldots & \ldots & \ldots \\
\hline
Einfluss & \ldots & \ldots & \ldots \\
\hline
Staat & \ldots & \ldots & \ldots \\
\hline
\end{tabularx}
\end{center}
\end{exercice}

\begin{exercice}[Conjugaison : le plus-que-parfait]
Complétez les phrases avec le verbe entre parenthèses au plus-que-parfait. Attention au choix entre \all{haben} et \all{sein}.
\begin{enumerate}
\item Nachdem die UNO die Resolution \ldots\ (beschließen), reagierte die Regierung.
\item Der Botschafter \ldots\ das Land schon \ldots\ (verlassen), als die Krise begann.
\item Bevor die Wahl stattfand, \ldots\ die Bevölkerung lange \ldots\ (demonstrieren).
\item Nachdem die Sanktionen \ldots\ \ldots\ (werden — au passif !), änderte das Parlament seine Position.
\item Die Journalisten \ldots\ schon nach Yaoundé \ldots\ (reisen), als der Konflikt eskalierte.
\end{enumerate}
\end{exercice}

\courssub{Je m'entraîne}

\begin{exercice}[Transformation : de l'actif au passif]
Mettez les phrases suivantes au passif, en respectant le temps indiqué entre parenthèses.
\begin{enumerate}
\item Die internationale Gemeinschaft verurteilt die Einmischung. (Präsens)
\item Der Präsident unterzeichnete das Abkommen. (Präteritum)
\item Die Medien haben die Wahlergebnisse kritisiert. (Perfekt)
\item Die Großmächte beeinflussen die Politik des Landes. (Präsens)
\item Man respektierte die Souveränität des Staates nicht. (Präteritum)
\item Die Nachbarstaaten haben die Grenzen geschlossen. (Perfekt)
\end{enumerate}
\end{exercice}

\begin{exercice}[Texte à trous : passif et plus-que-parfait]
Complétez le texte avec la forme correcte du verbe entre parenthèses (passif ou plus-que-parfait selon le sens).

\begin{textebox}[Eine Krise in der Region]
Nachdem der Konflikt im Nachbarland \ldots\ (ausbrechen), \ldots\ tausende Flüchtlinge an der Grenze \ldots\ (empfangen — passif, Präteritum). Die humanitäre Hilfe \ldots\ schnell \ldots\ (organisieren — passif, Perfekt), aber viele Familien \ldots\ schon alles \ldots\ (verlieren). Erst nach langen Verhandlungen \ldots\ ein Waffenstillstand \ldots\ (unterzeichnen — passif, Präteritum). Heute \ldots\ die Region noch immer von der internationalen Gemeinschaft \ldots\ (beobachten — passif, Präsens).
\end{textebox}
\end{exercice}

\begin{exercice}[Appariements : familles de mots]
Complétez le tableau des familles de mots, puis rédigez une phrase complète avec un mot de chaque famille.

\begin{center}
\begin{tabularx}{\linewidth}{|l|X|X|}
\hline
\rowcolor{popblueL} Verbe & Nomen & Adjektiv \\
\hline
sich einmischen & die \ldots & \ldots \\
\hline
\ldots & die Macht & \ldots \\
\hline
regieren & die \ldots & staatlich \\
\hline
\ldots & die Souveränität & souverän \\
\hline
beeinflussen & der \ldots & einflussreich \\
\hline
\end{tabularx}
\end{center}
\end{exercice}

\begin{exercice}[Thème : traduction en allemand]
Traduisez les phrases suivantes en allemand.
\begin{enumerate}
\item Après que les sanctions eurent été décidées, le gouvernement changea d'avis.
\item Si les droits de l'homme étaient respectés, aucune ingérence ne serait nécessaire.
\item La souveraineté de ce pays a été violée par les grandes puissances.
\item Le conflit avait déjà commencé quand les observateurs sont arrivés.
\item On ne doit pas s'immiscer dans les affaires intérieures d'un État souverain.
\end{enumerate}
\end{exercice}

\courssub{Je me prépare au Bac}

\begin{exercice}[Compréhension de l'écrit : texte inconnu]
Lisez le texte suivant, puis répondez aux questions \textbf{en allemand, par des phrases complètes}.

\begin{textebox}[Die UNO und die Souveränität der Staaten]
Seit ihrer Gründung im Jahr 1945 steht die Vereinten Nationen vor einem schwierigen Dilemma: Einerseits garantiert ihre Charta die Souveränität aller Mitgliedstaaten, andererseits soll sie den Weltfrieden und die Menschenrechte schützen. Nachdem im Sicherheitsrat lange über eine Resolution debattiert worden war, wurde schließlich entschieden, dass in bestimmten Fällen ein Eingreifen möglich sein soll. Kritiker argumentieren, dass solche Resolutionen oft von den Interessen der Großmächte beeinflusst werden. Kleine Staaten fühlen sich manchmal bedroht, weil ihre inneren Angelegenheiten international diskutiert werden. Befürworter dagegen betonen, dass die internationale Gemeinschaft nicht schweigen darf, wenn die Bevölkerung eines Landes unterdrückt wird. Wenn die Menschenrechte überall geachtet würden, wäre keine Einmischung nötig — doch die Wirklichkeit sieht anders aus. Die Debatte über das richtige Gleichgewicht zwischen Souveränität und Schutzverantwortung wird auch in Zukunft weitergeführt werden.
\end{textebox}

\textbf{Fragen zum Text :}
\begin{enumerate}
\item Welches Dilemma hat die UNO seit ihrer Gründung?
\item Was wurde nach den langen Debatten im Sicherheitsrat entschieden?
\item Warum fühlen sich kleine Staaten manchmal bedroht?
\item Was betonen die Befürworter eines Eingreifens?
\item Unter welcher Bedingung wäre keine Einmischung nötig?
\end{enumerate}

\textbf{Fragen zur Sprache :}
\begin{enumerate}
\item Relevez dans le texte deux phrases au passif et indiquez leur temps.
\item Relevez une phrase au plus-que-parfait et justifiez l'emploi de ce temps.
\item Relevez une phrase au subjonctif II et expliquez sa valeur (irréel, hypothèse, souhait).
\end{enumerate}
\end{exercice}

\begin{exercice}[Version type Bac]
Traduisez en français l'extrait suivant (environ 6 lignes), tiré du texte support de la leçon.

\begin{textebox}[Auszug]
Nachdem die Sanktionen von der internationalen Gemeinschaft beschlossen worden waren, änderte die Regierung ihre Politik. Die Souveränität des Landes war lange nicht respektiert worden, und die Bevölkerung hatte unter dem Druck der Großmächte gelitten. Heute wird die Einmischung in die inneren Angelegenheiten eines Staates von vielen Politikern kritisiert, denn jeder Staat soll über seine eigene Zukunft entscheiden können.
\end{textebox}

\textbf{Consignes :} respectez le sens global du texte ; traduisez les formes passives par des tournures naturelles en français ; rendez le plus-que-parfait allemand par un passé antérieur ou un plus-que-parfait français correct.
\end{exercice}

\begin{exercice}[Production écrite type Bac]
\textbf{Thema :} \all{Soll die UNO in die inneren Angelegenheiten eines Staates eingreifen?}

Rédigez un texte argumenté d'environ \textbf{12 lignes} (120--150 mots) en allemand. Structurez votre production selon le plan suivant :
\begin{enumerate}
\item \textbf{Einleitung :} présentez le problème (souveraineté contre protection des droits de l'homme).
\item \textbf{Pro :} donnez au moins deux arguments en faveur d'une intervention de l'ONU.
\item \textbf{Contra :} donnez au moins deux arguments contre l'ingérence.
\item \textbf{Schluss :} exprimez votre opinion personnelle, éventuellement en lien avec le contexte africain ou camerounais.
\end{enumerate}

\textbf{Consignes de langue :} employez au moins deux phrases au passif, une phrase au plus-que-parfait et une phrase au subjonctif II ; utilisez le vocabulaire du chapitre (\all{die Einmischung, die Souveränität, der Einfluss, der Staat, die Macht}).
\end{exercice}

\newpage

\coursec{Corrigés des exercices}

\coursec{Corrigés détaillés des exercices — Leçon AL26}

\begin{corrige}[Exercice 1 — QCM : grammaire et vocabulaire]
\begin{enumerate}
\item \textbf{b) respektiert} — \all{Die Souveränität eines Staates muss respektiert werden.} Passif avec modal : \cle{Modalverb + Partizip II + werden}. « La souveraineté d'un État doit être respectée. »
\item \textbf{b) hatte entschieden} — Plus-que-parfait : \all{haben} au Präteritum (\all{hatte}) + Partizip II (\all{entschieden}, verbe fort irrégulier : \all{entscheiden – entschied – hat entschieden}).
\item \textbf{a) wurden} — Passif Präteritum : \all{werden – wurde – ist worden}. Attention : au passif Perfekt, on dit \all{worden} (et non \all{geworden}).
\item \textbf{b) l'ingérence} — \all{die Einmischung, -en} = l'ingérence, l'immixtion (verbe : \all{sich einmischen in + Akk.}).
\item \textbf{b) wäre} — \all{Wenn die Menschenrechte geachtet würden, wäre keine Einmischung nötig.} Subjonctif II dans les deux membres (irréel du présent). \textbf{Analyse de la subordonnée} : \all{geachtet würden} = Partizip II (\all{geachtet}) + \all{würden} (Konjunktiv II de \all{werden} comme auxiliaire de passif). Ce n'est \textbf{PAS} \all{würden + Infinitif} (qui serait la forme de substitution pour l'actif : \all{würden achten}). La principale utilise \all{wäre} (Konjunktiv II de \all{sein}).
\end{enumerate}
\end{corrige}

\begin{corrige}[Exercice 2 — Vrai ou faux ? Justifiez]
\begin{enumerate}
\item \textbf{Richtig.} \all{Das Land wurde von den Sanktionen betroffen.} : auxiliaire \all{werden} au Präteritum (\all{wurde}) + Partizip II (\all{betroffen}) + complément d'agent introduit par \all{von} : c'est bien un passif Präteritum (Vorgangspassiv). « Le pays a été touché par les sanctions. »
\item \textbf{Falsch.} Les verbes de mouvement ou de changement d'état se conjuguent avec \all{sein} : \all{Nachdem der Botschafter abgereist war, \ldots} ; \all{Nachdem der Konflikt ausgebrochen war, \ldots}. Correction : « Au plus-que-parfait, la plupart des verbes se conjuguent avec \all{haben}, mais les verbes de déplacement (\all{reisen, gehen, fliehen}) et de changement d'état (\all{ausbrechen, sterben, werden}) prennent \all{sein}. »
\item \textbf{Richtig.} \all{sich einmischen in + Akkusativ} est un verbe réfléchi à particule séparable : \all{Der Staat mischt sich in die inneren Angelegenheiten ein.} Noter la particule \all{ein-} rejetée en fin de phrase (position finale).
\item \textbf{Falsch.} Le subjonctif II exprime l'irréel, l'hypothèse ou le souhait, jamais un fait certain. Correction : « Le subjonctif II exprime un fait irréel, hypothétique ou souhaité » (ex. \all{Wenn die Menschenrechte geachtet würden, wäre keine Einmischung nötig.}).
\item \textbf{Richtig.} \all{die Macht, die Mächte} : féminin en \all{-e} au pluriel avec Umlaut (\all{a} $\rightarrow$ \all{ä}), comme \all{die Stadt – die Städte}, \all{die Nacht – die Nächte}.
\end{enumerate}
\end{corrige}

\begin{corrige}[Exercice 3 — Vocabulaire : articles et pluriels]
\begin{center}
\begin{tabularx}{\linewidth}{|X|l|l|X|}
\hline
\rowcolor{popblueL} Nom & Article & Pluriel & Traduction \\
\hline
Einmischung & die & die Einmischungen & l'ingérence, l'immixtion \\
\hline
Macht & die & die Mächte & la puissance, le pouvoir \\
\hline
Land & das & die Länder & le pays \\
\hline
Souveränität & die & die Souveränitäten & la souveraineté \\
\hline
Einfluss & der & die Einflüsse & l'influence \\
\hline
Staat & der & die Staaten & l'État \\
\hline
\end{tabularx}
\end{center}
\textbf{Points de vigilance :} \all{das Land} est neutre (faux ami du genre français : « le pays » masculin) ; \all{der Einfluss} est masculin alors que « l'influence » est féminin ; \all{der Staat} est un nom faible (N-Deklination) : il prend \all{-en} au pluriel (\all{die Staaten}) et à tous les cas du singulier sauf nominatif (\all{des Staates, dem Staate, den Staat}). Apprenez toujours chaque nom avec son article et son pluriel.
\end{corrige}

\begin{corrige}[Exercice 4 — Conjugaison : le plus-que-parfait]
\begin{enumerate}
\item \all{Nachdem die UNO die Resolution \textbf{beschlossen hatte}, reagierte die Regierung.} — Auxiliaire \all{haben} (verbe transitif) ; \all{beschließen – beschloss – hat beschlossen} (verbe fort, voyelle radicale \all{ie} $\rightarrow$ \all{o} $\rightarrow$ \all{o}). « Après que l'ONU eut adopté la résolution, le gouvernement réagit. »
\item \all{Der Botschafter \textbf{hatte} das Land schon \textbf{verlassen}, als die Krise begann.} — \all{verlassen} est transitif (avec accusatif) : auxiliaire \all{haben}. Partizip II \all{verlassen} (verbe fort : \all{verlassen – verließ – hat verlassen}). « L'ambassadeur avait déjà quitté le pays quand la crise a commencé. »
\item \all{Bevor die Wahl stattfand, \textbf{hatte} die Bevölkerung lange \textbf{demonstriert}.} — Verbe faible en \all{-ieren} : Partizip II \all{demonstriert} (sans \all{ge-}). « Avant que l'élection n'ait lieu, la population avait manifesté longtemps. »
\item \all{Nachdem die Sanktionen \textbf{beschlossen worden waren}, änderte das Parlament seine Position.} — Passif au plus-que-parfait : \all{waren} (auxiliaire \all{sein} au Präteritum) + Partizip II \all{beschlossen} + \all{worden}. « Après que les sanctions eurent été décidées, le Parlement changea de position. »
\item \all{Die Journalisten \textbf{waren} schon nach Yaoundé \textbf{gereist}, als der Konflikt eskalierte.} — Verbe de déplacement (\all{reisen}) : auxiliaire \all{sein} (\all{waren}). Partizip II \all{gereist}. « Les journalistes étaient déjà partis pour Yaoundé quand le conflit s'est aggravé. »
\end{enumerate}
\textbf{Rappel :} après \all{nachdem} (et \all{bevor} si antériorité), l'allemand exige le plus-que-parfait quand la principale est au Präteritum — comme le français avec le passé antérieur (\all{après que} + passé antérieur).
\end{corrige}

\begin{corrige}[Exercice 5 — Transformation : de l'actif au passif]
\begin{enumerate}
\item \all{Die Einmischung wird von der internationalen Gemeinschaft verurteilt.} (Präsens : \all{wird} + Partizip II)
\item \all{Das Abkommen wurde vom Präsidenten unterzeichnet.} (Präteritum : \all{wurde} ; \all{vom} = \all{von dem} ; \all{der Président} est un nom faible : \all{dem Präsidenten})
\item \all{Die Wahlergebnisse sind von den Medien kritisiert worden.} (Perfekt : \all{sind} + Partizip II + \all{worden})
\item \all{Die Politik des Landes wird von den Großmächten beeinflusst.} (Präsens ; Partizip II de \all{beeinflussen} : \all{beeinflusst}, sans \all{ge-} car préfixe inséparable \all{be-})
\item \all{Die Souveränität des Staates wurde nicht respektiert.} (Präteritum ; \all{man} disparaît au passif ; la négation \all{nicht} se place avant le Partizip II)
\item \all{Die Grenzen sind von den Nachbarstaaten geschlossen worden.} (Perfekt ; \all{schließen – schloss – hat geschlossen}, verbe fort)
\end{enumerate}
\textbf{Méthode :} le complément d'objet direct (Akkusativobjekt) de l'actif devient le sujet (Nominativ) du passif ; le verbe \all{werden} prend le temps de la phrase active ; le sujet actif devient un complément d'agent avec \all{von} + Dativ.
\end{corrige}

\begin{corrige}[Exercice 6 — Texte à trous : passif et plus-que-parfait]
\begin{textebox}[Eine Krise in der Region — Lösung]
Nachdem der Konflikt im Nachbarland \textbf{ausgebrochen war}, \textbf{wurden} tausende Flüchtlinge an der Grenze \textbf{empfangen}. Die humanitäre Hilfe \textbf{ist} schnell \textbf{organisiert worden}, aber viele Familien \textbf{hatten} schon alles \textbf{verloren}. Erst nach langen Verhandlungen \textbf{wurde} ein Waffenstillstand \textbf{unterzeichnet}. Heute \textbf{wird} die Region noch immer von der internationalen Gemeinschaft \textbf{beobachtet}.
\end{textebox}
\textbf{Justifications :}
\begin{itemize}
\item \all{ausgebrochen war} : Plus-que-parfait avec \all{sein} (changement d'état : \all{ausbrechen}), exigé par \all{nachdem} ; \all{ausbrechen – brach aus – ist ausgebrochen}.
\item \all{wurden \ldots\ empfangen} : Passif Präteritum, pluriel (sujet \all{Flüchtlinge}).
\item \all{ist \ldots\ organisiert worden} : Passif Perfekt ; \all{organisiert} sans \all{ge-} (verbe en \all{-ieren}).
\item \all{hatten \ldots\ verloren} : Plus-que-parfait avec \all{haben} (verbe transitif \all{verlieren}).
\item \all{wurde \ldots\ unterzeichnet} : Passif Präteritum, singulier (sujet \all{ein Waffenstillstand}).
\item \all{wird \ldots\ beobachtet} : Passif Präsens (marqueur temporel \all{heute}).
\end{itemize}
\end{corrige}

\begin{corrige}[Exercice 7 — Appariements : familles de mots]
\begin{center}
\begin{tabularx}{\linewidth}{|l|X|X|}
\hline
\rowcolor{popblueL} Verbe & Nomen & Adjektiv / Partizip \\
\hline
sich einmischen & die Einmischung & eingemischt / einmischend \\
\hline
herrschen / mächtig sein & die Macht & mächtig \\
\hline
regieren & die Regierung & regierend / regierbar \\
\hline
souverän handeln & die Souveränität & souverän \\
\hline
beeinflussen & der Einfluss & einflussreich \\
\hline
\end{tabularx}
\end{center}
\textbf{Correction de la ligne « regieren » :} l'adjectif dérivé est \all{regierend} (participe présent) ou \all{regierbar} (gouvernable). \all{Staatlich} (étatique) dérive de \all{der Staat}, pas de \all{regieren}.
\textbf{Exemples de phrases :}
\begin{itemize}
\item \all{Die Einmischung der Großmächte wird von vielen Staaten kritisiert.} « L'ingérence des grandes puissances est critiquée par de nombreux États. »
\item \all{Die mächtigen Staaten haben oft einen großen Einfluss auf kleine Länder.} « Les États puissants ont souvent une grande influence sur les petits pays. »
\item \all{Die Regierung eines souveränen Staates entscheidet über ihre eigene Politik.} « Le gouvernement d'un État souverain décide de sa propre politique. »
\end{itemize}
\end{corrige}

\begin{corrige}[Exercice 8 — Thème : traduction en allemand]
\begin{enumerate}
\item \all{Nachdem die Sanktionen beschlossen worden waren, änderte die Regierung ihre Meinung.} — Passif au plus-que-parfait après \all{nachdem} ; verbe conjugué en position 2 dans la principale (\all{änderte} avant le sujet \all{die Regierung}).
\item \all{Wenn die Menschenrechte geachtet würden, wäre keine Einmischung nötig.} — Irréel du présent : Subjonctif II dans les deux propositions (passif \all{geachtet würden}, \all{wäre}).
\item \all{Die Souveränität dieses Landes ist von den Großmächten verletzt worden.} — Passif Perfekt ; \all{verletzen} = violer (un droit, une loi, une frontière), ne pas confondre avec « blesser » au sens physique qui se dit aussi \all{verletzen} (contexte différent).
\item \all{Der Konflikt hatte schon begonnen, als die Beobachter ankamen.} — Plus-que-parfait avec \all{haben} (\all{beginnen} prend \all{haben} au sens de « commencer » intransitif) ; \all{als} + Präteritum pour l'événement unique ponctuel.
\item \all{Man darf sich nicht in die inneren Angelegenheiten eines souveränen Staates einmischen.} — Modal \all{dürfen} + infinitif ; verbe réfléchi à particule séparable (\all{sich einmischen}) ; \all{nicht} se place avant le complément prépositionnel (\all{in die inneren Angelegenheiten}).
\end{enumerate}
\textbf{Points de vigilance :} place du verbe conjugué en deuxième position dans la principale après une subordonnée initiale ; particule séparable en fin de phrase (ou avant \all{zu} à l'infinitif) ; majuscule à tous les noms.
\end{corrige}

\begin{corrige}[Exercice 9 — Compréhension de l'écrit : texte inconnu]
\textbf{Fragen zum Text — Antworten :}
\begin{enumerate}
\item \all{Die UNO hat seit ihrer Gründung das Dilemma, dass ihre Charta einerseits die Souveränität aller Mitgliedstaaten garantiert, dass sie aber andererseits den Weltfrieden und die Menschenrechte schützen soll.}
\item \all{Nach den langen Debatten wurde entschieden, dass in bestimmten Fällen ein Eingreifen möglich sein soll.}
\item \all{Kleine Staaten fühlen sich manchmal bedroht, weil ihre inneren Angelegenheiten international diskutiert werden.}
\item \all{Die Befürworter betonen, dass die internationale Gemeinschaft nicht schweigen darf, wenn die Bevölkerung eines Landes unterdrückt wird.}
\item \all{Keine Einmischung wäre nötig, wenn die Menschenrechte überall geachtet würden.}
\end{enumerate}
\textbf{Fragen zur Sprache :}
\begin{enumerate}
\item Deux phrases au passif : \all{dass solche Resolutionen oft von den Interessen der Großmächte beeinflusst werden} (Präsens Passiv : \all{werden} + Partizip II) ; \all{wenn die Bevölkerung eines Landes unterdrückt wird} (Präsens Passiv). Autre possibilité : \all{Die Debatte \ldots\ wird auch in Zukunft weitergeführt werden} (Futur I Passiv).
\item Plus-que-parfait : \all{Nachdem im Sicherheitsrat lange über eine Resolution debattiert worden war, \ldots} — Passif au plus-que-parfait (\all{war} + Partizip II + \all{worden}) ; il exprime l'antériorité par rapport au Präteritum de la principale (\all{wurde entschieden}), conformément à la règle de \all{nachdem}.
\item Subjonctif II : \all{Wenn die Menschenrechte überall geachtet würden, wäre keine Einmischung nötig.} — Valeur d'irréel du présent : hypothèse contraire à la réalité, comme le confirme la phrase suivante (\all{doch die Wirklichkeit sieht anders aus}).
\end{enumerate}
\textbf{Barème indicatif (20 points) :} compréhension 10 (2 par réponse), questions de langue 6 (2 par question), qualité de la langue 4.
\end{corrige}

\begin{corrige}[Exercice 10 — Version type Bac]
\textbf{Traduction de référence :}

« Après que les sanctions eurent été décidées par la communauté internationale, le gouvernement changea de politique. La souveraineté du pays n'avait pas été respectée pendant longtemps, et la population avait souffert de la pression des grandes puissances. Aujourd'hui, l'ingérence dans les affaires intérieures d'un État est critiquée par de nombreux hommes politiques, car chaque État doit pouvoir décider de son propre avenir. »

\textbf{Points de vigilance :}
\begin{itemize}
\item \all{Nachdem \ldots\ beschlossen worden waren} : Passif au plus-que-parfait $\rightarrow$ Passé antérieur français (« eurent été décidées »), tournure soutenue attendue en version.
\item \all{war \ldots\ nicht respektiert worden} : Plus-que-parfait passif $\rightarrow$ « n'avait pas été respectée ».
\item \all{wird \ldots\ kritisiert} : Présent passif $\rightarrow$ « est critiquée » ; éviter la tournure lourde « est en train d'être critiquée ».
\item \all{die Großmächte} : « Les grandes puissances » (et non « les grands pouvoirs »).
\item \all{die inneren Angelegenheiten} : « Les affaires intérieures » ; \all{die Zukunft} : « L'avenir » (plus naturel que « le futur »).
\end{itemize}
\textbf{Barème indicatif (20 points) :} Respect du sens global 8, traduction des formes passives et des temps 6, richesse et correction du français 6.
\end{corrige}

\begin{corrige}[Exercice 11 — Production écrite type Bac]
\textbf{Proposition de rédaction modèle :}

\begin{quote}
\all{Die Frage, ob die UNO in die inneren Angelegenheiten eines Staates eingreifen soll, wird heute oft diskutiert. Einerseits ist die Souveränität jedes Staates ein wichtiges Prinzip, andererseits müssen die Menschenrechte überall geschützt werden.}

\all{Für ein Eingreifen spricht zuerst, dass die internationale Gemeinschaft nicht schweigen darf, wenn die Bevölkerung eines Landes unterdrückt wird. Außerdem können Konflikte oft verhindert werden, bevor sie eskalieren. Wenn man früher eingegriffen hätte, hätten viele Kriege vermieden werden können.}

\all{Gegen die Einmischung spricht jedoch, dass solche Interventionen oft von den Interessen der Großmächte beeinflusst werden. Kleine Staaten, zum Beispiel in Afrika, fühlen sich manchmal bedroht, weil ihre Souveränität nicht respektiert wird. Nachdem viele afrikanische Länder ihre Unabhängigkeit erlangt hatten, wollten sie ihre eigene Zukunft selbst bestimmen.}

\all{Meiner Meinung nach soll die UNO nur eingreifen, wenn die Menschenrechte massiv verletzt werden. Wenn alle Staaten die Rechte ihrer Bürger achten würden, wäre keine Einmischung nötig. Auch Kamerun und die anderen afrikanischen Staaten müssen über ihre eigenen Angelegenheiten entscheiden können.}
\end{quote}

\textbf{Analyse grammaticale du modèle (vérification des formes imposées) :}
\begin{itemize}
\item \textbf{Passif Präsens} : \all{wird \ldots\ diskutiert}, \all{müssen \ldots\ geschützt werden}, \all{unterdrückt wird}, \all{können \ldots\ verhindert werden}, \all{beeinflusst werden}, \all{nicht respektiert wird}, \all{verletzt werden} (7 occurrences).
\item \textbf{Plus-que-parfait} : \all{hatten \ldots\ erlangt} (après \all{nachdem}).
\item \textbf{Subjonctif II (Irréel présent)} : \all{wäre \ldots\ nötig}, \all{achten würden}.
\item \textbf{Subjonctif II (Irréel passé / conditionnel passé)} : \all{hätte \ldots\ eingegriffen} (au lieu de la forme passive impersonnelle lourde \all{eingegriffen worden wäre}), \all{hätten \ldots\ vermeiden können} (structure modale active) / \all{hätten \ldots\ vermieden werden können} (structure modale passive).
\end{itemize}

\textbf{Barème indicatif (20 points) :} Respect du sujet et de la structure (Einleitung, Pro, Contra, Schluss) 5 ; Correction grammaticale (passif, plus-que-parfait, subjonctif II, ordre des mots) 6 ; Vocabulaire du chapitre 4 ; Cohérence et argumentation 3 ; Longueur respectée (env. 120-150 mots) 2.

\textbf{Points de vigilance pour l'élève :}
\begin{itemize}
\item Vérifier que les formes imposées (passif, plus-que-parfait, subjonctif II) sont bien présentes et correctes.
\item Ordre des mots : verbe conjugué en position 2 dans la principale ; dans les subordonnées (\all{wenn}, \all{dass}, \all{weil}, \all{nachdem}, \all{ob}), le verbe conjugué va à la fin.
\item Vocabulaire attendu : \all{die Einmischung, die Souveränität, der Einfluss, der Staat, die Macht, eingreifen, sich einmischen, unterdrücken, respektieren, die Unabhängigkeit, die Menschenrechte}.
\item Éviter les calques du français : « décider soi-même » se dit \all{selbst entscheiden} ou \all{selbst bestimmen}, jamais \all{*selber decidieren} ; « ingérence dans » = \all{Einmischung in + Akk.} (pas \all{*in + Dat.}).
\end{itemize}
\end{corrige}

\newpage

\coursec{Fiche bilan}

\begin{popfigure}
\begin{center}
\begin{tikzpicture}[mindmap, grow cyclic, every node/.style={concept}, concept color=popblue!60, text=white,
level 1/.append style={level distance=3.4cm, sibling angle=60},
level 2/.append style={level distance=2.2cm},
scale=0.92, transform shape]
\node[concept, font=\bfseries] {Les ingérences en politique}
  child[concept color=poporange!80] { node[concept] {Texte support}
    child { node[concept, font=\small] {Compréhension globale puis détaillée} }
    child { node[concept, font=\small] {Qui ? Quoi ? Pourquoi ?} }
  }
  child[concept color=popgreen!70!black] { node[concept] {Lexique}
    child { node[concept, font=\small] {die Einmischung, die Souveränität} }
    child { node[concept, font=\small] {der Einfluss, die Macht, der Staat} }
  }
  child[concept color=poppurple!80] { node[concept] {Le passif}
    child { node[concept, font=\small] {werden + Partizip II} }
    child { node[concept, font=\small] {von / durch + Dativ} }
  }
  child[concept color=poppink!80] { node[concept] {Plus-que-parfait}
    child { node[concept, font=\small] {hatte / war + Partizip II} }
    child { node[concept, font=\small] {antériorité : nachdem} }
  }
  child[concept color=popgold!80!black] { node[concept] {Konjunktiv II}
    child { node[concept, font=\small] {würde + Infinitiv} }
    child { node[concept, font=\small] {hätte, wäre, könnte} }
  }
  child[concept color=popblue!80] { node[concept] {Méthode}
    child { node[concept, font=\small] {Commentaire : 5 étapes} }
    child { node[concept, font=\small] {Résumé + production} }
  };
\end{tikzpicture}
\end{center}
\legende{Carte mentale de la leçon : les ingérences en politique}
\end{popfigure}

\begin{bilan}
\textbf{1. Lexique clé à connaître par cœur}
\begin{itemize}
  \item \all{die Einmischung, -en} : l'ingérence, l'immixtion ; \all{sich einmischen in + Akk.} : s'ingérer dans
  \item \all{die Souveränität, -en} : la souveraineté ; \all{souverän} : souverain
  \item \all{der Einfluss, -"e} : l'influence ; \all{Einfluss nehmen auf + Akk.} : exercer une influence sur
  \item \all{die Macht, -"e} : le pouvoir ; \all{mächtig} : puissant ; \all{die Großmacht, -"e} : la grande puissance
  \item \all{der Staat, -en} : l'État ; \all{staatlich} : étatique ; \all{die Staatsangehörigkeit} : la nationalité
  \item \all{das Land, -"er} : le pays ; \all{die Politik} : la politique ; \all{der Politiker, - / die Politikerin, -nen} : l'homme / la femme politique
  \item \all{die Wahl, -en} : l'élection ; \all{wählen} : élire, choisir ; \all{der Wähler, -} : l'électeur
  \item \all{die Unabhängigkeit} : l'indépendance ; \all{unabhängig} : indépendant ; \all{abhängig von} : dépendant de
  \item \all{die Intervention, -en} : l'intervention ; \all{die Sanktion, -en} : la sanction ; \all{der Druck} : la pression
\end{itemize}

\textbf{2. Grammaire à maîtriser}
\begin{center}
\begin{tabularx}{\linewidth}{|l|X|X|}
\hline
\rowcolor{popblueL} \textbf{Point} & \textbf{Règle} & \textbf{Exemple} \\
\hline
Passif Präsens & \all{werden} (conjugué) + Partizip II en fin & \all{Die Wahl wird beeinflusst.} (L'élection est influencée.) \\
\hline
Passif Präteritum & \all{wurde} + Partizip II & \all{Der Staat wurde erpresst.} (L'État a été soumis à un chantage.) \\
\hline
Passif Perfekt & \all{ist} + Partizip II + \all{worden} & \all{Es ist viel Geld gezahlt worden.} \\
\hline
Agent du passif & \all{von} + Datif (personne), \all{durch} + Akkusativ (moyen) & \all{von der Regierung} / \all{durch Propaganda} \\
\hline
Plus-que-parfait & \all{hatte} (haben) ou \all{war} (sein) + Partizip II & \all{Nachdem er gewählt worden war, ...} (Après qu'il eut été élu, ...) \\
\hline
Konjunktiv II & \all{würde} + Infinitif ; formes propres : \all{hätte, wäre, könnte, sollte} & \all{Man sollte sich nicht einmischen.} (On ne devrait pas s'ingérer.) \\
\hline
\end{tabularx}
\end{center}

\textbf{3. Méthodes express}
\begin{itemize}
  \item \textbf{Comprendre le texte} : lire le titre et la source, repérer le thème (ingérence), dégager la thèse de l'auteur, relever les arguments et les exemples.
  \item \textbf{Commenter} : 5 étapes — introduction (présenter le texte), résumé, analyse des idées, avis personnel argumenté, conclusion.
  \item \textbf{Résumé} : phrases courtes, au Präsens, sans citer le texte mot à mot, avec ses propres mots.
  \item \textbf{Traduire un passif français} : vérifier si l'agent est exprimé ; sinon, préférer \all{man} + actif : « On a critiqué la politique » = \all{Man hat die Politik kritisiert.}
\end{itemize}
\end{bilan}

\begin{aretenir}[Checklist avant le Bac]
\begin{itemize}
  \item $\square$ Je connais l'article et le pluriel des noms du thème : \all{die Einmischung, -en ; der Einfluss, -"e ; die Macht, -"e ; der Staat, -en ; die Souveränität, -en ; das Land, -"er}.
  \item $\square$ Je sais conjuguer \all{werden} au Präsens (\all{ich werde, du wirst, er wird}) et au Präteritum (\all{ich wurde}).
  \item $\square$ Je place le Partizip II à la fin de la phrase au passif.
  \item $\square$ Je distingue \all{von} + Datif (agent) et \all{durch} + Akkusativ (moyen).
  \item $\square$ Je forme le passif Perfekt : auxiliaire \all{sein} + Partizip II + \all{worden} (jamais \all{geworden}).
  \item $\square$ Je forme le plus-que-parfait : \all{hatte} ou \all{war} + Partizip II, et je choisis le bon auxiliaire (verbes de mouvement ou de changement d'état : \all{sein}).
  \item $\square$ J'utilise le plus-que-parfait après \all{nachdem} quand le principal est au passé.
  \item $\square$ Je connais les formes du Konjunktiv II : \all{würde, hätte, wäre, könnte, sollte, müsste}.
  \item $\square$ Je mets une majuscule à tous les noms allemands et les guillemets allemands „ … “ dans mes citations.
  \item $\square$ Je sais remplacer un passif sans agent par \all{man} + actif.
  \item $\square$ Je sais rédiger un résumé en 4 à 6 phrases avec mes propres mots.
  \item $\square$ Je sais structurer un commentaire : introduction, résumé, analyse, avis personnel, conclusion.
\end{itemize}
\end{aretenir}

\findecours

\end{document}
